\documentclass[12pt,oneside]{book}

% =====================================================
% METADATOS
% =====================================================
\newcommand{\universidad}{Universidad de Buenos Aires}
\newcommand{\facultad}{Facultad de Derecho}
\newcommand{\catedra}{Cátedra de Derecho Procesal}
\newcommand{\materia}{Juicio de Desalojo}
\newcommand{\submateria}{Derecho Procesal --- CPCCN}
\newcommand{\titulo}{Juicio de Desalojo: proceso, contratos y práctica profesional}
\newcommand{\titulocorto}{Juicio de Desalojo}
\newcommand{\autor}{Isaac Edgar Camacho Ocampo}
\newcommand{\anio}{2026}

% =====================================================
% PAQUETES
% =====================================================
\usepackage[utf8]{inputenc}
\usepackage[T1]{fontenc}
\usepackage[spanish,es-nodecimaldot]{babel}
\usepackage[top=2.5cm,bottom=2.5cm,left=3cm,right=2.5cm,headheight=15pt]{geometry}
\usepackage{amsmath}
\usepackage{amssymb}
\usepackage{mathtools}
\usepackage{graphicx}
\usepackage{float}
\usepackage{caption}
\usepackage{subcaption}
\usepackage{booktabs}
\usepackage{array}
\usepackage{longtable}
\usepackage{tabularx}
\usepackage{enumitem}
\usepackage{xcolor}
\usepackage[most]{tcolorbox}
\usepackage{fancyhdr}
\usepackage{ragged2e}
\usepackage[colorlinks=true,linkcolor=blue!50!black,citecolor=green!40!black,urlcolor=blue!60!black,bookmarks=true,pdftitle={Juicio de Desalojo: proceso, contratos y práctica profesional},pdfauthor={Isaac Edgar Camacho Ocampo},pdfsubject={Juicio de Desalojo --- Derecho Procesal}]{hyperref}

% =====================================================
% CONFIGURACIÓN GENERAL
% =====================================================
\setcounter{tocdepth}{2}
\setcounter{secnumdepth}{3}
\setlength{\parindent}{0pt}
\setlength{\parskip}{0.5em}
\setlist[itemize]{topsep=0.3em,itemsep=0.2em}
\setlist[enumerate]{topsep=0.3em,itemsep=0.2em}
\clubpenalty=10000
\widowpenalty=10000
\renewcommand{\arraystretch}{1.2}

% =====================================================
% ENCABEZADOS Y PIES
% =====================================================
\pagestyle{fancy}
\fancyhf{}
\fancyhead[L]{\small\nouppercase{\leftmark}}
\fancyhead[R]{\small\titulocorto}
\fancyfoot[C]{\thepage}
\renewcommand{\headrulewidth}{0.4pt}
\renewcommand{\footrulewidth}{0pt}
\fancypagestyle{plain}{\fancyhf{}\fancyfoot[C]{\thepage}\renewcommand{\headrulewidth}{0pt}}

% =====================================================
% COMANDOS MATEMÁTICOS (disponibles; no se usan en este material)
% =====================================================
\newcommand{\abs}[1]{\left|#1\right|}
\newcommand{\norm}[1]{\left\lVert#1\right\rVert}
\newcommand{\vect}[1]{\mathbf{#1}}
\newcommand{\deriv}[2]{\frac{d#1}{d#2}}
\newcommand{\pderiv}[2]{\frac{\partial#1}{\partial#2}}

% =====================================================
% CAJAS SEMÁNTICAS
% =====================================================
\newtcolorbox{definition}[1]{enhanced,breakable,title={Definición: #1},fonttitle=\bfseries,colback=blue!3,colframe=blue!50!black,boxrule=0.8pt,arc=2mm}
\newtcolorbox{important}{enhanced,breakable,title={Importante},fonttitle=\bfseries,colback=yellow!8,colframe=orange!70!black,boxrule=0.8pt,arc=2mm}
\newtcolorbox{example}[1]{enhanced,breakable,title={Ejemplo: #1},fonttitle=\bfseries,colback=green!4,colframe=green!40!black,boxrule=0.8pt,arc=2mm}
\newtcolorbox{formula}[1]{enhanced,breakable,title={Regla: #1},fonttitle=\bfseries,colback=purple!4,colframe=purple!50!black,boxrule=0.8pt,arc=2mm}
\newtcolorbox{exercise}[1]{enhanced,breakable,title={Ejercicio: #1},fonttitle=\bfseries,colback=gray!5,colframe=gray!60!black,boxrule=0.8pt,arc=2mm}
\newtcolorbox{note}{enhanced,breakable,title={Nota},fonttitle=\bfseries,colback=white,colframe=gray!50,boxrule=0.6pt,arc=2mm}
\newtcolorbox{normacitada}{enhanced,breakable,title={Norma citada},fonttitle=\bfseries,colback=red!3,colframe=red!45!black,boxrule=0.8pt,arc=2mm}
\newtcolorbox{sintesis}{enhanced,breakable,title={Síntesis},fonttitle=\bfseries,colback=gray!4,colframe=black!60,boxrule=0.8pt,arc=2mm}

% =====================================================
% CUADRO CORNELL (tabla tabularx real)
% =====================================================
\newcommand{\cqa}[2]{\textbf{#1} & #2 \\ \midrule}
\newcommand{\cresumen}[1]{\multicolumn{2}{@{}X@{}}{\textbf{Resumen.} #1} \\ \bottomrule}
\newenvironment{cornell}[1]
 {\par\medskip\noindent{\large\bfseries #1}\par\smallskip
  \noindent\tabularx{\textwidth}{@{}>{\RaggedRight\arraybackslash}p{0.30\textwidth}|>{\RaggedRight\arraybackslash}X@{}}
  \toprule
  \textbf{Preguntas / palabras clave} & \textbf{Notas / respuestas} \\ \midrule}
 {\endtabularx\par\medskip}

% =====================================================
% DOCUMENTO
% =====================================================
\begin{document}

\begin{titlepage}
\thispagestyle{empty}
\begin{center}
\vspace*{1cm}
{\large \textsc{\universidad}}\\[0.2cm]
{\large \textsc{\facultad}}\\[0.2cm]
{\large \catedra}\\[3cm]
{\Huge \textbf{\materia}}\\[0.4cm]
{\normalsize \submateria}\\[1cm]
{\LARGE \titulo}\\[1.2cm]
\rule{0.70\textwidth}{0.5pt}\\[0.5cm]
{\large Apuntes de estudio integrados}\\[0.5cm]
\rule{0.70\textwidth}{0.5pt}
\vfill
{\large \autor}\\[0.3cm]
{\large \anio}
\end{center}
\vfill
\begin{flushleft}
\textit{Este es un modesto aporte a los estudiantes de <<\materia>>; no pretende sustituir el material oficial de la materia.}
\end{flushleft}
\end{titlepage}

\frontmatter
\tableofcontents

\mainmatter
\part{Marco general del desalojo}

\chapter{Presentación de la materia y principios rectores}\label{cap:principios}

El juicio de desalojo es el proceso mediante el cual se persigue la restitución de un inmueble que ocupa una persona que no tiene derecho a retenerlo, o cuyo derecho a retenerlo se ha extinguido. Este capítulo ubica la materia dentro del ordenamiento: explica por qué su estudio tiene aplicación inmediata, cuáles son los cuerpos normativos que se combinan en ella, qué principios constitucionales y convencionales la orientan y por qué la perspectiva de vulnerabilidad atraviesa todo el desarrollo posterior. Los capítulos siguientes parten de estas bases para estudiar el marco normativo de las locaciones (capítulo~\ref{cap:normativo}) y los contratos que sirven de título al desalojo (capítulo~\ref{cap:titulos}).

\section{Relevancia práctica del desalojo}

El desalojo constituye una de las materias de mayor aplicación práctica de la carrera: tarde o temprano, como abogado o abogada, un familiar, un amigo o un cliente atravesará un conflicto de esta naturaleza. Comprender la diferencia entre tenencia y posesión, saber elegir la vía procesal correcta y conocer qué firma es válida en cada instrumento son herramientas que se emplean desde el inicio del ejercicio profesional, más allá de la orientación que cada estudiante elija dentro del Derecho.

Se trata, además, de un contenido útil para la vida cotidiana de cualquier persona: como inquilino, propietario, comodatario o simplemente como quien acompaña a un familiar en un conflicto habitacional. Entender qué es un contrato de locación, qué diferencia existe entre poseer y tener, y qué firma resulta válida permite tomar decisiones más informadas y prudentes, incluso antes de ejercer la profesión.

\section{Cuerpos normativos que se combinan}

\begin{important}
El estudio del desalojo se apoya en tres cuerpos normativos que se entrecruzan permanentemente: el \textbf{Código Procesal Civil y Comercial de la Nación} (con sus diferencias respecto del Código Procesal de la provincia de Buenos Aires), el \textbf{Código Civil y Comercial de la Nación} ---en particular, la parte especial de contratos de locación--- y la \textbf{normativa constitucional y convencional}, incluidos los tratados internacionales de derechos humanos.
\end{important}

La referencia principal es el Código Procesal Civil y Comercial de la Nación (CPCCN). Sin embargo, el ejercicio profesional también se desarrolla en la provincia de Buenos Aires, por lo cual deben conocerse los comentarios y diferencias respecto del Código Procesal bonaerense; existen tantos códigos procesales como provincias. A su vez, el Código Civil y Comercial de la Nación (CCCN) contiene normas de contenido procesal relevantes para el desalojo de inmuebles urbanos, especialmente en la parte especial de contratos de locación. Esta interrelación llevó a citar la opinión de un reconocido procesalista, según la cual <<la última gran reforma del Código Procesal fue, en rigor, el Código Civil>>.

El estudio tampoco se limita al proceso judicial en sentido estricto: incorpora las normas constitucionales y convencionales aplicables y los criterios sentados por organismos internacionales, entre ellos la Corte Interamericana de Derechos Humanos.

\begin{note}
Resulta recomendable contar, en cuanto sea posible, con un código comentado ---aunque sea de varios tomos---, porque incorpora jurisprudencia, explicaciones y remisiones a otros artículos que un código sin comentar no ofrece, y debe actualizarse periódicamente. Hoy el acceso a la bibliografía jurídica es amplio gracias a los libros digitales y a las bibliotecas universitarias, incluso después de la graduación.
\end{note}

\section{Principios y sistemas}\label{sec:principios-sistemas}

Los principios y los sistemas constituyen categorías jurídicas diferentes.

\begin{definition}{Principios y sistemas}
Los \textbf{principios} son creados por el constituyente: son los presupuestos jurídicos fundantes de un ordenamiento y se encuentran en la Constitución Nacional y en los tratados internacionales, que a partir de la reforma constitucional de 1994 poseen igual jerarquía que ella. Los \textbf{sistemas} son creados por el legislador: son el método por el cual se traduce a la práctica, en normas positivas, un principio constitucional o convencional.
\end{definition}

El legislador debe considerar el ordenamiento jurídico, el presupuesto jurídico fundante y los principios vigentes en la República Argentina, y elaborar la ley en función de ellos. Cuando el sistema legal no respeta los principios constitucionales, los abogados y las abogadas terminan planteando cuestiones de inconstitucionalidad e inconvencionalidad ante los tribunales.

\begin{example}{El proyecto de ley de inviolabilidad de la propiedad privada}
Un sistema creado por el legislador ---como un eventual proyecto de ley sobre inviolabilidad de la propiedad privada--- debe someterse al Congreso para su aprobación y responder a los principios constitucionales.
\end{example}

En materia de desalojo se recurre con frecuencia a los tratados internacionales bajo una perspectiva de vulnerabilidad.

\begin{example}{Las miradas delante y detrás de la puerta}
Al analizar la complejidad del desalojo puede distinguirse la mirada de quien está delante de la puerta de una vivienda de la mirada de quien está detrás de ella. En un desalojo no solo está en juego el derecho de propiedad ---ubicado en un lugar jerárquicamente elevado entre los principios---, sino también otros derechos, como la igualdad y el respeto a la persona humana, que están en juego del otro lado de la puerta.
\end{example}

\section{La perspectiva de vulnerabilidad como eje transversal}\label{sec:vulnerabilidad}

La materia incorpora la perspectiva de vulnerabilidad como eje transversal, en tanto el operador jurídico puede encontrarse, según el caso, en cualquiera de los dos lados del conflicto: procurando el lanzamiento para recuperar la tenencia, o bien acompañando a un sujeto vulnerable amenazado de desalojo.

\subsection{Las Reglas de Brasilia}

\begin{formula}{Reglas de Brasilia sobre Acceso a la Justicia de las Personas en Condición de Vulnerabilidad}
Fueron incorporadas a la normativa argentina mediante la Acordada 5/2009 de la Corte Suprema de Justicia de la Nación. No provienen de una ley ni de un tratado internacional, sino de esa acordada, y son aplicables a todo tipo de proceso judicial, no solo al desalojo.
\end{formula}

Estas reglas se aplican al desalojo porque la condición de persona vulnerable no se limita al adulto mayor, al niño, niña o adolescente, ni a la persona con capacidad restringida. Cualquier persona puede resultar vulnerable, según quién esté del otro lado del conflicto: la posibilidad de educación, la historia personal y la situación económica de cada quien determinan su grado de vulnerabilidad frente al otro.

\begin{example}{El jubilado que vive de la renta de un inmueble}
Una persona adulta, con un trabajo modesto y jubilada, que tiene un pequeño departamento dado en locación, puede encontrarse en condición de vulnerabilidad si el inquilino deja de pagar o no desocupa el inmueble. Con su jubilación mínima y su trabajo mínimo, necesita el ingreso adicional del canon locativo para vivir con cierta comodidad. Si queda atrapada en un proceso de desalojo sin cobrar el alquiler, también puede ser vulnerable, aunque no encaje en ninguna de las categorías clásicas.
\end{example}

\begin{important}
Se prefiere hablar de <<perspectiva de vulnerabilidad>> antes que limitarse a la <<perspectiva de género>>, sin que esto implique dejarla de lado. La intención es abrir el abanico de análisis a cualquier sujeto vulnerable, sin focalizarlo exclusivamente en una única categoría.
\end{important}

\subsection{La vulnerabilidad depende del contexto}

La regulación de los alquileres ofrece un ejemplo de cómo la vulnerabilidad no depende de una categoría fija (véase la sección~\ref{sec:efectos-economicos}): tanto quien pretendía locar un inmueble para vivir como quien disponía de un inmueble para dar en locación resultaron, en distintos momentos y por distintas razones, sujetos vulnerables.

\section{El Código Civil y Comercial como código pro persona}

El Código Civil y Comercial, vigente desde agosto de 2015, es un código pro persona: abandona el positivismo puro para constituirse en un constitucionalismo privado, por el cual la normativa específica queda subordinada a la impronta de la norma constitucional e infraconstitucional. En esa lógica opera la autonomía de la voluntad, desarrollada en la sección~\ref{sec:autonomia}.

\section{La importancia de las formas procesales}

\begin{example}{Rechazo por incumplimiento formal ---art.\ 280 CPCCN}
En un fallo de la Corte Suprema de Justicia de la Nación, una comunidad indígena argentina reclamó la restitución de tierras invocando su carácter ancestral. La Corte no llegó a analizar el fondo del reclamo: rechazó el recurso por incumplimiento de los requisitos formales del artículo 280 del CPCCN. El caso ilustra que el desconocimiento de las normas procesales puede determinar la pérdida de un proceso, incluso cuando existan fundamentos sustanciales atendibles, sin que ello implique tomar posición sobre si la comunidad tenía o no razón en cuanto al fondo.
\end{example}

\begin{normacitada}
\textbf{Artículo 280 del CPCCN} ---regula los requisitos formales de admisibilidad del recurso extraordinario ante la Corte Suprema de Justicia de la Nación. Su incumplimiento formal puede determinar el rechazo del recurso sin analizar el fondo del planteo.
\end{normacitada}

\begin{sintesis}
El desalojo se estudia desde tres cuerpos normativos (CPCCN, CCCN y normativa constitucional y convencional); los principios provienen del constituyente y los sistemas del legislador; la perspectiva de vulnerabilidad, apoyada en las Reglas de Brasilia, es un eje transversal; y el respeto de las formas procesales puede decidir un caso antes de llegar al fondo.
\end{sintesis}

\begin{cornell}{Cuadro Cornell del capítulo~\ref{cap:principios}}
\cqa{¿Qué tres cuerpos normativos se combinan en el desalojo?}{El CPCCN (con sus diferencias respecto del Código Procesal bonaerense), el CCCN ---especialmente los contratos de locación--- y la normativa constitucional y convencional, incluidos los tratados de derechos humanos.}
\cqa{¿En qué se diferencian principios y sistemas?}{Los principios los crea el constituyente y son presupuestos fundantes del ordenamiento; los sistemas los crea el legislador y trasladan el principio a la realidad. Si el sistema no respeta el principio, se plantean cuestiones de inconstitucionalidad e inconvencionalidad.}
\cqa{¿Qué son las Reglas de Brasilia y de dónde provienen?}{Reglas de acceso a la justicia de personas en condición de vulnerabilidad, incorporadas por la Acordada 5/2009 de la Corte Suprema. No provienen de una ley ni de un tratado y se aplican a todo proceso judicial.}
\cqa{¿Por qué se habla de <<perspectiva de vulnerabilidad>> y no solo de género?}{Para abrir el análisis a cualquier sujeto vulnerable; cualquier persona puede serlo según quién esté del otro lado del conflicto. No implica dejar de lado la perspectiva de género.}
\cqa{¿Qué enseña el rechazo por el art.\ 280 CPCCN?}{Que el incumplimiento de requisitos formales puede hacer perder un proceso sin que se analice el fondo, aunque existan fundamentos sustanciales atendibles.}
\cresumen{El desalojo exige combinar normas procesales, civiles y constitucionales; leerlas a la luz de los principios y de la perspectiva de vulnerabilidad; y cuidar las formas, porque su incumplimiento puede decidir el resultado.}
\end{cornell}

\chapter{Marco normativo de las locaciones}\label{cap:normativo}

Antes de analizar un conflicto de desalojo fundado en una locación, es necesario saber qué normas rigen el contrato. Este capítulo recorre la sucesión legislativa reciente, el modo de resolver cuál es la norma aplicable a un contrato según su fecha, el papel de la autonomía de la voluntad y del orden público, los efectos económicos de la regulación anterior, la discusión sobre el derecho del consumidor y los requisitos de forma y de firma del instrumento contractual.

\section{Evolución legislativa de los alquileres}\label{sec:evolucion}

Lo primero que hay que ver cuando se tiene un juicio de desalojo fundado en un contrato de locación es la \textbf{fecha del contrato}, porque existen distintas normas con vigencia temporal diferenciada. Las normas que aparecen derogadas pueden, sin embargo, tener vigencia según la fecha del contrato que se esté analizando.

\begin{normacitada}
\textbf{Código Civil y Comercial de la Nación (Ley 26.994)} ---vigente desde agosto de 2015. Modificó aspectos de la anterior Ley 23.091 de Locaciones Urbanas e incorporó la regulación del contrato de locación (arts.\ 1187 y siguientes) al cuerpo unificado.
\end{normacitada}

\begin{normacitada}
\textbf{Ley 27.551 (Nueva Ley de Alquileres)} ---comenzó a regir el 1.º de julio de 2020 y modificó diversos aspectos del CCCN en materia de locación de inmuebles. Sus disposiciones rigen para los contratos celebrados a partir de su entrada en vigencia, sin operar retroactivamente sobre contratos anteriores. Fue posteriormente modificada por la Ley 27.737.
\end{normacitada}

\begin{normacitada}
\textbf{Decreto de Necesidad y Urgencia N.º 70/2023}, del 29 de diciembre de 2023 ---modificó de manera sustancial el régimen de locaciones urbanas, derogando buena parte del esquema establecido por las Leyes 27.551 y 27.737 y ampliando la autonomía de la voluntad de las partes. Es la norma operativa en la práctica profesional actual.
\end{normacitada}

\begin{table}[H]
\centering
\caption{Norma aplicable según la fecha de celebración del contrato de locación}
\begin{tabularx}{\textwidth}{>{\raggedright\arraybackslash}p{0.32\textwidth} X}
\toprule
\textbf{Período de celebración} & \textbf{Norma aplicable} \\
\midrule
Antes de agosto de 2015 & Ley 23.091 (Ley de Locaciones Urbanas) y régimen del Código Civil derogado. \\
Entre agosto de 2015 y junio de 2020 & CCCN (Ley 26.994), arts.\ 1187 y ss. \\
Entre julio de 2020 y diciembre de 2023 & CCCN reformado por la Ley 27.551 (Nueva Ley de Alquileres), luego modificada por la Ley 27.737. \\
A partir del 29 de diciembre de 2023 & Régimen resultante del DNU 70/2023, con mayor amplitud de la autonomía de la voluntad. \\
\bottomrule
\end{tabularx}
\end{table}

\section{Derecho transitorio: el artículo 7 del CCCN}

Un caso frecuente es el de un conflicto derivado de un contrato celebrado bajo las Leyes 27.551 y 27.737, hoy derogadas en gran parte por el DNU 70/2023. Corresponde entonces acudir al artículo 7 del CCCN para determinar qué norma se aplica a las relaciones jurídicas nacidas bajo el régimen anterior.

\begin{formula}{Artículo 7 del CCCN ---eficacia temporal}
A partir de su entrada en vigencia, las leyes se aplican a las consecuencias de las relaciones y situaciones jurídicas existentes, y las nuevas leyes no tienen efecto retroactivo, sean o no de orden público, salvo disposición en contrario.
\end{formula}

Aquellas leyes derogaron buena parte de las normas del Código Civil en cuanto al procedimiento, aunque no en su totalidad, y establecían además pautas de procedimiento en el Código Procesal local. El Código Procesal nacional se ve afectado por un proyecto de ley que, hasta el momento de los apuntes originales, no se encontraba vigente.

\begin{important}
Frente a cualquier conflicto de locación, el primer paso metodológico consiste en determinar la fecha de celebración del contrato para identificar la norma temporalmente aplicable; el segundo, en leer íntegramente el instrumento para verificar qué pactaron las partes, controlando que ninguna cláusula vulnere el orden público vigente al momento de la celebración.
\end{important}

\section{Autonomía de la voluntad, orden público y contrato como ley entre las partes}\label{sec:autonomia}

\begin{formula}{Artículo 958 del CCCN ---libertad de contratación}
Consagra la libertad de contratar: las partes son libres para celebrar un contrato y determinar su contenido, dentro de los límites impuestos por la ley, el orden público, la moral y las buenas costumbres.
\end{formula}

En materia de locación y comodato opera plenamente la autonomía de la voluntad: lo que las partes resuelvan, contraten y plasmen en el contrato vale como la ley misma, salvo que se vulneren normas de orden público. Por eso la aplicación del Código es supletoria de lo que las partes acordaron, y ante cada conflicto resulta indispensable leer con detenimiento el contrato: aunque suele percibirse como un documento simple, contiene cláusulas relevantes ---muchas veces en letra pequeña--- que la contraparte no leyó o no advirtió.

\begin{definition}{Cuestiones disponibles y orden público}
Existen cuestiones disponibles, sobre las cuales las partes pueden pactar libremente, y cuestiones donde interviene el orden público, que limitan lo que las partes pueden acordar por su sola voluntad. Cuando hay orden público, existe un límite a la autonomía de la voluntad; cuando no lo hay, las partes pueden disponer libremente.
\end{definition}

Si una cláusula implica la renuncia a un derecho, es preciso determinar si es válida: existen artículos del CCCN que son de orden público y, por lo tanto, irrenunciables, más allá de lo pactado. Recae en quien recibe y analiza el contrato la responsabilidad de detectar las cláusulas potencialmente nulas.

\begin{example}{La carta documento con norma derogada}
Al recibir un caso, lo primero que se solicitó a la persona que consultaba fue el contrato de locación, sin confiar en el relato oral del cliente sobre su fecha o contenido, y exigiendo siempre el documento original o una copia fehaciente. El contrato había sido celebrado en 2025, ya bajo la plena vigencia del DNU 70/2023, y la carta documento invocaba un artículo del CCCN derogado precisamente por ese decreto. Además, la cláusula contractual invocada contenía una renuncia de las partes, mayores y capaces, a valerse de una norma vinculada a los desperfectos que afectan el uso y goce del inmueble. La enseñanza es que hay que leer qué firmaron las partes: el contrato es ley para ellas.
\end{example}

% TODO: verificar referencia presente en las fuentes originales. La fuente (clase del 27/08) menciona oralmente el <<artículo 1223 del Código Civil>> y agrega una precisión editorial: según esa nota, el art. 1223 CCCN regula el <<Desalojo>> y los desperfectos que afectan el uso y goce se regulan en los arts. 1200 a 1204 CCCN (en particular el art. 1201). Se conserva la advertencia de la fuente sin resolverla.
\begin{note}
La fuente original advierte que el número de artículo mencionado oralmente en este caso (art.\ 1223) debe verificarse antes de invocarlo en un escrito judicial. La nota editorial de esa fuente ubica la cuestión de los desperfectos que afectan el uso y goce de la cosa en las normas sobre conservación y garantía por defectos (arts.\ 1200 a 1204 CCCN, en particular el art.\ 1201). El principio que el caso ilustra no depende de ese número: el contrato es ley para las partes.
\end{note}

\section{Efectos económicos de la regulación anterior}\label{sec:efectos-economicos}

Bajo la ley anterior, el problema principal fue que quienes tenían viviendas o locales para dar en locación retiraron buena parte de su oferta del mercado: aproximadamente el ochenta por ciento de la oferta se encontraba restringida. La actualización del canon locativo era anual, las propiedades muchas veces se devolvían deterioradas, y el costo de refacción resultaba equivalente al valor de la locación del año anterior.

\begin{example}{El monoambiente inaccesible}
Una persona proveniente de otra provincia que desea estudiar en la Facultad de Derecho de la Universidad de Buenos Aires y busca alquilar un monoambiente encuentra que, por la escasez de oferta generada por la regulación anterior, ese monoambiente llegaba a valer como una casa con pileta. Quienes disponían de inmuebles para locar preferían no alquilar antes que someterse a una ley con una impronta de mayor amparo a quien usaba y gozaba de la cosa; muchos preferían pagar las expensas de su bolsillo. La consecuencia paradójica fue que la norma pensada para proteger al inquilino terminó perjudicando también a quienes buscaban alquilar.
\end{example}

Cuando el DNU cambió el esquema y los contratos vencidos se actualizaron, muchos inquilinos se retiraron de esas locaciones porque resultaban más económicas las nuevas ofertas que la actualización de los contratos regidos por la ley anterior.

\begin{important}
La normativa que buscaba proteger a un sector terminó perjudicando a todos. Quien pretendía locar un inmueble para vivir y quien disponía de un inmueble para dar en locación resultaron, en distintos momentos y por distintas razones, sujetos vulnerables. La vulnerabilidad no depende de una categoría fija, sino del contexto y de la relación concreta entre las partes.
\end{important}

\section{Locación y derecho del consumidor}

Se plantea en qué casos podría aplicarse la ley de defensa del consumidor a la locación de inmuebles. Resulta relevante la postura del Dr.\ Ricardo Lorenzetti, ministro de la Corte Suprema de Justicia de la Nación, según la cual en la locación existe un consumidor y un usuario, dado que hay uso y goce de la cosa y contraprestación de un precio.

\begin{formula}{Artículo 1093 del CCCN ---contrato de consumo}
Es el celebrado entre un consumidor o usuario final, con una persona humana o jurídica que actúe profesional u ocasionalmente, para la adquisición, uso o goce de bienes o servicios, para uso privado, familiar o social.
\end{formula}

\begin{formula}{Artículo 1187 del CCCN ---contrato de locación}
Es aquel por el cual una parte se obliga a otorgar a otra el uso y goce temporario de una cosa, a cambio del pago de un precio en dinero.
\end{formula}

Existen distintas posturas doctrinarias: una tesis restringida, según la cual las normas de consumo nunca se aplican; una tesis amplia, según la cual siempre se aplican; y una postura intermedia, que mira la conflictiva concreta caso por caso.

\begin{definition}{La conflictiva por sobre el molde preconcebido}
No todo conflicto derivado de una locación debe encuadrarse automáticamente en el derecho del consumidor, ni descartarse siempre esa posibilidad. Lo importante ---no solo para el desalojo, sino para cualquier conflicto que llegue a un juez, jueza, abogado o abogada--- es escuchar primero la conflictiva concreta y no encajarla rápidamente en un molde preconcebido.
\end{definition}

\section{Forma del contrato y tipos de firma}\label{sec:firma}

\subsection{Forma escrita del contrato de locación de inmuebles}

\begin{formula}{Artículo 1188 del CCCN ---forma}
El contrato de locación de inmueble exige, para su consentimiento, la forma escrita. Cuando la locación tiene por objeto un inmueble, no se admite la contratación verbal: debe volcarse a un soporte documental.
\end{formula}

Para iniciar los procesos de desalojo fundados en vencimiento de contrato o en falta de pago se necesita el instrumento, porque se trata de un contrato consensual cuyo consentimiento, para inmuebles, exige por disposición del legislador la forma escrita. Este contenido se ubica dentro de los derechos personales: primero los contratos en general (autonomía de la voluntad, art.\ 958) y luego el contrato de locación en particular. En el Código Procesal, este proceso se ubica dentro de los procesos especiales.

\subsection{Firma hológrafa}

\begin{definition}{Firma hológrafa}
Es la firma inserta de puño y letra: la firma tradicional, que se utilizaba antes y se sigue utilizando.
\end{definition}

Frente al desconocimiento de la firma hológrafa por quien la suscribió, la única vía para acreditar su autenticidad es la \textbf{prueba pericial caligráfica}. Por ello resulta decisiva la etapa de instrucción previa a la demanda: al redactarla no solo deben describirse los hechos conforme a la normativa procesal aplicable, sino también anticipar si la contraparte negará o desconocerá ese hecho y qué medio de prueba se ofrecerá.

\begin{normacitada}
\textbf{Artículo 330 del CPCCN} ---establece los requisitos de la demanda, entre ellos los hechos en que se funda, explicados claramente, y el ofrecimiento de toda la prueba de que la parte intente valerse.
\end{normacitada}

\subsection{Firma electrónica y firma digital}

La firma digital está operativa en la República Argentina y se diferencia de la firma electrónica. Es frecuente que notarios y notarias utilicen este tipo de firma en la celebración de contratos de locación o comodato.

\begin{definition}{Firma electrónica y firma digital}
La \textbf{firma electrónica} está avalada por una institución privada. La \textbf{firma digital} está avalada por un organismo público o, en el caso del notario ---considerado funcionario público---, da fe pública de esa firma en el instrumento. No son equivalentes: cada una habilita un tipo de prueba distinto ante su desconocimiento.
\end{definition}

\begin{table}[H]
\centering
\caption{Tipos de firma en los instrumentos contractuales}
\begin{tabularx}{\textwidth}{>{\raggedright\arraybackslash}p{0.22\textwidth} X X}
\toprule
\textbf{Tipo de firma} & \textbf{Quién la avala} & \textbf{Régimen probatorio ante desconocimiento} \\
\midrule
Hológrafa (de puño y letra) & No requiere organismo; es la firma manuscrita tradicional. & Corresponde prueba pericial caligráfica; la carga de acreditar la autenticidad recae en quien invoca la firma. \\
Electrónica & Institución privada (por ejemplo, el Colegio de Abogados). & La carga de probar su autenticidad recae en quien la invoca; no se presume la autoría. \\
Digital & Organismo público, o notario en ejercicio de la fe pública. & Se presume la autoría y la integridad del documento; se invierte la carga de la prueba, que pasa a quien la desconoce. \\
Certificada ante notario & El escribano da fe únicamente de que la firma corresponde a esa persona, no del contenido íntegro del instrumento. & Queda constancia en el libro de requerimientos de la escribanía; el escribano certifica la firma, no el contenido. \\
\bottomrule
\end{tabularx}
\end{table}

\begin{example}{Firma digital y firma hológrafa frente al desconocimiento}
La firma hológrafa puede ser atacada y, hasta que un perito calígrafo determine que no corresponde a la persona, queda en tela de juicio. La firma digital también puede atacarse, pero con la carga de la prueba invertida: se presume la autoría de quien firmó y que el documento no fue adulterado, y queda a cargo de la otra parte probar lo contrario. La firma electrónica, si es desconocida, debe ser probada por quien la invoca.
\end{example}

\begin{note}
El tipo de firma utilizado en la práctica varía según el fuero: en el fuero nacional, los jueces y juezas no cuentan con firma digital, sino con firma electrónica; en la provincia de Buenos Aires, los jueces sí cuentan con firma digital. Al presentar un escrito ante el sistema del Poder Judicial de la Nación mediante firma electrónica, no es necesario volver a firmarlo en forma hológrafa.
\end{note}

\subsection{Firma del cliente y poder en la práctica}

Los tribunales suelen rechazar escritos firmados electrónicamente por el propio cliente, exigiendo en su lugar la firma hológrafa digitalizada.

\begin{example}{Firma hológrafa escaneada, no fotografiada}
Cuando el cliente no puede concurrir en persona, el escrito se envía a la persona designada para gestionar la firma; esa persona hace firmar de puño y letra al cliente, escanea el documento en formato PDF ---nunca como fotografía--- y lo reenvía al abogado responsable, quien recién entonces lo sube al expediente electrónico.
\end{example}

\begin{important}
\textbf{Artículo 47 del CPCCN y reforma del CCCN (2015).} El Código Procesal exigía, en su artículo 47, escritura de poder general, judicial o especial. El CCCN, vigente desde agosto de 2015, flexibiliza ese requisito y permite que el poder para el abogado se otorgue por instrumento privado o por instrumento privado con firma certificada. En cualquier caso, el poder debe acompañarse en la primera presentación junto con la demanda, para que el juez verifique la personería; de lo contrario, el demandado puede oponer la excepción de falta de personería.
\end{important}

En cuestiones de familia y en demandas de inicio en materia patrimonial, resulta preferible que el propio cliente firme la demanda ---presentándose primero con patrocinio letrado--- y que recién después se incorpore el poder, de modo que quede constancia de su voluntad desde el inicio. Los hechos que el cliente relata en la entrevista suelen ser incompletos o se omiten por vergüenza o por desconocimiento de su relevancia jurídica; en cuestiones de familia o de orden público conviene requerir un relato escrito y firmado, o un correo electrónico enviado por el propio cliente.

\begin{example}{El divorcio de común acuerdo que derivó en una demanda contra la propia abogada}
En un caso anterior a la reforma de 2015 se patrocinó un divorcio que se tramitaba en apariencia de común acuerdo. Tiempo después, una de las partes planteó la nulidad de lo actuado y demandó también a la profesional interviniente, alegando haber sido obligada a firmar el convenio. El reclamo pudo desestimarse porque la clienta había firmado como patrocinada, habiéndolo leído, y esa firma había sido puesta ante un escribano público que certificó su autenticidad.
\end{example}

\begin{important}
La firma hológrafa certificada protege tanto al cliente como al abogado: exigir una firma manuscrita, certificada o acompañada de un relato escrito del cliente no es solo un requisito formal para que el tribunal acepte el escrito, sino también una herramienta de resguardo profesional frente a reclamos posteriores.
\end{important}

\begin{example}{El escrito que nadie había despachado}
Se sostenía de memoria que determinado escrito ya había sido resuelto por el juzgado. Al revisar el expediente se comprobó que la resolución que se creía reciente correspondía a un escrito de años atrás, y que el planteo verdaderamente pendiente nunca había sido despachado, lo que se advirtió recién cuando el caso llegó a la Cámara por otro planteo. La conclusión es que la memoria puede fallar: conviene dejar todo por escrito y llevar una carpeta paralela de seguimiento del expediente, sin confiar exclusivamente en el sistema informático del Poder Judicial.
\end{example}

\begin{important}
El instrumento contractual base del desalojo puede estar suscripto mediante firma hológrafa, electrónica, digital o certificada ante notario. Cada forma exige, ante su desconocimiento, una estrategia probatoria distinta: pericial caligráfica para la hológrafa; carga de la prueba en cabeza de quien invoca la firma electrónica; y carga de la prueba en cabeza de quien desconoce la firma digital, que goza de presunción de autoría e integridad.
\end{important}

\begin{cornell}{Cuadro Cornell del capítulo~\ref{cap:normativo}}
\cqa{¿Qué es lo primero que hay que verificar en un desalojo fundado en locación?}{La fecha del contrato, porque distintas normas tienen vigencia temporal diferenciada y las derogadas pueden seguir aplicándose según esa fecha.}
\cqa{¿Qué norma rige según el período de celebración?}{Ley 23.091 (antes de agosto de 2015); CCCN, Ley 26.994 (agosto de 2015 a junio de 2020); CCCN reformado por la Ley 27.551, luego modificada por la Ley 27.737 (julio de 2020 a diciembre de 2023); DNU 70/2023 (desde el 29 de diciembre de 2023).}
\cqa{¿Qué establece el art.\ 7 CCCN y para qué sirve acá?}{Que las leyes se aplican a las consecuencias de las relaciones existentes y no tienen efecto retroactivo, salvo disposición en contrario. Sirve para decidir qué norma rige un contrato nacido bajo el régimen anterior.}
\cqa{¿Qué papel cumplen el art.\ 958 y el orden público?}{El art.\ 958 consagra la libertad de contratar; lo pactado vale como ley para las partes salvo que se vulneren normas de orden público, que limitan la autonomía de la voluntad y hacen irrenunciables ciertos derechos.}
\cqa{¿Por qué la regulación anterior perjudicó también a los inquilinos?}{Porque cerca del ochenta por ciento de la oferta se retiró del mercado, lo que encareció el acceso a la vivienda; esto muestra que la vulnerabilidad depende del contexto y no de una categoría fija.}
\cqa{¿Se aplica siempre el derecho del consumidor a la locación?}{Hay tres posturas: restringida (nunca), amplia (siempre) e intermedia (caso por caso). Lo central es escuchar primero la conflictiva concreta y no encajarla en un molde.}
\cqa{¿Qué diferencia hay entre firma hológrafa, electrónica y digital ante su desconocimiento?}{Hológrafa: pericial caligráfica. Electrónica: la carga de probarla recae en quien la invoca. Digital: se presume autoría e integridad y la carga recae en quien la desconoce.}
\cresumen{Todo análisis de locación comienza por la fecha del contrato y por su lectura íntegra, porque rige la autonomía de la voluntad dentro del orden público. La forma escrita es exigida para inmuebles y el tipo de firma determina la estrategia probatoria.}
\end{cornell}

\chapter{Locación y comodato como títulos del desalojo}\label{cap:titulos}

El desalojo no se limita al vencimiento de un contrato de locación o a la falta de pago del canon: otros títulos habilitan la vía, como el comodato o la sublocación. Este capítulo parte de la distinción entre tenencia y posesión, que explica qué se discute en el juicio de desalojo, y estudia luego los caracteres de los dos contratos que habitualmente sirven de título: la locación y el comodato. Cierra con las formas de extinción anticipada y con la conveniencia de agotar las vías alternativas antes de litigar.

\section{Tenencia y posesión}\label{sec:tenencia-posesion}

El desalojo suele asociarse al vencimiento de un contrato de locación o a la falta de pago, pero existen otros motivos que habilitan la vía: el comodato, el subarrendamiento, entre otros. La ejecución de una hipoteca, en cambio, no constituye uno de esos supuestos, porque persigue un objetivo distinto: la subasta del bien.

En el juicio de desalojo se persigue la restitución del bien que ocupa una persona que no tiene derecho a retenerlo, o cuyo derecho a retenerlo se ha extinguido y debe restituirlo. En ese proceso se discute la \textbf{tenencia}, a diferencia de la \textbf{posesión}.

\begin{definition}{Tenencia y posesión}
Ambas son relaciones de poder que una persona puede tener respecto de una cosa. En la \textbf{posesión}, quien tiene la relación de poder lo hace a título de dueño: se comporta como dueño (\textit{animus domini}), excluyendo a los demás de ese carácter. En la \textbf{tenencia}, quien tiene la cosa lo hace en nombre de otro: reconoce en otro el carácter de poseedor, aunque detente materialmente la cosa. El locatario no recibe la posesión, sino la tenencia: recibe la facultad de usar y gozar de la cosa a cambio de un precio, de forma temporal.
\end{definition}

\begin{important}
El elemento diferencial entre poseer y tener no está en el hecho material de ocupar la cosa, sino en el \textbf{ánimo} con que se la ocupa. El dominio, como derecho real, requiere ambos elementos: el \textit{corpus} (la posesión material) y el \textit{animus domini} (el ánimo de dueño). Quien sostiene una cosa y se comporta como su dueño es, de hecho, poseedor; la otra parte podrá discutir luego si tiene un mejor derecho a esa posesión.
\end{important}

\subsection{La interversión del título}

Una relación que comenzó como mera tenencia puede, con el transcurso del tiempo y mediante actos externos inequívocos, mutar en posesión. En derechos reales esto se conoce como \textbf{intervertir el título}.

\begin{note}
La jurisprudencia sostiene que no basta la mera manifestación personal de que se ocupa la cosa. Quien comenzó como tenedor y luego pretende comportarse como poseedor debe haberlo manifestado mediante actos que impliquen esa nueva atribución: por ejemplo, excluir al dueño o no dejarlo entrar o, en un inmueble rural, marcar los límites, plantar árboles o cultivar el campo. Deben existir actos que no realiza quien es meramente tenedor y reconoce ajena la posesión.
\end{note}

\begin{example}{El locatario que se comporta como dueño}
Con el correr del tiempo, un locatario puede dejar de pagar el alquiler, no permitir el ingreso del propietario, modificar el inmueble sin autorización y negar el carácter de propietario a quien lo era. En ese supuesto la persona <<se dice poseedora>> y, si logra acreditar mínimamente esos actos, el desalojo deja de ser la vía adecuada.
\end{example}

\subsection{Desalojo y acciones posesorias: procesos distintos}

El desalojo es distinto de las acciones reales, de las acciones posesorias y de los interdictos regulados en el Código Procesal: no tutelan el mismo derecho.

\begin{important}
La acción de despojo parte de un derecho real: hay señorío, hay \textit{animus domini} más tradición de la cosa. El desalojo, donde se busca recuperar la tenencia, parte de un \textbf{derecho personal}: la locación y el comodato generan una relación de derecho personal. En la tenencia el \textit{animus domini} no está, porque se reconoce el señorío en otra persona.
\end{important}

\textbf{Consecuencia procesal.} El desalojo es una acción personal que solo sirve para discutir la tenencia. Si el demandado reviste el carácter de poseedor, la posesión no puede discutirse por esta vía. Si el ocupante invoca y acredita ---aunque sea sumariamente--- actos posesorios (excluir al dueño, negarse a reconocerlo, dejar de pagar y modificar el inmueble como propio), corresponde discutir la posesión por la vía real o posesoria (interdictos, acción reivindicatoria). Antes de encuadrar la demanda es fundamental determinar con qué carácter ocupa el inmueble quien será demandado; ante la duda, conviene intimar previamente al ocupante mediante una diligencia preliminar, para conocer en qué calidad ocupa el bien.

\begin{example}{Comodato dentro de un conflicto sucesorio}
Una persona ocupa un inmueble perteneciente al acervo hereditario del padre de los titulares, habiendo sido admitida por uno de los herederos, sin pagar contraprestación, en el marco de un préstamo de uso gratuito. La parte contraria promueve una acción reivindicatoria. La figura aplicable es el comodato, y la vía adecuada frente al comodato vencido, o frente a la reclamación de restitución, sería en principio el desalojo por comodato vencido. Además, existía una obligación alimentaria del padre hacia la ocupante, lo que agregaba complejidad al conflicto.
\end{example}

\section{El contrato de locación}\label{sec:locacion}

\begin{normacitada}
\textbf{Código Civil y Comercial, artículo 1187 y siguientes} ---regulan el contrato de locación, definiendo las figuras del locador y del locatario y estableciendo sus caracteres, obligaciones y modo de extinción.
\end{normacitada}

La calidad de locador no se identifica necesariamente con la titularidad de dominio: la legitimación para peticionar el desalojo y recuperar la tenencia es más amplia que la mera titularidad registral. Para iniciar un desalojo no es necesario ser propietario; basta revestir el carácter de locador, incluso de un bien ajeno. Puede demandar incluso un locatario que, por una facultad de sublocar prevista en su contrato, resulta locador de un sublocatario que no le paga. La forma escrita exigida para el consentimiento se estudió en la sección~\ref{sec:firma}.

\subsection{Caracteres esenciales}

\begin{definition}{Caracteres del contrato de locación}
La locación es:
\begin{itemize}
\item \textbf{Onerosa}: una persona usa y goza del inmueble y otra permite ese uso y goce a cambio de un precio acordado; la gratuidad es propia del comodato.
\item \textbf{De plazo}, hoy librado en gran medida a la autonomía de la voluntad. Las normativas anteriores establecían plazos mínimos y máximos; actualmente quedan abiertos a que las partes los determinen libremente. En la práctica suele mantenerse la fórmula de dos años para vivienda y tres para destino comercial, aunque ya no constituya una exigencia legal indicativa.
\end{itemize}
\end{definition}

\begin{example}{Los gastos que integran la obligación de pago del locatario}
Además del canon locativo, integran la obligación de pago los servicios de luz, gas y agua, las expensas ---ordinarias o, según lo pactado, completas--- y el impuesto de rentas. Estos rubros pueden integrar el reclamo tanto en el proceso de desalojo como en el proceso ejecutivo o de cobro de alquileres, y su exigibilidad surge de lo pactado, lo que exige la lectura completa del contrato.
\end{example}

\subsection{Incumplimiento y causales de desalojo}

El incumplimiento que habilita el desalojo no se limita al impago del canon, sino que puede derivar de cualquier obligación pactada.

\begin{note}
El objeto del proceso puede ser desalojo por vencimiento de contrato, por cambio de destino, entre otras causales, agrupadas administrativamente bajo la carátula genérica de <<otras causales>>. Según esta fuente, el incumplimiento del pago del alquiler o de las expensas se canaliza por la vía de ejecución o cobro de alquileres y no por la de desalojo.
% TODO: verificar consistencia entre fuentes originales. Otras clases (31/08, 03/09, 24/09) desarrollan el desalojo por falta de pago como causal; la clase del 12/08 indica que el impago se canaliza por la vía de cobro o ejecución de alquileres. Se conservan ambas posiciones.
\end{note}

\subsection{Cesión y sublocación}

También deben considerarse la cesión y la sublocación, verificando qué se pactó al respecto en el contrato. La práctica contractual suele incluir una cláusula que prohíbe la sublocación y la cesión, aunque lo efectivamente practicado por las partes puede diferir de lo pactado formalmente.

\section{El contrato de comodato}

\begin{normacitada}
\textbf{Código Civil y Comercial, artículo 1533 y siguientes} ---regulan el contrato de comodato, definiendo las figuras del comodante y del comodatario y estableciendo la gratuidad como carácter esencial.
\end{normacitada}

\begin{definition}{Comodato}
Es una figura afín a la locación, pero estructuralmente distinta: sus sujetos son el \textbf{comodante} y el \textbf{comodatario}, y se otorga el uso y goce de la cosa sin ninguna prestación de canon locativo. Es gratuito en cuanto al uso, aunque puede pactarse que el comodatario asuma los servicios del inmueble.
\end{definition}

Como en la locación, puede promoverse el desalojo por comodato vencido, y el contrato exige en principio la forma escrita. El comodato suele celebrarse entre personas unidas por una relación de confianza ---familiares, amigos--- y, a diferencia de la locación, no reconoce plazos mínimos ni máximos legales: rige plenamente la autonomía de la voluntad.

\begin{important}
El Código de Vélez admitía el \textbf{comodato precario}, es decir, el pactado verbalmente. El CCCN vigente desde 2015 no reflejó esta figura: aunque el conflicto exista en la práctica, el préstamo de uso debe probarse con otros medios si no se cuenta con el instrumento escrito.
\end{important}

\begin{table}[H]
\centering
\caption{Comparación entre locación y comodato}
\begin{tabularx}{\textwidth}{>{\raggedright\arraybackslash}p{0.20\textwidth} X X}
\toprule
\textbf{Aspecto} & \textbf{Locación} & \textbf{Comodato} \\
\midrule
Sujetos & Locador y locatario. & Comodante y comodatario. \\
Contraprestación & Onerosa: precio (canon locativo). & Gratuito; puede pactarse el pago de servicios. \\
Forma & Escrita, exigida para inmuebles. & En principio escrita; ya no existe el comodato precario verbal. \\
Plazos legales & Antes con mínimos y máximos; hoy librados a la autonomía de la voluntad. & Sin mínimos ni máximos; plena autonomía de la voluntad. \\
Causal de desalojo & Vencimiento, falta de pago, cambio de destino, otras. & Comodato vencido. \\
\bottomrule
\end{tabularx}
\end{table}

\begin{example}{Comodato derivado de una compraventa}
El comodato puede derivar de una compraventa inmobiliaria. Si dos personas celebran una escritura pública de compraventa traslativa de dominio, el comprador puede permitir al vendedor permanecer en el inmueble por un plazo breve ---por ejemplo, mientras finaliza la mudanza---, configurándose así un comodato posterior a la compraventa.
\end{example}

\section{Extinción anticipada y alternativas al litigio}

El convenio de desocupación, cláusula de uso frecuente en el pasado, no se encuentra derogado; sin embargo, existen hoy otras vías para pactar la salida del inmueble. En particular, las partes pueden, de común acuerdo, renunciar al plazo pactado y hacer la restitución anticipada.

\begin{example}{Renuncia anticipada al plazo pactado}
Un contrato con plazo original de dos años, que vencería por ejemplo en marzo de determinado año, puede ser objeto de una renuncia anticipada de ese plazo por ambos contratantes, de común acuerdo, permitiendo la entrega anticipada del inmueble.
\end{example}

\begin{important}
Puede pactarse, por ejemplo, que la persona permanezca dos meses sin abonar el canon y, a cambio, entregue la tenencia, documentando el acuerdo ante un mediador para que tenga fuerza de acuerdo homologado. Los desalojos son procesos muy extensos; aun cuando con un arreglo pueda perderse algún dinero, suele ganarse en tiempo y en desgaste evitado. Agotar las vías alternativas ---incluida la mediación--- es una buena práctica profesional antes de acudir al Poder Judicial.
\end{important}

\begin{cornell}{Cuadro Cornell del capítulo~\ref{cap:titulos}}
\cqa{¿Qué diferencia a la tenencia de la posesión?}{La posesión ejerce el poder sobre la cosa con \textit{animus domini}; la tenencia lo ejerce reconociendo la posesión ajena. El locatario recibe la tenencia, no la posesión.}
\cqa{¿Qué es intervertir el título y cómo se acredita?}{Es mutar de tenedor a poseedor mediante actos externos inequívocos (por ejemplo, excluir al dueño); no basta la mera manifestación personal.}
\cqa{¿Por qué importa distinguir poseedor de tenedor antes de demandar?}{Porque el desalojo, acción personal, solo discute la tenencia; si el ocupante acredita sumariamente actos posesorios corresponde la vía real o posesoria (interdictos, reivindicación).}
\cqa{¿Quién puede demandar el desalojo en cuanto al carácter de locador?}{No hace falta ser propietario: basta ser locador, aun de un bien ajeno, o locatario que sublocó a quien no le paga.}
\cqa{¿Qué caracteres esenciales tiene la locación?}{Es onerosa, de plazo hoy librado a la autonomía de la voluntad y de forma escrita para inmuebles; el pago comprende canon más servicios, expensas e impuestos según lo pactado.}
\cqa{¿En qué se diferencia el comodato de la locación?}{Es gratuito, sin plazos legales y sin comodato precario verbal desde el CCCN; su causal de desalojo es el comodato vencido.}
\cqa{¿Cómo puede extinguirse anticipadamente la locación?}{Por renuncia de común acuerdo al plazo pactado y restitución anticipada, documentada preferentemente ante un mediador.}
\cresumen{El desalojo discute la tenencia y parte de un derecho personal nacido de un título, principalmente la locación (onerosa) o el comodato (gratuito). Antes de litigar conviene verificar el carácter del ocupante y considerar acuerdos de extinción anticipada, que suelen ahorrar tiempo y desgaste.}
\end{cornell}

\part{Sujetos del proceso y legitimación}

\chapter{Partes y legitimación en el desalojo}\label{cap:partes}

Quien promueve o enfrenta un desalojo debe poder responder dos preguntas previas al fondo: quién puede demandar y contra quién, y con qué aptitud comparece cada uno en el proceso. Este capítulo distingue la legitimación sustancial de la capacidad procesal y de la personería, expone quiénes tienen legitimación activa y pasiva, y precisa cómo se plantean las excepciones correspondientes. Sobre esa base, los capítulos siguientes estudian los órganos de la jurisdicción (capítulo~\ref{cap:organos}) y los terceros que pueden estar presentes (capítulo~\ref{cap:terceros}).

\section{Más allá de la tríada actor, demandado y juez}\label{sec:triada}

Existe una tendencia a concebir el proceso de manera reducida, limitándolo a la tríada actor--demandado--juez. Frente a esa simplificación corresponde manejar un esquema más amplio. Los sujetos del proceso de desalojo pueden organizarse en tres categorías: los órganos de la jurisdicción, los auxiliares del juez, y las partes y sus colaboradores. Además de las partes pueden intervenir sujetos extrajudiciales: el mediador o la mediadora, cuya intervención en el desalojo es optativa y no obligatoria, y el árbitro, en el marco de un sistema adversarial de resolución del conflicto.

Atraviesan todos los procesos de desalojo y de ejecución de alquileres, además, los \textbf{subinquilinos} y los \textbf{ocupantes} (véase el capítulo~\ref{cap:terceros}).

\section{Legitimación sustancial, capacidad procesal y personería}\label{sec:legitimacion}

\begin{definition}{Legitimación sustancial y capacidad procesal}
La \textbf{legitimación sustancial} (para obrar) es la posibilidad que la ley otorga a una persona para reclamar una determinada pretensión: refiere a ser el titular de la relación jurídica sustancial que se pretende hacer valer. En la terminología clásica equivale a la capacidad de derecho. En un contrato de locación corresponde a quien figura como locador o locatario, y no a un tercero ajeno a esa relación.

La \textbf{capacidad procesal} (equivalente a la capacidad de hecho o de ejercicio) refiere a la aptitud para estar por sí mismo en el proceso. La \textbf{personería} es la aptitud de esa persona para comparecer válidamente, ya sea por capacidad propia o por representación.
\end{definition}

Un sujeto debe reunir, entonces, dos condiciones: capacidad de derecho (legitimación sustancial, ser titular de la relación) y capacidad de hecho, para ejercer esa titularidad dentro del proceso. Una persona puede ser titular sustancial del derecho y carecer de capacidad de ejercicio, caso en el cual debe comparecer a través de su representante legal. Una cosa es la \emph{calidad de parte} y otra la \emph{legitimidad}: un menor de edad no tiene calidad de parte porque no puede actuar por sí mismo, pero puede tener legitimidad sustancial ---por ejemplo, si el inmueble le fue donado por sus abuelos---.

\begin{important}
Las fuentes designan de dos modos la deficiencia en la capacidad o representación: como \textbf{falta de personería} y como falta de legitimación procesal. Ambas se distinguen de la \textbf{falta de legitimación} (activa o pasiva), que indica que quien reclama, o es demandado, no es el sujeto que la ley habilita o la relación sustancial involucra. Cuando falta capacidad o personería se opone la excepción de falta de personería; cuando falta la legitimación, la de falta de legitimación activa o pasiva. Son excepciones distintas, aunque en la práctica a veces se las denomine, de manera imprecisa, con el mismo nombre.
\end{important}

\begin{formula}{Artículo 24 del CCCN ---personas incapaces de ejercicio}
Son incapaces de ejercicio: a) la persona por nacer; b) la persona que no cuenta con la edad y grado de madurez suficiente, con el alcance dispuesto en la Sección 2.ª de ese Capítulo; c) la persona declarada incapaz por sentencia judicial, en la extensión dispuesta en esa decisión.
\end{formula}

\begin{example}{El locatario declarado insano}
Un locatario que, después de firmar el contrato, es declarado insano por sentencia judicial, conserva su legitimación sustancial (sigue siendo titular del contrato) pero pierde la capacidad procesal: debe comparecer su curador. Puede haber un curador provisorio mientras se sustancia el proceso de declaración de incapacidad, hasta que el Cuerpo Médico Forense produzca sus informes. Se señaló que este es el único caso que debe subir en consulta a la Cámara, aunque no se haya apelado, para confirmar la decisión de quien declaró la insania. Si el curador propuesto tiene intereses contrapuestos con la persona incapaz ---por ejemplo, si es heredero en una sucesión en la que también participa la persona bajo curatela---, corresponde designar a otro representante por existir un conflicto de intereses.
\end{example}

\subsection{Formas de representación y control de los poderes}

Se distinguen la \textbf{representación legal} (los padres o, si se trata de una persona con capacidad restringida, el curador) y la \textbf{representación convencional}, propia de la relación abogado--cliente mediante apoderamiento. Esta última puede instrumentarse por instrumento privado con firma certificada o por escritura pública, conforme a la mayor flexibilidad introducida por el CCCN desde 2015 (art.\ 1017 CCCN; véase la sección~\ref{sec:firma}).

Es necesario revisar con cuidado los poderes acompañados en el expediente, porque en ocasiones se presenta un apoderamiento que no habilita la representación pretendida ---por ejemplo, un poder otorgado para un proceso laboral, que nada tiene que ver con la representación convencional en un desalojo---. Como el tribunal no siempre controla esta cuestión de oficio, corresponde a los abogados de cada parte verificarla y, eventualmente, plantear la excepción. Cuando la parte actúa por derecho propio no corresponde acompañar instrumento de personería; si se actúa con apoderado y se omite acreditarla, la contraparte puede oponer la excepción de falta de personería, que corre con las costas del incidente.

\begin{example}{El menor propietario y el tío administrador}
Un menor de edad es titular de una propiedad heredada de sus padres, que se encuentra alquilada. Su tío, sin titularidad ni representación legal, promueve por sí mismo una demanda de desalojo contra el inquilino por falta de pago. Si en el contrato figura como locador el menor, el legitimado activo es su padre o madre, su representante legal, o el tutor, según corresponda. Aunque el tío haya administrado de hecho el departamento, no tiene legitimación activa: la ley no le reconoce aptitud para reclamar en ese carácter. Podría actuar como abogado patrocinante, pero no por derecho propio, salvo que el representante legal le otorgue poder. El menor, por su parte, no podría demandar por sí mismo, salvo que estuviera emancipado.
\end{example}

\begin{example}{Poder mal redactado: falta de personería, no de legitimación}
Si el tío del ejemplo anterior es abogado y actúa con un poder mal redactado, no se configura falta de legitimación sino \textbf{falta de personería}: existe una deficiencia en la personería o en la capacidad de quien se presenta. Si en cambio interpone la demanda invocando derecho propio, sin ser locador, ni propietario, ni tener poder válido, el demandado puede oponer la excepción de falta de legitimación activa, porque el sistema no lo reconoce como sujeto habilitado.
\end{example}

\begin{table}[H]
\centering
\caption{Falta de personería y falta de legitimación pasiva para obrar}
\begin{tabularx}{\textwidth}{>{\raggedright\arraybackslash}p{0.24\textwidth} >{\raggedright\arraybackslash}p{0.37\textwidth} X}
\toprule
\textbf{Excepción} & \textbf{Qué cuestiona} & \textbf{Ejemplo} \\
\midrule
Falta de personería & Que quien actúa (letrado o representante) no acreditó debidamente su representación o poder. & El abogado de los consultantes no acompaña el poder de dos de los coherederos junto con la demanda. \\
Falta de legitimación pasiva para obrar & Que el demandado no tiene ninguna relación jurídica con el objeto del litigio. & Se demanda a un amigo del locatario que solo visita la casa, pero no la ocupa ni la usa. \\
\bottomrule
\end{tabularx}
\end{table}

\section{Legitimación activa}\label{sec:leg-activa}

La ley no limita la posibilidad de demandar por desalojo al propietario ni exclusivamente al locador. El CPCCN no establece expresamente quiénes tienen legitimación activa, a diferencia de lo que ocurre con la pasiva.

\begin{definition}{Legitimación activa amplia}
No se limita al titular dominial (quien tiene \textit{animus domini} y tradición), sino que se extiende a todo aquel que ostente el derecho al uso y goce de la cosa: el locador; el locatario que tiene un sublocatario; el sublocador; el comodante cuyo préstamo venció; el condómino; el heredero; el administrador de la asociación o consorcio; el representante; el usuario; el usufructuario y el tenedor. Incluso el poseedor puede demandar cuando exista una persona con menos derecho que él a la tenencia que se haya introducido en el inmueble.
\end{definition}

No hace falta ser propietario o titular de dominio: basta revestir el carácter de locador, aun cuando se locara un bien ajeno. Cuando hay cotitulares, no importa cuál de ellos firmó el contrato sino su condición de propietario; en el contrato puede verificarse qué parte corresponde a cada cotitular, y a quien reviste esa condición se lo denomina locador.

\subsection{Acreditación frente a intrusos}

Cuando se demanda a un intruso ---persona que ocupó el inmueble sin título alguno que la vincule con el actor--- no existe contrato que sirva de título habilitante. Quien demanda debe acreditar su mejor derecho sobre la cosa, ya sea invocando el dominio o una posesión anterior. No basta el informe de dominio del Registro de la Propiedad Inmueble, porque acredita el título pero no la posesión efectiva: el dominio requiere \textit{corpus} más \textit{animus domini}. La escritura pública, en cambio, permite acreditar la posesión, en tanto hace constar que hubo tradición del bien, aunque sea una manifestación del escribano y no un hecho comprobado por el tribunal.

\begin{table}[H]
\centering
\caption{Acreditación de la posesión frente a intrusos: provincia de Buenos Aires y orden nacional}
\begin{tabularx}{\textwidth}{>{\raggedright\arraybackslash}p{0.22\textwidth} X X}
\toprule
\textbf{Aspecto} & \textbf{Departamentos Judiciales de la provincia de Buenos Aires} & \textbf{Fuero civil nacional} \\
\midrule
Exigencia probatoria & Se exige acreditar que el actor tuvo, en algún momento, la posesión efectiva, aun cuando hoy reclame solo la tenencia. & No se exige posesión anterior; basta acompañar el título (por ejemplo, el contrato de locación) y acreditar el carácter invocado. \\
Riesgo procesal & La demanda puede ser rechazada si no se acredita esa posesión pretérita. & Menor exigencia probatoria en este punto. \\
\bottomrule
\end{tabularx}
\end{table}

Esta exigencia bonaerense debe tenerse presente al redactar los escritos postulatorios en esa jurisdicción. En la práctica, el apartado de legitimación activa de la demanda es optativo: se aclara como estrategia, para anticiparse a la excepción dilatoria de falta de legitimación que suele oponer el demandado (véase el capítulo~\ref{cap:demanda}).

\subsection{Vías civil y penal frente a intrusos}

Contra los intrusos existen dos vías posibles, la civil y la penal, aunque en la práctica la penal resulta poco efectiva. No se excluyen entre sí, pero cada una tramita por un fuero distinto, y en la práctica suele optarse por una sola. El desalojo inmediato, propio de la vía civil, se desarrolla en el capítulo~\ref{cap:vias}.

\section{Legitimación pasiva}\label{sec:leg-pasiva}

A diferencia de la activa, la legitimación pasiva sí se encuentra prevista expresamente en el Código Procesal.

\begin{formula}{Artículo 680 del CPCCN ---legitimación pasiva}
La acción de desalojo puede dirigirse contra todo aquel que use u ocupe el inmueble, aun sin derecho, cuando su título es meramente precario o si carece de título, así como contra los sucesores de ese ocupante, los sublocatarios y los ocupantes de cualquier especie. Comprende al locatario, al sublocatario, al intruso y a toda persona que tenga la obligación de restituir la cosa cuya tenencia se reclama.
\end{formula}

La legitimación pasiva se dirige, en síntesis, contra todo aquel que no tenga derecho al uso y goce de la cosa y, sin embargo, la use u ocupe. En un juicio de desalojo corresponde demandar no solo al locatario formal, sino también a los subinquilinos y a todo ocupante real del inmueble, aunque no figure como parte firmante del contrato. Lo contrario ---demandar a quien solo visita la casa sin ocuparla--- expone a la excepción de falta de legitimación pasiva.

\begin{sintesis}
La legitimación sustancial exige ser titular de la relación jurídica; la capacidad procesal y la personería, poder comparecer por sí o por representante. Para demandar basta el uso y goce de la cosa; se demanda a quien use u ocupe sin derecho, incluidos subinquilinos y ocupantes. La falta de personería y la de legitimación son excepciones distintas.
\end{sintesis}

\begin{cornell}{Cuadro Cornell del capítulo~\ref{cap:partes}}
\cqa{¿Qué diferencia hay entre legitimación sustancial y capacidad procesal?}{La legitimación sustancial es ser titular de la relación jurídica que se hace valer (capacidad de derecho); la capacidad procesal es poder estar por sí en el proceso (capacidad de hecho). Quien carece de esta última comparece por su representante legal.}
\cqa{¿Cuándo se opone falta de personería y cuándo falta de legitimación?}{Falta de personería, ante deficiencia en la capacidad o la representación (por ejemplo, un poder mal redactado). Falta de legitimación, cuando quien reclama o es demandado no es el sujeto habilitado por la relación sustancial.}
\cqa{¿Quiénes pueden demandar el desalojo?}{Todo aquel que ostente el derecho al uso y goce: locador, locatario con sublocatario, sublocador, comodante, condómino, heredero, usufructuario, usuario, tenedor y aun el poseedor frente a quien tenga menos derecho. No se exige ser propietario.}
\cqa{¿Qué exige la provincia de Buenos Aires frente a intrusos y qué el fuero nacional?}{Los Departamentos Judiciales bonaerenses exigen acreditar posesión efectiva anterior; el fuero civil nacional basta con acompañar el título invocado.}
\cqa{¿Contra quién se dirige la acción según el art.\ 680 CPCCN?}{Contra todo aquel que use u ocupe el inmueble sin derecho, título precario o sin título, sus sucesores, sublocatarios y ocupantes de cualquier especie.}
\cqa{¿Por qué conviene demandar a subinquilinos y ocupantes?}{Porque todo ocupante real debe ser demandado y notificado para resguardar la defensa en juicio, aunque no esté individualizado ni firmara el contrato.}
\cresumen{Legitimarse exige titularidad de la relación sustancial y aptitud para comparecer. La legitimación activa es amplia y la pasiva alcanza a todo ocupante sin derecho. Confundir legitimación con personería lleva a oponer o contestar la excepción equivocada.}
\end{cornell}

\chapter{Órganos de la jurisdicción, auxiliares y colaboradores}\label{cap:organos}

Los sujetos que intervienen en un desalojo, además de las partes, desempeñan funciones diferenciadas. Este capítulo recorre el juez o jueza y sus límites (conciliación intraprocesal y recusación), los auxiliares del juez (Ministerio Público, oficial notificador y oficial de justicia) y los colaboradores de las partes (patrocinio, apoderamiento y consultores técnicos). El capítulo~\ref{cap:terceros} estudia luego a los terceros vinculados con el inmueble.

\section{El juez o jueza: deberes y facultades}

Los deberes y facultades del juez se ubican en el CPCCN, particularmente en los artículos 34, 35 y 36.

\begin{important}
Los artículos 34 y 35 del CPCCN regulan los deberes y facultades del juez, incluida su facultad de disponer, en cualquier momento del proceso, la comparecencia personal de las partes para intentar una conciliación. El artículo 36 ---la denominada <<audiencia del artículo 36>>--- habilita al juez a convocar esta audiencia en cualquier etapa del proceso, de oficio o a petición de parte, a diferencia de otras audiencias que solo pueden fijarse en momentos procesales determinados.
\end{important}

Esta conciliación intraprocesal se da muy pocas veces en el desalojo: cuando el demandado solicita la audiencia, el juzgado en general no la concede, porque el instituto suele utilizarse como herramienta dilatoria antes que con fines genuinamente conciliatorios, en mayor medida que en otros procesos de conocimiento. Una negociación genuina suele darse fuera del tribunal: los letrados coordinan directamente entre sí, por videollamada o en algún estudio jurídico, sin depender del tribunal.

\subsection{Recusación con y sin expresión de causa}

\begin{note}
El artículo 17 del CPCCN regula la recusación con expresión de causa, cuando existe un motivo concreto ---como amistad o enemistad con alguna de las partes--- que justifica apartar al juez de la causa.
\end{note}

\begin{normacitada}
\textbf{Artículo 14 del CPCCN ---recusación sin expresión de causa.} Permite recusar a los jueces de primera instancia sin expresar la causa. El actor puede ejercer esta facultad al entablar la demanda o en su primera presentación; el demandado, en su primera presentación, antes o al tiempo de contestarla, u oponer excepciones en el juicio ejecutivo, o al comparecer a la audiencia fijada como primer acto procesal. El último apartado, incorporado por la Ley 25.488 (B.O.\ 22/11/2001), dispone que no procede en el proceso sumarísimo, en las tercerías, en el juicio de desalojo y en los procesos de ejecución, ampliando la prohibición que antes alcanzaba al juicio ejecutivo.
\end{normacitada}

La prohibición se aplica de manera directa y expresa al desalojo, por el texto literal del artículo. Se planteó un debate conexo: si la prohibición respecto de los <<procesos de ejecución>> alcanza también a la ejecución de alquileres. Si se debe alquileres u otros rubros hay que iniciar otro juicio, la ejecución o cobro de alquileres, que es un proceso distinto del desalojo. Aunque el texto del artículo 14 no menciona expresamente esa ejecución, tanto la Cámara Nacional de Apelaciones en lo Comercial en pleno (plenario del 27 de mayo de 2003, <<Recusación sin causa en juicios ejecutivos s/convocatoria a plenario>>) como la Cámara Nacional de Apelaciones en lo Civil, con criterio análogo, resolvieron que la prohibición alcanza a los juicios ejecutivos, categoría en la que se comprendió la ejecución de alquileres.

\begin{example}{La recusación como estrategia dilatoria}
La recusación sin causa se ejercita electrónicamente, con la sola presentación en el sistema informático, y puede utilizarse para dilatar el proceso, ya que quien recusa consigue que el expediente vuelva al sistema de sorteo o asignación informática. El límite del último apartado del artículo 14 busca evitar esa dilación en procesos que, por su naturaleza, requieren celeridad. Aunque los abogados de la contraparte a veces intenten igualmente la recusación sin causa en estos procesos, el juzgado debe rechazarla en virtud de los plenarios citados.
\end{example}

\section{Los auxiliares del juez}

Los auxiliares del juez son los sujetos que colaboran con el órgano jurisdiccional para llevar adelante el proceso.

\subsection{Ministerio Público}

Dentro de esta institución interesa especialmente el Ministerio Público de la Defensa, que interviene cuando están presentes personas en situación de vulnerabilidad: niños, niñas y adolescentes, personas con capacidad restringida o adultos mayores (véase la sección~\ref{sec:vulnerabilidad}). Cuando alguno de ellos está involucrado, corresponde correr vista al Ministerio Público para que emita su dictamen, y recién luego el expediente vuelve al juez para resolver.

\begin{example}{Desalojo con menores en el inmueble (Ciudadela, partido de Tres de Febrero)}
Un proceso de desalojo tramitaba en el fuero nacional civil respecto de un inmueble ubicado en Ciudadela, partido de Tres de Febrero, provincia de Buenos Aires. Al llegar el lanzamiento, el oficial de justicia del departamento judicial de Tres de Febrero advirtió, en el recorrido previo al inmueble realizado antes de fijar fecha, la presencia de menores de edad. Como el proceso tramitaba en el fuero nacional civil pero el inmueble estaba en jurisdicción provincial, la solución consistió en abrir un oficio bajo la Ley 22.172 dirigido al juzgado civil y comercial de San Martín, para que habilitara la intervención y la consulta al Ministerio Público de la Defensa. El trasfondo era un locador de edad avanzada, viudo y sin hijos, que vivía de su jubilación y del ingreso del alquiler.
\end{example}

\begin{important}
Verificar quiénes ocupan efectivamente el inmueble ---incluidos menores, personas con capacidad restringida o adultos mayores--- no es un trámite formal secundario: puede modificar el curso del proceso, obligando a articular mecanismos interjurisdiccionales (como la Ley 22.172) y a dar intervención al Ministerio Público de la Defensa antes de concretar el lanzamiento.
\end{important}

\subsection{Oficial notificador}

Es el auxiliar encargado de llevar la notificación de una resolución judicial. Su competencia territorial es relevante: debe pertenecer al departamento judicial correspondiente al domicilio donde se practica la notificación, que es el domicilio real que se denuncie o se corrobore mediante el Registro Nacional de las Personas, y no necesariamente el del tribunal. Si la persona a notificar reside en un departamento judicial distinto ---por ejemplo, en un partido bonaerense y no en la Ciudad Autónoma---, no corresponde el oficial notificador de la Ciudad, sino la \textbf{cédula ley}.

\begin{important}
La Ley 22.172 (convenio de comunicaciones entre tribunales) regula la comunicación entre tribunales de distintas jurisdicciones mediante la denominada <<cédula ley>>, que se utiliza cuando la persona a notificar tiene domicilio en un departamento judicial distinto de aquel donde tramita el proceso.
\end{important}

Para diligenciarla puede llevarla el propio letrado, alguien del estudio, o un gestor que cobra un arancel; muchas veces es el cliente quien asume el costo. En el desalojo deben notificarse también, en muchos casos, los fiadores o garantes, como obligados principales pagadores (véase la sección~\ref{sec:fiador}).

Cuando el oficial notificador no logra notificar al demandado en el domicilio denunciado ---por ejemplo, porque se niega su presencia o la numeración no coincide---, corresponde acreditar que la persona efectivamente reside allí y solicitar que se notifique bajo responsabilidad de la parte actora; si nuevamente se niega la residencia, el oficial fija la cédula en la puerta de acceso. Para reducir el margen de error conviene adjuntar a la cédula fotografías y un croquis detallado con referencias visuales del entorno inmediato.

\subsection{Oficial de justicia y lanzamiento}

El oficial de justicia se diferencia con claridad del oficial notificador. Su tarea principal es cumplir la diligencia del lanzamiento: la desocupación efectiva del inmueble respecto de quienes lo ocupan, sean demandados, subinquilinos u ocupantes. Es una figura contratada por el poder judicial local de cada departamento judicial, que actúa por orden del juez ---nunca interviene el juez en persona ni el secretario--- y que puede requerir, según lo peticionado y ordenado, el auxilio de la fuerza pública, de un cerrajero, de una ambulancia o incluso de un servicio veterinario cuando hay animales en el inmueble.

\begin{table}[H]
\centering
\caption{Oficial notificador y oficial de justicia}
\begin{tabularx}{\textwidth}{>{\raggedright\arraybackslash}p{0.22\textwidth} >{\raggedright\arraybackslash}X >{\raggedright\arraybackslash}X}
\toprule
\textbf{Criterio} & \textbf{Oficial notificador} & \textbf{Oficial de justicia} \\
\midrule
Función principal & Lleva la notificación de resoluciones judiciales (cédulas). & Ejecuta el lanzamiento: hace efectiva la desocupación. \\
Competencia territorial & Departamento judicial del domicilio a notificar. & Departamento judicial donde tramita el lanzamiento. \\
Puede requerir & --- & Fuerza pública, cerrajero, ambulancia, servicio veterinario (si se peticiona y el juez lo ordena). \\
Relación con el flete & No interviene. & Actúa junto con el fletero, contratado aparte. \\
\bottomrule
\end{tabularx}
\end{table}

El flete no es una condición procesal sino una persona ---el fletero--- contratada aparte del oficial de justicia y del cerrajero. Es necesaria porque, al abrirse la puerta, los bienes que allí se encuentren no pueden amontonarse junto a la puerta, sino que deben trasladarse a un depósito, notificando al ocupante, con un plazo mínimo de un mes, que sus bienes se encuentran allí; si no se presenta a retirarlos, se dispone su entrega a una institución de bien público. El lanzamiento se estudia en detalle en el capítulo~\ref{cap:lanzamiento}.

\begin{note}
Como estrategia frente al lanzamiento conviene requerir la totalidad de los auxiliares disponibles (cerrajero, fletero, fuerza pública), porque muchas veces no puede preverse la situación dentro del inmueble. Debe explicarse al cliente que, si no se contratan, el lanzamiento puede no prosperar en la fecha fijada y reprogramarse varios meses después, con costo adicional.
\end{note}

\section{Colaboradores del proceso}

\subsection{Patrocinio letrado y apoderamiento}

El abogado puede actuar como patrocinante letrado, siempre con la firma de la parte, o como apoderado.

\begin{important}
Si un escrito se presenta bajo patrocinio letrado pero sin la firma de la parte, se tiene por no presentado y el juzgado lo desglosa.
\end{important}

El apoderado puede ser general judicial o especial para el desalojo en particular. Las fuentes presentan dos momentos normativos y una práctica:

\begin{important}
El \textbf{artículo 47 del CPCCN} exigía escritura pública (<<escritura poder>>) para el poder judicial. El \textbf{artículo 1017 del CCCN} (vigente desde agosto de 2015) flexibiliza la exigencia: según una de las fuentes admite que el mandato se otorgue bajo cualquier forma; según otra fuente, que se otorgue por instrumento privado o por instrumento privado con firma certificada. En la práctica procesal suele mantenerse la escritura pública para el poder judicial. En todo caso el poder se acompaña con la primera presentación (véase la sección~\ref{sec:firma}).
% TODO: verificar consistencia entre fuentes originales. Las clases del 31/08 y del 10/09 describen de manera distinta el alcance de la flexibilización del poder (<<cualquier forma>> frente a <<instrumento privado o con firma certificada>>).
\end{important}

\subsection{Consultores técnicos}

El consultor técnico, también llamado perito de parte, es una figura distinta del perito de oficio. Resulta especialmente útil cuando la liquidación del canon locativo ---que suma alquileres, intereses, servicios e impuestos--- alcanza montos elevados y comienza a controvertirse: conviene entonces acompañarse de un consultor técnico, con formación en matemáticas o en herramientas de análisis, para no afrontar en soledad esa tarea.

\begin{example}{Falta de legitimación activa del fiador}
En un caso de la práctica, una demanda de ejecución de alquileres fue promovida por el fiador ---y no por el locador---, luego de que aquel pagara la deuda y pretendiera accionar contra la locataria. Se planteó si correspondía oponer la excepción de falta de legitimación activa, porque quien inicia la ejecución no reviste, en principio, el carácter de acreedor directo de la relación locativa. Cabe distinguir lo técnico-jurídico de la realidad del ejercicio profesional, en la que un planteo débil puede utilizarse con fines dilatorios: si el tribunal lo admite y la contraparte no lo cuestiona a tiempo, el acto queda firme. Esta demora estructural motivó llamados de atención de la Corte Interamericana de Derechos Humanos hacia la Argentina en materia de plazo razonable.
\end{example}

\begin{sintesis}
El juez dirige el proceso con los límites del artículo 14 (sin recusación sin causa en el desalojo); el Ministerio Público de la Defensa interviene ante personas vulnerables; el oficial notificador y el de justicia cumplen funciones y competencias distintas; y las partes cuentan con patrocinio, apoderados y consultores técnicos.
\end{sintesis}

\begin{cornell}{Cuadro Cornell del capítulo~\ref{cap:organos}}
\cqa{¿Qué es la audiencia del artículo 36 y por qué casi no se concede en el desalojo?}{Es una audiencia conciliatoria que el juez puede convocar en cualquier momento, de oficio o a pedido de parte. En el desalojo suele denegarse porque se interpreta que se pide con fines dilatorios más que conciliatorios.}
\cqa{¿Puede recusarse sin causa en el desalojo?}{No: el último apartado del art.\ 14 CPCCN (Ley 25.488) lo prohíbe en el desalojo y en los procesos de ejecución; los plenarios de 2003 extendieron el criterio a los juicios ejecutivos, incluida la ejecución de alquileres.}
\cqa{¿Cuándo interviene el Ministerio Público de la Defensa?}{Cuando hay niños, niñas y adolescentes, personas con capacidad restringida o adultos mayores; se corre vista para el dictamen antes de que el juez resuelva, y puede exigir oficios bajo la Ley 22.172 si el inmueble está en otra jurisdicción.}
\cqa{¿Qué diferencia al oficial notificador del oficial de justicia?}{El notificador lleva cédulas y depende del departamento judicial del domicilio a notificar; el oficial de justicia ejecuta el lanzamiento y puede requerir fuerza pública, cerrajero, ambulancia o veterinario.}
\cqa{¿Qué es la cédula ley?}{La notificación regulada por la Ley 22.172, usada cuando el domicilio del notificado está en un departamento judicial distinto del del proceso.}
\cqa{¿Qué ocurre con un escrito con patrocinio sin firma de la parte?}{Se tiene por no presentado y se desglosa.}
\cqa{¿Para qué sirve un consultor técnico?}{Para controvertir liquidaciones complejas de canon, intereses y servicios; es perito de parte, distinto del perito de oficio.}
\cresumen{Los órganos y auxiliares cumplen funciones diferenciadas, con límites como la prohibición de recusar sin causa. Verificar quiénes ocupan el inmueble puede activar al Ministerio Público y mecanismos interjurisdiccionales, y la forma de los escritos y poderes condiciona su validez.}
\end{cornell}

\chapter{Terceros, pluralidad de partes y el fiador}\label{cap:terceros}

Junto con las partes principales, el desalojo involucra casi siempre a personas que no firmaron el contrato pero ocupan el inmueble, o que garantizan su cumplimiento. Este capítulo estudia a los subinquilinos y ocupantes, la pluralidad de partes (litisconsorcio), y de manera especial al fiador: su régimen de responsabilidad, su extinción, su notificación y las medidas sobre sus bienes.

\section{Subinquilinos y ocupantes}

La figura de subinquilinos y ocupantes está siempre presente en la demanda de desalojo. Es habitual incluir una fórmula genérica que da cuenta de que se desconoce si existen subinquilinos u ocupantes distintos del demandado principal y que, si se los conociera, se los denunciaría para su identificación. Aunque el proceso identifique con precisión a la persona demandada (por ejemplo, la locataria), siempre debe notificarse también a los supuestos subinquilinos y ocupantes, aun cuando no estén individualizados.

\begin{definition}{Fórmula genérica <<subinquilinos y ocupantes>>}
Se admite esta imprecisión por una razón procesal: identificar a cada persona por nombre y apellido obligaría a notificar correctamente a cada una, y omitir a alguna comprometería la debida defensa en juicio. No corresponde, entonces, individualizar a cada conviviente: se demanda genéricamente contra el locatario y contra sus subinquilinos y ocupantes.
\end{definition}

El lanzamiento se dirige tanto contra el demandado identificado como contra todo tercero que ocupe el inmueble sin haber sido individualizado. Entre los terceros del proceso corresponde mencionar también a peritos y testigos.

\section{Litisconsorcio y proceso colectivo}

El \textbf{litisconsorcio} es la pluralidad de partes que puede presentarse del lado actor, del demandado o de ambos. En el desalojo puede darse, por ejemplo, cuando varios condóminos revisten en conjunto el carácter de locadores frente al locatario y sus ocupantes.

\begin{definition}{Litisconsorcio voluntario y necesario}
El \textbf{litisconsorcio voluntario} es la regla general: la pluralidad se da de manera facultativa. El \textbf{litisconsorcio necesario} se reserva a casos muy particulares en los que la ley exige demandar conjuntamente a todos los sujetos involucrados; esta modalidad no se presenta en el desalojo.
\end{definition}

\begin{example}{Litisconsorcio necesario en una acción de escrituración}
Si dos personas son condóminas de un mismo inmueble, quien pretende la escrituración debe demandar a ambas, porque cada una solo puede escriturar su propia proporción (por ejemplo, el cincuenta por ciento).
\end{example}

El \textbf{proceso colectivo} no tiene aplicación directa en materia de desalojo. Fue incorporado por la reforma constitucional de 1994 y, sin embargo, no cuenta todavía con una regulación procesal operativa integral; la Corte Suprema, mediante una acordada, avanzó en la creación de un registro de procesos colectivos, pero el instituto sigue sin desarrollo legislativo pleno. Son ejemplos ajenos al desalojo, útiles para no confundir figuras, los casos ambientales (como el vertido de desechos industriales en un curso de agua que contamina las napas de un barrio o una ciudad) y los reclamos colectivos por comisiones bancarias cobradas indebidamente en cajas de ahorro.

\section{El fiador}\label{sec:fiador}

\subsection{Su condición de parte y su notificación}

El \textbf{fiador es parte} y, cuando concurre, debe poder intervenir; la litis debe integrarse para poder llegar a un acuerdo con todas las partes. Su notificación, sin embargo, requiere precisiones.

\begin{important}
No existe obligación legal de notificar al fiador el inicio del desalojo: es una \textbf{buena práctica} que algunos juzgados adoptan según el criterio de cada tribunal. El <<debe>> es una obligación cuya inobservancia acarrea sanción; la notificación al fiador no está impuesta bajo sanción, sino que resulta útil, porque apenas notificado formalmente el fiador suele intentar un entendimiento antes de afrontar una condena.
\end{important}

Con independencia de la vía judicial suele intentarse contactar a la contraparte por todos los medios antes de iniciar el proceso, pero la notificación judicial formal al fiador tiende a generar una reacción distinta de la de los acercamientos informales. En la realidad judicial, si el fiador estuvo debidamente notificado y no concurrió a la audiencia, esta se celebra igual: la jueza manda notificar el acuerdo al fiador, que está presentado en el proceso, para que preste conformidad u oposición, y el acuerdo queda cerrado una vez vencido el plazo. Esto respeta los principios de defensa en juicio e igualdad de las partes: si el inquilino cumple y se provoca un desalojo, el fiador es en definitiva el más perjudicado.

\subsection{Alcance de la responsabilidad y extinción de la fianza}

\begin{formula}{Régimen de la fianza en la locación}
El fiador responde, en principio, hasta el vencimiento del plazo del contrato de locación. \textbf{Excepción}: si al vencer el contrato el locatario no restituye el inmueble, el locador puede mantener vigente la fianza más allá del vencimiento, siempre que intime fehacientemente al locatario a restituir y comunique esa intimación al fiador. Si el locador omite ambas cosas, la fianza cesa automáticamente al vencimiento y el fiador deja de responder por los períodos posteriores, cualquiera sea la deuda acumulada.
\end{formula}

% TODO: verificar contenido incompleto o ambiguo en las notas originales: la clase del 03/09 señala que, bajo el régimen anterior a la reforma del Código Civil y Comercial, si el fiador no rescindía expresamente su compromiso seguía respondiendo indefinidamente, pero remite a una clase anterior para el desarrollo completo de ese régimen previo.

\begin{formula}{Artículo 1225 del CCCN ---extinción de la fianza}
La responsabilidad del fiador se limita al plazo contractual pactado. Si el contrato de locación continúa más allá de ese plazo sin que se celebre un nuevo contrato o una renovación formal en la que el fiador vuelva a obligarse expresamente, la fianza no se renueva automáticamente junto con la locación de hecho.
\end{formula}

\begin{example}{Extinción de la fianza por falta de renovación formal}
Un contrato de locación por dos años, con el hermano del locatario como fiador, vence en diciembre de 2025. Como la relación entre locador y locatario es buena y este es un buen pagador, el locador le propone continuar de palabra ---actualizando el alquiler según la evolución del mercado--- sin renovar por escrito ni formalizar nuevamente la fianza. La locación de hecho continúa con pagos mensuales, pero sin la fianza renovada. Si luego el locatario deja de pagar y el locador inicia el desalojo y, en paralelo, la ejecución de la deuda, la posibilidad de demandar también al hermano fiador depende de que su abogado conozca y oponga oportunamente el artículo 1225 CCCN.
\end{example}

Si el abogado del fiador desconoce el límite del artículo 1225, puede terminar aconsejándole pagar, o contestar la demanda sin oponer la excepción correspondiente, asumiendo una responsabilidad de la que podría liberarse. Si lo conoce y ejerce a tiempo la carga de comparecer y defenderse, puede oponer la \textbf{excepción de falta de legitimación pasiva} invocando el artículo 1225, con fundamento en que la responsabilidad se limitaba al plazo original, vencido sin renovación formal. Corresponde también un rol activo del abogado del locador: si asesora a quien quiere demandar al fiador, debe advertirle honestamente el riesgo de que la excepción prospere, dejando a criterio del cliente si de todas formas desea afrontar la tasa de justicia.

\begin{important}
El artículo 1225 CCCN no depende de una cláusula contractual que extienda la fianza <<hasta que el locatario devuelva la tenencia al locador>>: si el locador no fue diligente en renovar el contrato por escrito ni en iniciar oportunamente el desalojo, el fiador puede liberarse de responder por el período posterior al vencimiento del plazo original. Hay que conocer la norma para aplicarla oportunamente.
% TODO: verificar consistencia entre fuentes originales. La clase del 03/09 admite que el locador mantenga la fianza tras el vencimiento mediante intimación fehaciente al locatario y comunicación al fiador; la clase del 31/08 sostiene que la fianza no se extiende de hecho si no hay renovación formal. Se conservan ambas formulaciones.
\end{important}

\begin{note}
Conviene hacer la intimación con antelación ---tres meses antes, o al menos dos meses antes del vencimiento---, contactando al inquilino para saber si habrá renovación o prórroga. Si la hay, el fiador debe volver a firmar con firma certificada, o directamente armarse un nuevo contrato certificado con nueva fianza.
\end{note}

\begin{example}{El fiador que terminó pagando la deuda}
En un caso real, el inquilino desapareció, dejó de pagar veinte meses y no quiso saber más nada. Se recurrió al fiador: se obtuvo una cautelar genérica ---no fue posible embargar directamente por no existir firma certificada---. El fiador, desesperado, terminó pagando las cuotas adeudadas, si bien lo que había estado mal era la conducta del inquilino, que dejó el inmueble sin avisar al locador y le cerró la puerta, quedando ese incumplimiento a cargo del fiador.
\end{example}

\subsection{Medidas cautelares sobre bienes del fiador}

Las medidas cautelares constituyen un subsistema autónomo del sistema procesal y requieren siempre reunir sus propios presupuestos ---verosimilitud del derecho y peligro en la demora--- con independencia del proceso principal.

\begin{table}[H]
\centering
\caption{Medidas cautelares sobre bienes del fiador según el tipo de proceso}
\begin{tabularx}{\textwidth}{>{\raggedright\arraybackslash}p{0.30\textwidth} X}
\toprule
\textbf{Proceso} & \textbf{¿Procede cautelar sobre bienes del fiador?} \\
\midrule
Desalojo & No corresponde: lo que se persigue es recuperar la tenencia del inmueble, no una prestación dineraria del fiador. \\
Ejecución de alquileres & Puede jugar el subsistema de medidas cautelares, aunque en la generalidad de los tribunales no se conceden hasta que la litis está trabada, es decir, hasta que la contraparte esté debidamente notificada, haya contestado o no. \\
\bottomrule
\end{tabularx}
\end{table}

Entre las alternativas se mencionan la inhibición general de bienes, cuando no hay bienes identificados del deudor, y más excepcionalmente la prohibición de contratar, útil por ejemplo si se advierte que el fiador está por enajenar el inmueble que constituye su garantía patrimonial.

\begin{example}{Simulación de una compraventa para eludir una fianza}
Un local comercial estaba alquilado por un canon mensual de varios millones de pesos, garantizado por una fianza. El inquilino dejó de pagar y se insolventó; el fiador, para evitar su obligación, transfirió simuladamente un local de su propiedad a su sobrino mediante un boleto de compraventa. Según sospechaba el acreedor, detrás habría un contradocumento ---un pacto por el cual el sobrino reconoce que la venta no es real--- que nadie exhibiría en el juicio. Si no hay forma de acceder a él, la simulación se prueba reuniendo varios indicios precisos y concordantes ---nunca uno solo--- que, interpretados en conjunto, permitan presumir el hecho oculto.
\end{example}

\subsection{Garantías locativas en la práctica}

La garantía puede haberse constituido varios años antes de que el conflicto se materialice, y en ese tiempo el garante puede haber transferido los bienes que la respaldaban. Por ello se recomienda solicitar un informe de dominio o de titularidad actualizado antes de aceptar la garantía y, eventualmente, al iniciar su ejecución.

\begin{definition}{Tipos de garantía mencionados}
\textbf{Fianza}: garantía personal por excelencia en las locaciones; el fiador se obliga a responder por las deudas del locatario, sin que la obligación siga a un bien determinado, a diferencia de una garantía real.
\textbf{Prohibición personal de disponer}: no debe confundirse con una hipoteca, que es una garantía real sobre un inmueble determinado y sigue a la cosa; la prohibición personal no tiene ese efecto de persecución.
\textbf{Inhibición general de bienes}: hoy se practica exclusivamente como medida cautelar ordenada por un juez, previo cumplimiento de verosimilitud del derecho y peligro en la demora, y sujeta a contracautela. La <<autoinhibición>> voluntaria existió hace más de un siglo, pero fue derogada y no se practica en la actualidad, sin proyecto legislativo para reintroducirla.
\end{definition}

Calificar como insolvencia fraudulenta la venta de un bien por el garante requeriría acreditarlo en sede penal, lo que suele implicar costos superiores al beneficio que se busca proteger. Como práctica de gestión contractual ---no como exigencia legal---, se recomienda exigir, además del depósito habitual de garantía (por lo general un mes de alquiler por año de contrato, o un mes fijo), una caución real adicional.

\begin{example}{La caución real como incentivo a la negociación}
En los contratos en que se pactó que el locatario o su fiador entreguen al locador una caución real de una suma determinada de dinero, rara vez se llega a juicio: las partes prefieren sentarse a negociar y resolver el conflicto ---incluso cuestiones menores, como el desperfecto de un calefón--- antes que perder la suma entregada en garantía. Su eficacia se apoya en el incentivo económico que genera en ambas partes para evitar la vía judicial.
\end{example}

\begin{sintesis}
Subinquilinos y ocupantes se demandan y notifican genéricamente; el litisconsorcio necesario no se da en el desalojo ni el proceso colectivo tiene aplicación directa; el fiador es parte, responde en principio hasta el vencimiento del contrato, y su fianza no se extiende sin intimación o renovación formal.
\end{sintesis}

\begin{cornell}{Cuadro Cornell del capítulo~\ref{cap:terceros}}
\cqa{¿Por qué se demanda genéricamente a subinquilinos y ocupantes?}{Porque individualizar a cada persona obligaría a notificar a todas, y omitir a alguna comprometería la defensa en juicio. El lanzamiento alcanza a todo ocupante, esté o no individualizado.}
\cqa{¿Hay litisconsorcio necesario en el desalojo?}{No: rige el litisconsorcio voluntario como regla general; el necesario se reserva a casos como la escrituración contra condóminos.}
\cqa{¿Es obligatorio notificar al fiador?}{No hay obligación legal bajo sanción; es una buena práctica de algunos juzgados que suele generar entendimientos antes de la condena. Si notificado no concurre, la audiencia se celebra igual y se le notifica el acuerdo.}
\cqa{¿Hasta cuándo responde el fiador?}{En principio hasta el vencimiento del contrato; puede extenderse si el locador intima fehacientemente al locatario a restituir y notifica al fiador. Según el art.\ 1225 CCCN, la fianza no se renueva con la locación de hecho sin renovación formal.}
\cqa{¿Qué excepción puede oponer el fiador cuando venció el plazo?}{La de falta de legitimación pasiva, fundada en el art.\ 1225 CCCN, siempre que se oponga oportunamente.}
\cqa{¿Procede una cautelar sobre bienes del fiador en el desalojo?}{No, porque se persigue la tenencia y no una prestación dineraria; en la ejecución de alquileres puede evaluarse, aunque suele concederse recién con la litis trabada.}
\cqa{¿Cómo se prueba una simulación sin el contradocumento?}{Con varios indicios precisos y concordantes, nunca con uno solo, que permitan presumir el hecho oculto.}
\cresumen{Los terceros ---subinquilinos, ocupantes y fiador--- integran la litis sin siempre estar individualizados. El fiador es parte con régimen propio: su responsabilidad está limitada por el plazo contractual y por el art.\ 1225 CCCN, y su garantía debe verificarse y reforzarse desde la celebración del contrato.}
\end{cornell}

\chapter{Competencia, fuero de atracción y domicilios}\label{cap:competencia}

Determinar ante qué tribunal se promueve el desalojo y dónde se notifica a cada parte son decisiones previas que condicionan todo el proceso. Este capítulo expone las reglas de competencia y su prórroga, el fuero de atracción de los procesos universales, los distintos tipos de domicilio y las reglas de notificación propias del desalojo.

\section{Competencia territorial}

El primer paso consiste siempre en revisar el contrato completo, si existe, porque en cuestiones patrimoniales en las que no interviene el orden público las partes pueden haber pactado una \textbf{prórroga de competencia}: un \textbf{pacto de foro prorrogado}, que fija de antemano el tribunal competente, o una \textbf{cláusula compromisoria}, que deriva el conflicto a un tribunal arbitral ---por ejemplo, el de un colegio profesional, el de la Bolsa de Comercio, o el de otro departamento judicial---.

\begin{example}{El pacto de foro prorrogado hacia otra jurisdicción}
Un contrato puede pactar que, ante cualquier conflicto, entiendan tribunales distintos de los del lugar donde está ubicado el inmueble ---por ejemplo, de otra jurisdicción---, aun cuando el inmueble se encuentre en la Ciudad de Buenos Aires.
\end{example}

A falta de pacto o de contrato disponible se aplican las reglas generales.

\begin{formula}{Artículo 5 del CPCCN ---competencia (fragmento)}
La competencia se determina por la naturaleza de las pretensiones deducidas en la demanda y no por las defensas opuestas por el demandado. Salvo los casos de prórroga expresa o tácita, cuando procediere, y sin perjuicio de las reglas especiales de este Código y de otras leyes, será juez competente, cuando se ejerciten acciones reales sobre bienes inmuebles, el del lugar donde esté situada la cosa litigiosa.
\end{formula}

En los casos trabajados, a falta de pacto de foro prorrogado o cláusula compromisoria, la competencia se determina por el lugar de ubicación del inmueble ---arts.\ 5 y siguientes del CPCCN--- o, en su defecto, por el domicilio del demandado. Otra fuente la formula como el lugar donde está el inmueble o donde se celebró el contrato, salvo prórroga de jurisdicción pactada.
% TODO: verificar consistencia entre fuentes originales. Las clases formulan la regla general de competencia de tres maneras: lugar del inmueble o, en su defecto, domicilio del demandado (31/08); lugar del inmueble (14/09); lugar del inmueble o del contrato, salvo prórroga (24/09).

\begin{definition}{La prórroga depende de la voluntad de las partes en cada instancia}
Aunque el contrato establezca un pacto de foro prorrogado o una cláusula compromisoria, es posible iniciar la demanda ante el juez nacional en lo civil. El límite está en el momento en que se corre traslado: la parte demandada puede oponer la excepción de incompetencia al contestar la demanda o aceptar la competencia del fuero elegido por el actor. Si el demandado contesta sin oponer incompetencia, se produce una \textbf{prórroga tácita} que consolida la competencia de ese tribunal. La decisión queda, en definitiva, en manos del demandado.
\end{definition}

\section{Fuero de atracción de los procesos universales}

En los procesos universales ---sucesiones, quiebras y concursos--- todos los procesos vinculados al patrimonio de la persona física o jurídica deben ser atraídos y tramitados ante el mismo juez, para concentrar el tratamiento del patrimonio en un único proceso.

\begin{important}
Cuando el locatario utiliza el inmueble exclusivamente para vivienda propia o de su familia, el contrato de locación queda ajeno al concurso o quiebra: no opera el fuero de atracción. En cambio, cuando la locación tiene destino comercial, sí resulta atraída por el fuero del concurso o la quiebra. Esta distinción no siempre es advertida de oficio por el tribunal y depende en buena medida de que la contraparte la plantee.
\end{important}

% TODO: verificar referencia presente en las fuentes originales. La clase cita la norma como <<artículo 157 de la ley de concursos y quiebras>> (Ley 24.522); la propia fuente indica cotejar la referencia con el texto vigente antes de invocarla en un escrito.

\begin{formula}{Regla práctica}
Contrato con destino comercial: opera el fuero de atracción y se acude al juez del concurso o la quiebra. Contrato con destino vivienda: queda afuera y puede desalojarse.
\end{formula}

\section{Los distintos domicilios}\label{sec:domicilios}

Es un error frecuente confundir domicilios que cumplen funciones procesales diferentes. En los escritos suele emplearse la fórmula <<constituyo domicilio legal en\ldots>> cuando en realidad se constituye un domicilio procesal o \emph{ad litem}.

\begin{table}[H]
\centering
\caption{Tipos de domicilio en el proceso}
\begin{tabularx}{\textwidth}{>{\raggedright\arraybackslash}p{0.26\textwidth} X}
\toprule
\textbf{Tipo de domicilio} & \textbf{Función y características} \\
\midrule
Especial & Pactado en el contrato de locación o comodato, para las notificaciones relativas al vínculo contractual. No equivale al domicilio procesal ni al legal. \\
Procesal o \emph{ad litem} (constituido) & Constituido a los fines del proceso; debe ubicarse dentro del radio del juzgado donde tramita la causa. \\
Legal & El que fija la ley para ciertos sujetos: por ejemplo, el de una sociedad registrada, o el de una persona con capacidad restringida, que es el de su curador. \\
Electrónico constituido & Obligatorio en todo proceso judicial. En el orden nacional, el sistema es el M.E.V.; en los departamentos judiciales de la provincia de Buenos Aires, Notificaciones CBA. \\
Real & Lugar de residencia efectiva. Se utiliza para actos específicos: la sentencia y el traslado de la demanda. \\
\bottomrule
\end{tabularx}
\end{table}

\subsection{Domicilio especial en el inmueble}

Es habitual fijar en los contratos un domicilio especial en el propio inmueble locado, distinto del domicilio real o del domicilio legal de una eventual sociedad locataria, para asegurar que todas las notificaciones ---incluidas las de una futura demanda--- lleguen a destino.

\begin{important}
El domicilio especial fijado en el inmueble no es el domicilio real: se establece, para el locatario, en el propio inmueble y, para el locador, en el domicilio donde sabe que recibirá notificaciones. Si el inquilino se traslada a otra provincia o a un municipio más lejano, las notificaciones igualmente se practican en el inmueble: si el destinatario está presente, se notifica en persona; si no, queda de todos modos notificado. Si el locatario es una sociedad, conviene el del inmueble y no la sede social, porque esta puede modificarse por estatuto en cualquier momento.
\end{important}

\subsection{Reglas especiales de notificación en el desalojo}

\begin{formula}{Notificación en el inmueble objeto del desalojo}
A diferencia de otros procesos (por ejemplo, una compraventa o un reclamo de daños), en el desalojo se considera válidamente notificada a la parte demandada en el domicilio del propio inmueble. Es una regla especial de esta clase de proceso.
\end{formula}

En el desalojo el traslado de la demanda se dirige siempre contra el sujeto pasivo que figura en el contrato ---persona física o jurídica--- y siempre se notifica además a subinquilinos y ocupantes, aun cuando estén expresamente prohibidos por el contrato (véase el capítulo~\ref{cap:terceros}).

\begin{important}
Solo puede notificarse por vía electrónica cuando la contraparte ya constituyó su domicilio electrónico en el proceso, es decir, cuando ya se presentó con su letrado. Si todavía no se presentó, corresponde cédula en papel al domicilio procesal constituido ---que en la generalidad de los casos coincide con el estudio jurídico del abogado--- o, si no hay domicilio constituido, al domicilio real, donde debe haber alguien que reciba la cédula identificándose con su documento. Muy pocos actos se notifican al domicilio real: la sentencia y el traslado de la demanda, porque en esas etapas la contraparte no tiene todavía domicilio procesal. El domicilio electrónico reemplazó en gran medida la cédula en papel que antes diligenciaba el oficial notificador.
\end{important}

En el patrocinio jurídico gratuito es frecuente constituir domicilio en los estrados del juzgado, lo que habilita que las notificaciones corran automáticamente los días de nota, con el riesgo de que el letrado pierda el plazo si no controla esos días.

Ante la dificultad de contactar a una parte se recurre a todos los canales disponibles ---llamado telefónico, mensaje de texto, correo electrónico y carta simple---, para no depender únicamente de que el correo postal entregue en mano. También se utiliza el domicilio electrónico constituido, si existe. Los datos mínimos del primer escrito (domicilio real de cada actor, domicilio electrónico del letrado constituido en la sede electrónica del Colegio de Abogados, teléfono optativo) se tratan en el capítulo~\ref{cap:demanda}.

\subsection{El domicilio electrónico pactado en el contrato}

\begin{note}
Una ley anterior ---hoy derogada por el DNU 70/2023--- establecía como obligatorio incluir en los contratos de locación un domicilio de correo electrónico, y disponía que toda notificación a esa dirección servía como prueba definitoria; con el DNU, muchos contratos dejaron de incluir la cláusula. Un proyecto de ley pretende reincorporar la obligatoriedad solo para el domicilio electrónico constituido por la parte, no el del abogado. Cuando las partes pactaban sus correos, ninguna podía alegar desconocimiento, dado que también debía comunicar cualquier cambio de dirección. Esa norma, pensada en beneficio del locatario ---para notificar al locador una filtración o un desperfecto---, terminó perjudicándolo al habilitar notificar por esa vía actos procesales relevantes que antes requerían notificación formal.
\end{note}

% TODO: verificar consistencia entre fuentes originales. La clase del 20/08 indica que el domicilio electrónico <<dejó de ser obligatorio>> con el DNU 70/2023, mientras la del 31/08 lo describe como obligatorio en todo proceso judicial. Parecen referirse, respectivamente, al correo pactado en el contrato y al domicilio electrónico procesal; se conservan ambas afirmaciones.

\subsection{Cómputo de los plazos}

\begin{formula}{Días corridos y días hábiles}
Los plazos del Código Civil y Comercial ---derecho de fondo--- se cuentan en \textbf{días corridos}. Los plazos del Código Procesal ---derecho de forma--- se cuentan en \textbf{días hábiles}.
\end{formula}

\begin{sintesis}
La competencia surge del contrato (foro prorrogado o cláusula compromisoria) y, a falta de él, del lugar del inmueble; el demandado decide si acepta o no la prórroga. El fuero de atracción alcanza a los contratos de destino comercial y no a los de vivienda. Cada domicilio cumple una función distinta y el desalojo admite la notificación en el propio inmueble.
\end{sintesis}

\begin{cornell}{Cuadro Cornell del capítulo~\ref{cap:competencia}}
\cqa{¿Qué se revisa primero para determinar la competencia?}{El contrato completo, por si hay pacto de foro prorrogado o cláusula compromisoria; a falta de ellos, rigen los arts.\ 5 y ss.\ CPCCN (lugar del inmueble).}
\cqa{¿Qué es la prórroga tácita?}{La que se produce cuando se demanda ante un tribunal distinto del pactado y el demandado contesta sin oponer incompetencia.}
\cqa{¿Cuándo opera el fuero de atracción respecto de una locación?}{Cuando el destino es comercial, el contrato es atraído por el concurso o la quiebra; si es vivienda propia o familiar, no. Conviene verificar la norma citada (art.\ 157 Ley 24.522).}
\cqa{¿Cuáles son los domicilios y para qué sirve cada uno?}{Especial (contrato), procesal o \emph{ad litem} (proceso, dentro del radio del juzgado), legal (lo fija la ley), electrónico constituido (notificaciones) y real (sentencia y traslado de demanda).}
\cqa{¿Dónde se notifica válidamente al demandado en el desalojo?}{En el domicilio del propio inmueble objeto del desalojo, regla especial de este proceso.}
\cqa{¿Cuándo se puede notificar electrónicamente?}{Solo si la contraparte ya constituyó domicilio electrónico en el proceso; antes, por cédula en papel al domicilio procesal o, en su defecto, al real.}
\cqa{¿Cómo se cuentan los plazos?}{Los de fondo (CCCN), en días corridos; los procesales, en días hábiles.}
\cresumen{La competencia y los domicilios se resuelven leyendo el contrato: foro, cláusula compromisoria y domicilio especial en el inmueble. El destino comercial activa el fuero de atracción, y cada domicilio tiene una función procesal propia que no debe confundirse.}
\end{cornell}

\part{Resolución de conflictos antes del juicio}

\chapter{Sistemas de resolución de conflictos: negociación, mediación, conciliación y arbitraje}\label{cap:masc}

Antes de analizar el proceso de desalojo conviene preguntarse si el proceso judicial es siempre la mejor respuesta al conflicto. Este capítulo ubica a la jurisdicción dentro de un mapa más amplio de sistemas de resolución de conflictos, explica la negociación por intereses y estudia la mediación, la conciliación y el arbitraje, con sus alcances y límites frente al desalojo. El capítulo~\ref{cap:vias-de-hecho} marca el límite legal de lo que el locador puede hacer por sí mismo mientras negocia.

\section{El proceso como instrumento y los métodos alternativos}

La percepción generalizada sobre el sistema judicial destaca su lentitud y el peso de la burocracia procesal. No toda cuestión tiene, sin embargo, la misma complejidad: un conflicto sencillo y uno complejo no deberían tramitar de manera idéntica, y por eso determinados procesos ---como el desalojo por falta de pago--- cuentan con un trámite más ágil que el ordinario. El sistema procesal tiende, aun así, a operar como un \textbf{molde único} en el que se encajan todos los conflictos, cuando debería existir una diversidad de moldes. Sería razonable otorgar al juez mayor flexibilidad para adecuar las formas al conflicto.

\begin{definition}{Derecho procesal instrumental y MASC}
Las formas procesales son una \textbf{garantía} de certeza y seguridad jurídica, no un fin en sí mismo; el derecho procesal es \textbf{instrumental} respecto del derecho de fondo, al que sirve. Cuando el proceso, tal como está diseñado, no da una respuesta adecuada al conflicto, aparecen los \textbf{métodos alternativos de resolución de conflictos (MASC)} como vía complementaria ---no sustitutiva--- de la jurisdicción estatal.
\end{definition}

\begin{important}
La forma es garantía, no finalidad: el derecho de fondo no debe subordinarse a la forma procesal. Si el proceso no permite realizar el derecho de fondo, no está cumpliendo su función. Aun así, la vía judicial no debe abandonarse: los MASC se utilizan como complemento y, de ser posible, en paralelo.
\end{important}

El abogado o abogada debe analizar la conflictiva a partir del concepto de conflicto de Remo Entelman: la interferencia intersubjetiva de intereses entre dos o más partes, que puede derivar del incumplimiento de un contrato de locación, y trabajarla desde una etapa temprana, para que no escale. El vínculo entre las partes puede pensarse como un canal de comunicación (emisor, receptor, código), de modo que corresponde buscar distintas maneras de llegar a la otra parte, aun cuando en un primer momento parezca imposible dialogar. Aunque los desalojos son extensos, engorrosos y lentos, siempre conviene explorar ese canal: citar a todas las partes que se pueda, de manera virtual o presencial, porque de allí pueden surgir alternativas de todo tipo, incluidas económicas.

\section{Mapa de los sistemas de resolución de conflictos}

Los mecanismos se ordenan en dos familias, según la estructura del enfrentamiento y la capacidad del tercero para imponer una solución.

\begin{definition}{Autocomposición y heterocomposición asistida}
En los sistemas \textbf{no adversariales} no hay enfrentamiento directo entre las partes. En la \textbf{autocomposición}, las propias partes resuelven el conflicto sin intervención de un tercero: el caso paradigmático es la \textbf{negociación directa} entre ellas o entre sus representantes legales. En la \textbf{heterocomposición asistida} interviene un tercero que ayuda a las partes a acercar posiciones, sin poder imponerles una solución: es el caso de la \textbf{mediación} y la \textbf{conciliación}.
\end{definition}

En los sistemas \textbf{adversariales} hay enfrentamiento directo y el tercero sí puede imponer una solución obligatoria: el \textbf{arbitraje} y el \textbf{proceso judicial}.

\begin{important}
La diferencia esencial entre los sistemas no adversariales y los adversariales no reside en la presencia de un tercero, sino en si ese tercero tiene o no \textbf{poder para imponer una solución obligatoria} a las partes.
\end{important}

\begin{table}[H]
\centering
\caption{Clasificación de los sistemas de resolución de conflictos}
\begin{tabularx}{\textwidth}{>{\raggedright\arraybackslash}p{0.22\textwidth} >{\raggedright\arraybackslash}p{0.22\textwidth} X >{\raggedright\arraybackslash}p{0.18\textwidth}}
\toprule
\textbf{Sistema} & \textbf{Modalidad} & \textbf{Mecanismo} & \textbf{¿El tercero impone?} \\
\midrule
No adversarial & Autocomposición & Negociación directa & No \\
No adversarial & Heterocomposición & Mediación & No \\
No adversarial & Heterocomposición & Conciliación & No \\
Adversarial & Heterocomposición & Arbitraje & Sí \\
Adversarial & Heterocomposición & Proceso judicial & Sí \\
\bottomrule
\end{tabularx}
\end{table}

En términos de vías, suele distinguirse la \textbf{vía composicional, cooperativa y colaborativa} ---en la que las partes, llamadas aquí <<partícipes>> y no parte actora y parte demandada, construyen juntas la solución cediendo cada una algo--- de la \textbf{vía adversarial}, que incluye el proceso judicial y el arbitraje. Esta última se contrapone al diseño confrontativo del Código Procesal, donde quien demanda afirma los hechos y el demandado niega, desconoce y rechaza cada punto. Si bien el arbitraje podría aplicarse en teoría a un desalojo, en la práctica no es habitual: la vía más frecuente es la composicional cuando hay margen de acuerdo, o la judicial cuando no lo hay.

\section{La negociación por intereses}

La negociación eficaz entre partes o entre sus abogados se rige por cuatro principios del método de negociación basado en intereses.

\begin{enumerate}
\item \textbf{Separar a las personas del problema.} Los conflictos interpersonales y los de intereses son dimensiones distintas que no deben mezclarse; las relaciones personales deben preservarse. En el contexto profesional, pelearse con un colega cada vez que se litiga duplica el problema. Se trabaja sobre los intereses \emph{con} el apoyo de las personas, no contra ellas.
\item \textbf{Negociar sobre intereses y no sobre posiciones.} La \textbf{posición} es la demanda concreta de cada parte (<<quiero esto>>, <<te ofrezco tanto>>); detrás hay un \textbf{interés real}, la razón subyacente. Negociar por posiciones suele llevar a un punto medio; trabajar sobre intereses ---a veces comunes o compatibles--- abre acuerdos que aquel no permite.
\item \textbf{Generar opciones de beneficio mutuo.} Antes de cerrar una posición conviene una \textbf{tormenta de ideas}: poner sobre la mesa varias soluciones y evaluarlas. Adoptar una posición de entrada es riesgoso porque luego es difícil retroceder sin perder credibilidad.
\item \textbf{Usar criterios objetivos.} Un valor de mercado, un antecedente, una tasación, un índice de referencia permiten evaluar las propuestas racionalmente.
\end{enumerate}

\begin{example}{Las dos empresas y la naranja}
Dos empresas disputaban la misma materia prima, cada una sosteniendo quedarse con el cien por ciento. Al indagar los intereses reales se descubrió que una necesitaba la \textbf{cáscara} (para mermelada) y la otra el \textbf{jugo}. Resuelto el conflicto sobre los intereses ---y no sobre las posiciones---, ambas obtuvieron el cien por ciento de lo que necesitaban. Cuando se negocia por posiciones el resultado suele ser un reparto a mitades; cuando se negocia por intereses, ambas partes pueden obtener mucho más.
\end{example}

\subsection{Técnicas en la práctica profesional}

\begin{itemize}
\item \textbf{Escucha activa}: distinguir entre oír y escuchar; la comunicación escrita con el cliente, preferible a los audios, permite documentar con precisión las instrucciones y evitar reproches posteriores.
\item \textbf{Análisis de los intereses en juego}: explicar siempre al cliente los costos y riesgos de un proceso, incluida la posibilidad de que las costas se inviertan y deban afrontarse.
\item \textbf{Parámetros objetivos}: por ejemplo, una pericial arbitral, distinta de la pericial dentro de un proceso judicial, para dirimir cuestiones técnicas puntuales.
\item \textbf{Ampliar las opciones} (<<agrandar la torta>>): buscar alternativas creativas antes de litigar.
\item \textbf{<<Una cosa por otra>>}: un intercambio compensatorio, como ofrecer al ocupante una suma para la mudanza, e incluso dejarle adeudada la deuda de alquileres, a cambio de que se retire voluntariamente, evitando un proceso que puede extenderse más de un año y medio.
\end{itemize}

\begin{example}{La pericial arbitral y el calefón}
Un locatario se negaba a pagar el alquiler argumentando que el calefón no funcionaba bien. Se propuso designar a un plomero como perito arbitral, quien constató que el desperfecto se debía a un uso incorrecto de la palanca reguladora entre agua fría y caliente por el propio ocupante, y no a un defecto atribuible al locador. La pericial permitió despejar objetivamente el reclamo y evitar una discusión judicial prolongada.
\end{example}

\begin{example}{Pagar el alojamiento a cambio de la desocupación}
Ante la consulta de si puede ofrecerse al ocupante de una vivienda el pago de tres meses de alojamiento en otro lugar para que desocupe y luego se reinserte en un nuevo contrato, corresponde advertir que difícilmente se le alquile un lugar si no cuenta con una garantía propietaria o un seguro de caución. Una alternativa consiste en que el propio locador le ofrezca alquilarle directamente por tres meses un lugar para que desocupe el inmueble propio. Este tipo de acuerdo debe ser voluntariamente pactado: el juez no puede obligar a la parte demandada a aceptarlo, y lograrlo es tarea de los abogados y no del tribunal.
\end{example}

\begin{example}{El caso sin margen de negociación}
En un caso llevado hasta el final ---con sentencia firme, orden de lanzamiento y presencia de policía y cerrajero--- no hubo forma de negociar durante el trámite. Pero en el momento mismo del lanzamiento, con el flete ya presente, el ocupante aceptó retirarse a cambio de que se le pagaran treinta días de estadía en una habitación alquilada, permitiéndosele además llevar sus pertenencias. El acuerdo se cerró en el instante, frente al oficial de justicia. Todas las alternativas que el cliente permita son válidas, siempre que sean legales.
\end{example}

\section{Mediación y conciliación}

Ambos son mecanismos de heterocomposición asistida por un tercero, con una diferencia funcional esencial.

\begin{definition}{Mediación y conciliación}
El \textbf{mediador} asiste a las partes para que sean ellas mismas quienes encuentren la solución; no puede proponer fórmulas conciliatorias ni resolver el caso. El \textbf{conciliador} puede proponer \textbf{fórmulas conciliatorias}, pero tampoco puede resolver el caso por sí mismo.
\end{definition}

\begin{table}[H]
\centering
\caption{Comparación entre mediador y conciliador}
\begin{tabularx}{\textwidth}{X >{\centering\arraybackslash}p{0.22\textwidth} >{\centering\arraybackslash}p{0.22\textwidth}}
\toprule
\textbf{Característica} & \textbf{Mediador} & \textbf{Conciliador} \\
\midrule
Asiste a las partes & Sí & Sí \\
Puede proponer fórmulas & No & Sí \\
Puede imponer una solución & No & No \\
Decide el conflicto & No & No \\
\bottomrule
\end{tabularx}
\end{table}

\subsection{El juez como conciliador: la audiencia del artículo 360}

En la \textbf{audiencia preliminar} del artículo 360 del CPCCN el juez actúa con facultades de conciliador y puede proponer fórmulas a las partes. Si la conciliación no prospera, fija los \textbf{hechos controvertidos} y ordena los medios de prueba que se producirán (véase el capítulo~\ref{cap:audiencia}). El propio Código resuelve la tensión con el riesgo de \textbf{prejuzgamiento}: establece que el juez no incurre en él por formular propuestas conciliatorias. En esta audiencia conviene que las partes hablen directamente, y no solo sus abogados, porque mediación y conciliación presuponen que las propias partes busquen la solución.

\begin{definition}{Principio de confidencialidad}
Todo lo manifestado por las partes durante una audiencia de mediación o conciliación es \textbf{confidencial} y no puede utilizarse como prueba en el proceso judicial posterior si no se llega a un acuerdo; el mediador tampoco puede ser citado como testigo sobre lo tratado. Esta garantía permite hablar con franqueza sin temor a perjudicarse en una etapa ulterior.
\end{definition}

\subsection{Tipos, plazos y efectos de la mediación prejudicial}

\begin{itemize}
\item \textbf{Mediación privada}: la parte requirente elige un mediador de la lista habilitada por el Ministerio de Justicia y debe proponer \textbf{tres mediadores} a la otra parte para que el requerido elija uno.
\item \textbf{Mediación oficial}: se presenta un formulario ante la mesa de entradas de la cámara correspondiente ---actualmente por vía electrónica--- y el mediador es designado \textbf{por sorteo}.
\end{itemize}

La mediación tiene un plazo de \textbf{sesenta días} para las audiencias, prorrogable. Si no hay acuerdo se labra un \textbf{acta de cierre}, que debe acompañarse como \textbf{requisito de admisibilidad} de la demanda: si el proceso se inicia sin mediación previa, el juez remite el expediente a mediación antes de continuar. La mediación prejudicial obligatoria se instituyó porque el proceso judicial estaba congestionado; dentro del proceso ya iniciado, la audiencia del artículo 360 es una segunda oportunidad de conciliación.

% TODO: verificar consistencia entre fuentes originales. La clase del 24/08 presenta la mediación prejudicial como requisito de admisibilidad con acta de cierre; las clases del 20/08 y del 10/09 indican que en los juicios de desalojo no es obligatoria. Se conservan ambas posiciones.
\begin{note}
Las fuentes distinguen el régimen general de la mediación prejudicial ---con acta de cierre como requisito de admisibilidad--- de lo que ocurre en el desalojo, donde se la describe como optativa y librada al criterio del profesional (véase el capítulo~\ref{cap:intimacion}). Debe verificarse la norma aplicable antes de acompañar o prescindir del acta de cierre.
\end{note}

\begin{note}
Una diferencia clave entre la mediación y la audiencia preliminar es el \textbf{principio de informalidad}: en la mediación no hay obligación de exhibir la prueba (las partes pueden mencionar que cuentan con determinados elementos sin mostrarlos), mientras que en la audiencia del artículo 360 ya se formuló la demanda, se fundó la pretensión y se ofreció la prueba: toda la evidencia está <<sobre la mesa>>, y el juez, que conoce la demanda, propone fórmulas con mayor conocimiento del caso.
\end{note}

\begin{normacitada}
\textbf{Sanción por incomparecencia a la mediación.} El régimen original de la Ley de Mediación Prejudicial Obligatoria preveía la imposición \emph{obligatoria} de una multa a la parte requerida que no comparecía sin causa justificada. La reforma de la \textbf{Ley 26.589} sustituyó ese mandato por una \emph{facultad} del juez (<<podrá>> imponer la multa), lo que diluyó su aplicación: casi no se aplican multas, aunque una multa disuasiva sería muy efectiva frente a quienes firman el acuerdo de comparecencia y luego no concurren.
\end{normacitada}

\subsection{Mediación, prescripción y caducidad}

\begin{definition}{Efecto de la mediación sobre la prescripción}
La \textbf{prescripción} extingue la acción ante la inacción del titular durante el plazo legal, y puede interrumpirse o suspenderse. Al iniciarse la mediación se \textbf{suspende el curso de la prescripción}, suspensión que se extiende hasta \textbf{veinte días después} de labrada el acta de cierre; desde entonces corre el plazo remanente.
\end{definition}

\begin{example}{Cálculo del plazo remanente}
A una acción le quedaban solo siete días de plazo cuando se inició la mediación, que dura noventa. Labrada el acta de cierre, transcurren veinte días más y recién entonces el plazo vuelve a correr: quedan \textbf{siete días}. Si no se demanda dentro de ellos, la acción prescribe y la demanda queda expuesta a la excepción de prescripción. La misma lógica de plazos se aplica a las medidas cautelares (véase el capítulo~\ref{cap:cautelares}).
\end{example}

\begin{important}
Cuando el cliente afirma que <<el deudor es de absoluta confianza y jamás dejará prescribir la acción>>, el abogado debe advertirle expresamente ---y por escrito, mediante consentimiento informado--- que el plazo está corriendo y que dejarlo vencer implica perder la acción. La confianza del cliente en su deudor no es oponible frente a la excepción de prescripción.
\end{important}

\section{Arbitraje}

En el arbitraje, a diferencia del proceso judicial, las propias partes se \textbf{sustraen voluntariamente de la jurisdicción estatal} y someten el conflicto a uno o más árbitros.

\begin{definition}{Cláusula compromisoria y compromiso arbitral}
La \textbf{cláusula compromisoria} es el pacto incluido en el contrato \emph{antes} de que exista el conflicto, por el cual las controversias futuras serán resueltas por arbitraje; puede designar al árbitro o remitir al procedimiento estandarizado de una institución (por ejemplo, una cámara de comercio). El \textbf{compromiso arbitral} se suscribe \emph{una vez que el conflicto ya existe} y fija las cuestiones concretas que el árbitro resolverá; los árbitros no pueden fallar sobre cuestiones no sometidas ni fuera del plazo establecido.
\end{definition}

La cláusula compromisoria cumple dos funciones. La \textbf{función positiva}: si surge el conflicto y la otra parte no quiere firmar el compromiso, puede demandarse que se constituya; si no comparece, el juez puede suplir esa voluntad y designar al árbitro. La \textbf{función negativa}: si pese a la cláusula se demanda ante la justicia ordinaria, el demandado puede oponer la excepción de incompetencia, ya que el competente es el tribunal arbitral.

\subsection{Jurisdicción, competencia y prórroga}

\begin{definition}{Jurisdicción y competencia}
La \textbf{jurisdicción} es la potestad de juzgar, atribuida a todos los jueces. La \textbf{competencia} es la medida de esa jurisdicción para un caso concreto: se divide por materia, territorio y personas.
\end{definition}

Solo puede prorrogarse la competencia \textbf{territorial} y en asuntos \textbf{patrimoniales}; no la que se determina por razón de la materia ni de las personas (véase el capítulo~\ref{cap:competencia}).

\begin{example}{Prórroga de competencia en una locación}
La locación es una cuestión patrimonial, por lo que es válido prorrogar la competencia territorial. Supóngase un inmueble ubicado en Catamarca, con contrato firmado en Córdoba y una cláusula que fija la competencia de los tribunales de la Ciudad de Buenos Aires: si la demanda de desalojo se interpone en Córdoba, la otra parte puede oponer excepción de incompetencia y el planteo prospera.
\end{example}

\begin{note}
La competencia territorial y patrimonial es \textbf{disponible} para las partes, por lo que puede prorrogarse y, en principio, el juez no puede declararla de oficio. Distinto es el caso de la competencia federal frente a la local: si se demanda en el fuero equivocado, el juez puede declararse incompetente de oficio, porque no es materia disponible. Del mismo modo, si hay cláusula compromisoria y se demanda judicialmente, la otra parte puede oponer la excepción correspondiente, que prosperará.
\end{note}

\subsection{Límites del árbitro frente al juez}

Se discute si la función del árbitro es jurisdiccional en sentido estricto. Si la jurisdicción consiste en fijar los hechos controvertidos, determinar si hay una norma que los contemple y dictar una decisión concreta, el árbitro cumple esa función: fija los hechos, los subsume en una norma y dicta un \textbf{laudo}, con la misma función que la sentencia.

\begin{definition}{Imperium}
Es el \textbf{poder coactivo} del Estado que permite al juez ejecutar sus decisiones de manera forzada, incluido el empleo de la fuerza pública.
\end{definition}

\begin{important}
El árbitro \textbf{carece de imperium}: no puede ejecutar coactivamente su laudo ni trabar medidas cautelares por sí mismo. Si la parte condenada no cumple voluntariamente, la parte interesada debe acudir al \textbf{juez estatal que hubiera sido competente} de no existir la cláusula arbitral. Cuando lo prioritario es recuperar un inmueble cuanto antes, como en el desalojo, el arbitraje puede suponer un paso adicional ---laudo y luego juez para ordenar el lanzamiento--- que a veces demanda más tiempo que ir directamente a la vía judicial.
\end{important}

Otra diferencia: mientras las decisiones judiciales deben fundarse en derecho, un \textbf{arbitraje de amigables componedores} puede resolver según la \textbf{equidad}, <<a verdad sabida y buena fe guardada>>, sin la misma exigencia de fundamentación normativa exhaustiva.

\begin{sintesis}
Los sistemas se distinguen por la facultad del tercero de imponer la solución: negociación, mediación y conciliación no la tienen; arbitraje y juicio, sí. Negociar por intereses abre soluciones que las posiciones cierran. La mediación suspende la prescripción y exige cuidar los plazos. El árbitro carece de imperium, lo que limita su utilidad en el desalojo.
\end{sintesis}

\begin{cornell}{Cuadro Cornell del capítulo~\ref{cap:masc}}
\cqa{¿Qué distingue a los sistemas adversariales de los no adversariales?}{No es la presencia de un tercero sino su poder de imponer una solución obligatoria: lo tienen el árbitro y el juez; no el mediador, el conciliador ni las partes que negocian.}
\cqa{¿Qué diferencia hay entre mediador y conciliador?}{El mediador solo asiste para que las partes encuentren la solución; el conciliador puede además proponer fórmulas. Ninguno decide el conflicto.}
\cqa{¿Cuáles son los cuatro principios de la negociación por intereses?}{Separar personas del problema; negociar sobre intereses y no posiciones; generar opciones de beneficio mutuo; usar criterios objetivos.}
\cqa{¿Qué efecto tiene la mediación sobre la prescripción?}{Suspende su curso hasta veinte días después del acta de cierre; luego corre el plazo remanente.}
\cqa{¿Qué diferencia hay entre cláusula compromisoria y compromiso arbitral?}{La primera se pacta antes del conflicto; el segundo, una vez surgido, y fija las cuestiones concretas a resolver.}
\cqa{¿Por qué el arbitraje puede ser lento en el desalojo?}{Porque el árbitro carece de imperium: el laudo debe ejecutarse ante el juez estatal competente, lo que añade un paso.}
\cqa{¿Qué protege el principio de confidencialidad?}{Que lo dicho en mediación o conciliación no pueda usarse como prueba en el juicio posterior si no hay acuerdo; el mediador no puede declarar como testigo.}
\cresumen{El desalojo puede abordarse por vías complementarias del juicio. Mediación y conciliación asisten a las partes sin imponer; el arbitraje impone pero sin imperium. Negociar por intereses y cuidar la prescripción y la confidencialidad son claves prácticas, y debe verificarse el carácter obligatorio u optativo de la mediación.}
\end{cornell}

\chapter{Vías de hecho prohibidas y desalojos administrativos}\label{cap:vias-de-hecho}

La negociación tiene un límite legal: el locador que no logra recuperar el inmueble no puede hacerlo por sus propios medios. Este capítulo estudia la prohibición de las vías de hecho, el margen de acción frente a una ocupación reciente, el caso del inmueble desalojado por una autoridad y las condiciones en que la administración puede, excepcionalmente, desalojar sin intervención judicial previa.

\section{La prohibición de las vías de hecho: el caso del candado}

Un caso frecuente en la consulta profesional: un locador que, ante la falta de pago del último alquiler, ingresa al local comercial y coloca un candado al locatario. La primera reacción no debe ser demandar directamente, sino intentar \textbf{negociar}; si la otra parte tiene representación legal, la negociación se conduce con su abogado y no directamente con el cliente. Pero, más allá de cómo termine la negociación, hay un principio indiscutible.

\begin{important}
Por más que exista incumplimiento del locatario ---incluso falta de pago---, el locador \textbf{no puede recuperar la tenencia por la fuerza}. Las vías de hecho contra quien tiene la tenencia están expresamente prohibidas: debe acudir a la \textbf{acción de desalojo} por la vía judicial correspondiente.
\end{important}

\begin{normacitada}
\textbf{Prohibición de vías de hecho.} El CCCN prohíbe al propietario o locador recuperar la tenencia de un inmueble por propia autoridad, aun cuando exista incumplimiento del locatario; solo puede recuperarse por la vía judicial (acción de desalojo). Actuar por la fuerza expone al locador a una \textbf{denuncia penal} por usurpación y a la obligación de restituir la tenencia.
\end{normacitada}

\begin{note}
El cambio de cerradura o el candado puede ser fácilmente objeto de denuncia penal y resulta sencillo de comprobar. La situación difiere en una \textbf{ocupación ilegal reciente}: en el momento mismo del hecho el titular puede actuar en defensa de la posesión, pero cuanto más tiempo transcurre y más se consolida la ocupación, menor es el margen para actuar por medios propios y mayor la exposición a responsabilidad por la propia agresión.
\end{note}

\subsection{Inmueble desalojado por una autoridad}

Si el inmueble fue desalojado por una autoridad (por ejemplo, clausura administrativa por riesgo de derrumbe) y ya no hay resistencia del ocupante, quien era poseedor o dueño puede en principio recuperarlo, salvo que haya intervenido como parte en el procedimiento que originó la clausura: en ese caso debe pedir autorización judicial. Si no fue parte, puede recuperarlo libremente, siempre que no exista otro impedimento (proceso penal en curso, problema de violencia, investigación activa que mantenga afectado el inmueble).

\section{Desalojos administrativos y sus límites}

Existen desalojos realizados por el Gobierno de la Ciudad de Buenos Aires al amparo de un decreto administrativo. Ese decreto existe y es real, pero su aplicación está limitada a causales específicas y extremas: situaciones de emergencia habitacional grave, como el riesgo de derrumbe estructural.

\begin{example}{Edificio abandonado con riesgo de derrumbe}
Un edificio abandonado por la constructora que lo levantó (que quebró y lo dejó inconcluso) fue tomado por numerosas familias, sin servicios mínimos y con riesgo edilicio en losas y columnas. En esa situación extrema el decreto habilita el llamado <<desalojo administrativo>>, sin intervención judicial previa.
\end{example}

\begin{important}
El decreto solo puede fundarse en causales extremas de emergencia (riesgo de derrumbe, ausencia de condiciones mínimas de habitabilidad) y no puede utilizarse frente a la mera ocupación irregular de un inmueble, situación que requiere la intervención del Poder Judicial y de las defensorías.
\end{important}

\begin{example}{Medida cautelar judicial frente a un desalojo administrativo}
Un juez de la Ciudad de Buenos Aires consideró que la Ciudad invocaba el decreto de emergencia fuera de sus causales y dictó una medida cautelar, dentro de un proceso colectivo, para frenar ese desalojo administrativo hasta que la situación fuera revisada judicialmente.
\end{example}

Frente a un edificio verdaderamente en riesgo, el Poder Ejecutivo debe actuar para resolver un peligro real; pero el mismo mecanismo no puede emplearse como atajo en conflictos de ocupación sin ese riesgo extremo, porque para ellos están el Poder Judicial y las defensorías, que analizan la situación de cada persona involucrada, sea o no ocupante irregular.

\begin{sintesis}
El locador no puede recuperar la tenencia por la fuerza aunque el locatario incumpla: debe acudir al desalojo. La ocupación reciente admite defensa inmediata, pero ese margen se reduce con el tiempo. El desalojo administrativo es excepcional, limitado a emergencias graves y sujeto a control judicial.
\end{sintesis}

\begin{cornell}{Cuadro Cornell del capítulo~\ref{cap:vias-de-hecho}}
\cqa{¿Puede el locador recuperar el inmueble por la fuerza si hay falta de pago?}{No. Las vías de hecho están prohibidas; debe acudir al desalojo judicial. Actuar por la fuerza expone a denuncia penal por usurpación y a restituir la tenencia.}
\cqa{¿Qué primer paso corresponde ante el locatario que no paga?}{Intentar negociar, con su abogado si tiene representación, antes de demandar.}
\cqa{¿Cambia algo si la ocupación es reciente?}{Sí: en el momento del hecho el titular puede defender la posesión, pero ese margen se reduce a medida que la ocupación se consolida.}
\cqa{¿Cuándo puede recuperar el inmueble quien fue desalojado por una autoridad?}{Si no fue parte en el procedimiento de clausura, libremente (sin otros impedimentos); si fue parte, necesita autorización judicial.}
\cqa{¿Cuándo es válido el desalojo administrativo?}{Solo ante causales extremas de emergencia, como riesgo de derrumbe; la mera ocupación irregular requiere intervención judicial y de las defensorías.}
\cresumen{La autotutela está vedada al locador: aun con incumplimiento del locatario, el camino es negociar y, si no hay acuerdo, el desalojo judicial. La administración solo puede desalojar sin juez ante emergencias graves y bajo control judicial.}
\end{cornell}

\chapter{Diligencias preliminares e intimación previa}\label{cap:intimacion}

Antes de demandar, el abogado debe saber contra quién dirigirá la acción y haber constituido en mora al deudor cuando la ley lo exige. Este capítulo estudia las diligencias preliminares que permiten conocer el objeto del proceso, su diferencia con la prueba anticipada, la intimación previa ---obligatoria en la falta de pago, recomendable en los demás casos--- y la redacción de la carta documento, incluidos sus domicilios y su contenido mínimo.

\section{Diligencias preliminares}

Las expresiones <<diligencias preliminares>> y <<medidas preparatorias>> no son sinónimos y conviene no confundirlas. Las diligencias preliminares se regulan en el Código Procesal Civil y Comercial de la Nación a partir del artículo 323 y siguientes, dentro del capítulo dedicado a ellas, en el título de los procesos de conocimiento. En la práctica del desalojo se usan poco, pero serían útiles para \textbf{focalizar el objeto} del proceso: un mismo inmueble puede ser objeto de distintos procesos simultáneos o sucesivos ---un desalojo, una acción reivindicatoria, una diligencia preliminar---, y conocer con precisión a quién se va a demandar evita que el demandado oponga una defensa de falta de legitimación por haber sido demandado incorrectamente.

\begin{formula}{Artículo 323, inciso 1.º, del CPCCN}
El proceso de conocimiento podrá prepararse pidiendo el que pretenda demandar, o quien con fundamento prevea que será demandado, que la persona contra quien se proponga dirigir la demanda preste declaración jurada, por escrito y dentro del plazo que fije el juez, sobre algún hecho relativo a su personalidad, sin cuya comprobación no pueda entrarse en juicio. El mismo artículo enumera, en sus incisos siguientes, otras diligencias posibles (exhibición de la cosa mueble, de testamento, de títulos en caso de evicción, de documentación societaria, entre otras), no todas aplicables al desalojo.
\end{formula}

El mecanismo consiste en una declaración jurada por escrito que el propio abogado redacta de antemano, en formato de preguntas y respuestas, sobre lo que resulte útil para el desalojo (por ejemplo, en qué carácter ocupa el inmueble la persona notificada). Se solicita al juzgado que ordene al oficial de justicia constituirse en el inmueble para llevar adelante el interrogatorio. Aunque el ocupante no firme ---porque no atiende o no abre la puerta---, el oficial deja constancia al devolver el mandamiento y el juzgado puede tener por cierto el contenido, habilitando el inicio del desalojo por la causal correspondiente (falta de pago, vencimiento de contrato, comodato vencido, etc.). El desalojo debe dirigirse siempre contra el sujeto identificado como ocupante, subinquilino u ocupante en general, porque de ello depende la legitimación pasiva, y debe notificarse fehacientemente ---por cédula o acta notarial---, para que nadie pueda luego alegar la nulidad por falta de notificación.

\begin{example}{La sobrina heredera única}
Una sobrina, heredera única de un tío fallecido, soltero y sin hijos, que llevaba una vida ordenada y reservada, no encontró documentación sobre la situación de varios inmuebles que el tío aparentemente poseía: no se sabía si estaban alquilados, prestados ni en qué carácter estaban ocupados. No pudo aportar datos elementales, como el teléfono del ocupante, una caja de seguridad o algún amigo cercano. El perfil del causante permitía descartar razonablemente que hubiera prestado el inmueble sin contrato, pero no saber si existía un boleto de compraventa, una seña, un comodato, una locación o incluso una donación. Antes de la vía judicial se intenta un contacto directo y no conflictivo con el ocupante mediante <<el papelito>>: una citación no formal, entregada personalmente, con múltiples canales de contacto (teléfono, correo, videollamada), que evita una carta documento que muchas veces ni siquiera se retira del correo. Si el ocupante no responde, se recurre a la diligencia preliminar. Este protocolo previo no es obligatorio sino una práctica de estilo personal del profesional.
\end{example}

\begin{definition}{Informe de dominio e índice de titularidad}
El \textbf{informe de dominio} se solicita cuando ya se cuenta con el número de matrícula del inmueble. El \textbf{índice de titularidad} sirve para determinar qué inmuebles figuran a nombre de una persona, sin conocer la matrícula; se solicita en el registro de cada jurisdicción, con su propio costo, y solo pueden pedirlo profesionales habilitados: abogados, escribanos o martilleros.
\end{definition}

Si el heredero sospecha que el causante tenía bienes en otra provincia, no existe una vía directa para relevar todo el país: solo cabría solicitar al juez que oficie al registro de esa provincia, siempre que exista algún indicio razonable, ya que el juez no libraría oficios genéricos a las veinticuatro provincias.

\subsection{Diligencia preliminar y prueba anticipada}

\begin{definition}{Diligencia preliminar y prueba anticipada: dos institutos distintos}
La \textbf{diligencia preliminar} es restrictiva y residual: se limita a los supuestos del artículo 323 y siguientes (taxativos o, según cierta doctrina, enunciativos) y queda sujeta a un \textbf{plazo de caducidad de treinta días} para iniciar el proceso principal. La \textbf{prueba anticipada}, en cambio, es un sistema asegurativo probatorio: una herramienta para no perder una prueba que podría desaparecer antes de la etapa procesal correspondiente ---el ejemplo típico es el de un testigo próximo a fallecer---. Es un instituto autónomo, no subordinado a la diligencia preliminar, y puede describirse como <<una mega cautelar>>.
\end{definition}

Una vez sustanciada la diligencia, el interesado cuenta con treinta días para iniciar el proceso principal; si no lo hace, la diligencia pierde sentido y sus efectos caducan. No siempre es necesario recurrir a ellas, pero cuando faltan datos imprescindibles deben usarse y no dejarse <<enterradas>>.

\begin{note}
En el derecho comparado, el \emph{discovery} (Estados Unidos) y el \emph{disclosure} (Reino Unido) cumplen un trabajo preparatorio equivalente, con una etapa extrajudicial de mayor amplitud, mientras que en la Argentina las diligencias preliminares son netamente judiciales. La fuente cita un artículo del Dr.\ Salgado, publicado en la revista \emph{Derecho Procesal}, que desarrolla la comparación, y advierte contra trasplantar mecánicamente institutos extranjeros: aplicarlos con mentalidad local no funcionaría.
\end{note}

\section{La intimación previa: cuándo es obligatoria}

\begin{table}[H]
\centering
\caption{Obligatoriedad de la intimación previa según la causal}
\begin{tabularx}{\textwidth}{>{\raggedright\arraybackslash}p{0.30\textwidth} >{\raggedright\arraybackslash}p{0.26\textwidth} X}
\toprule
\textbf{Causal de desalojo} & \textbf{¿Obligatoria?} & \textbf{Observación} \\
\midrule
Falta de pago (destino habitacional) & Sí & Es requisito previo de la demanda. \\
Cambio de destino & No & Se recomienda enviarla igual. \\
Vencimiento de contrato & No & Es de los objetos más típicos; se recomienda igualmente. \\
\bottomrule
\end{tabularx}
\end{table}

\begin{normacitada}
\textbf{Artículo 1222 del CCCN ---intimación de pago.} Si el destino es habitacional, previamente a la demanda de desalojo por falta de pago de alquileres, el locador debe intimar fehacientemente al locatario el pago de la cantidad debida, otorgando un plazo que nunca debe ser inferior a diez días corridos contados desde la recepción de la intimación, consignando el lugar de pago. La notificación remitida al domicilio denunciado en el contrato por el locatario se tiene por válida, aun si este se negara a recibirla o no pudiese perfeccionarse por motivos imputables a él. La fuente advierte que existe un proyecto legislativo orientado a acotar o eliminar esta exigencia.
\end{normacitada}

La obligación de intimar rige específicamente para la causal de falta de pago y no está prevista de igual manera para las otras causales. Aun así, se recomienda enviar la carta documento: transmite una seriedad que una llamada, un correo o un mensaje no logran, deja constancia documental de la buena fe de quien la envía y permite intentar una solución más rápida y económica. Intimar no siempre es obligatorio, pero casi siempre conviene.

\subsection{Falta de pago: períodos adeudados y orden público}

\begin{definition}{Orden público en la falta de pago}
El CCCN exige, como mínimo, \textbf{dos períodos} de canon locativo adeudados para habilitar el desalojo por falta de pago. Es una norma de orden público: el contrato puede pactar un plazo mayor de tolerancia (por ejemplo, tres meses), pero nunca uno menor a los dos períodos. Esta defensa debe ser invocada por la parte interesada: si el locatario no plantea que el contrato viola el orden público al exigir un solo mes de mora, el juez no lo declarará de oficio en todos los casos.
\end{definition}

\begin{formula}{Plazos de la intimación por falta de pago}
Cumplida la intimación fehaciente, deben transcurrir los plazos legales antes de demandar: dos meses de mora y diez días corridos desde el vencimiento de la intimación, conforme a la fuente, que remite a los artículos 1219 a 1227 del CCCN como marco de fondo aplicable.
\end{formula}

\begin{example}{Imputación de pagos entre expensas y canon}
Si el locatario adeuda expensas y paga puntualmente el canon de agosto, se recomienda al cliente locador imputar ese pago primero a la deuda de expensas, dejando constancia expresa del concepto imputado. Cuando llegue septiembre sin pago del alquiler, el locador queda habilitado a demandar por los dos períodos adeudados (agosto y septiembre), vencido el plazo contractual de pago (por ejemplo, del día 1 al día 10 de cada mes).
\end{example}

La obligación de pagar las expensas es una obligación \emph{propter rem}, que sigue a la cosa: si el consorcio las reclama, dirige su ejecución contra el titular de dominio y no contra el ocupante, aunque en la relación interna el locatario se haya comprometido contractualmente a pagarlas. Tras la derogación de la Ley de Alquileres por el DNU 70/2023, el precio del alquiler queda sujeto, en principio, al valor de mercado y a la libre negociación, pudiendo incluir o no las expensas.

\subsection{Alternativa menos drástica: la intimación de cese}

Cuando el cliente no busca recuperar el inmueble sino hacer cesar la actividad indebida y mantener la locación, existe una alternativa menos drástica: intimar por carta documento el cese de la actividad, invocando la cláusula contractual de destino, bajo apercibimiento de iniciar el desalojo por cambio de destino. Antes de demandar conviene evaluar siempre si el cliente quiere recuperar el inmueble (desalojo) o solo detener el uso indebido (intimación de cese). Las tres causales ---falta de pago, vencimiento y cambio de destino--- comparten la misma vía procesal pero requieren pruebas y planteos distintos.

\begin{important}
La mediación previa no es obligatoria en los juicios de desalojo: a diferencia de otros procesos, queda a criterio de cada profesional proponerla como vía alternativa antes de litigar (véase el capítulo~\ref{cap:masc}).
\end{important}

\section{La carta documento}

La carta documento, junto con el contrato de locación, constituye la prueba documental insustituible del expediente y ambos deben incorporarse siempre. Puede quedar incontestada ---nadie está obligado a responderla---: su función es intimar a pagar o a acercarse a dialogar con los representantes del actor. Si el locatario no contesta ni se comunica, queda de todos modos notificado y, vencido el plazo de espera, puede iniciarse la demanda. Si se discute su autenticidad, se ofrece prueba informativa al correo, que informe sobre la autenticidad de la pieza y la fecha de recepción.

\subsection{Formato papel y digital}

\begin{table}[H]
\centering
\caption{Carta documento en papel y carta documento digital}
\begin{tabularx}{\textwidth}{>{\raggedright\arraybackslash}p{0.18\textwidth} >{\raggedright\arraybackslash}X >{\raggedright\arraybackslash}X}
\toprule
\textbf{Aspecto} & \textbf{Papel} & \textbf{Digital} \\
\midrule
Elaboración & Tres ejemplares: uno queda en el correo, otro sellado para quien la despacha y el tercero para el destinatario. & Se hace en el sistema del Correo Argentino, sin papel. \\
Trámite & Alguien debe concurrir al correo y firmar. & No hace falta ir personalmente. \\
A quién conviene & Personas mayores, de bajos recursos o que no usan celular. & Quien maneja cómodamente medios digitales. \\
Costo & Arancelado: la carta y el envío. & También arancelado. \\
\bottomrule
\end{tabularx}
\end{table}

La gratuidad corresponde al ámbito laboral; en el ámbito civil la carta documento es arancelada tanto en su confección como en su envío. Para encuadrar correctamente la intimación antes de redactarla se necesita tener a mano el Código Procesal y el Código Civil y Comercial.

\subsection{Remitente, destinatario y domicilios}

Antes de redactar hay que identificar correctamente a los sujetos del contrato leyéndolo íntegramente. Cuando hay varios remitentes ---por ejemplo, tres locadores---, cada uno debe despachar su propia carta; de igual modo, cuando hay varios destinatarios, es una carta por persona, aunque vivan todos juntos.

\begin{table}[H]
\centering
\caption{Domicilios a consignar en la carta documento}
\begin{tabularx}{\textwidth}{>{\raggedright\arraybackslash}p{0.22\textwidth} >{\raggedright\arraybackslash}X >{\raggedright\arraybackslash}X}
\toprule
\textbf{Parte de la carta} & \textbf{Domicilio} & \textbf{Razón} \\
\midrule
Remitente (locador) & El del contrato, o el del estudio jurídico si el cliente lo pide. & Se pregunta siempre al cliente: muchos no quieren que figure su domicilio. \\
Constitución de domicilio & El del estudio (<<constituyo domicilio a los efectos del presente en\ldots>>). & Si se repite el del contrato, no se aclara nada. \\
Destinatario (locatario) & El inmueble locado. & Es un desalojo: se intima donde está el inmueble. \\
\bottomrule
\end{tabularx}
\end{table}

Antes de usar el domicilio del remitente hay que confirmar con el cliente que continúa viviendo allí.

\subsection{Contenido de la intimación por cambio de destino}

El objeto del reclamo no es que el locatario cambie el destino ni que rescinda el contrato, sino que \textbf{cese} el cambio de destino: si el contrato establece <<destino familiar únicamente>> y prohíbe otro uso, la carta cita la cláusula que fija el destino pactado y la que prohíbe modificarlo, otorga un plazo para revertir la situación y apercibe con iniciar la acción de desalojo por cambio de destino e incumplimiento contractual, con costas. Los \textbf{costos} son lo que genera de gasto; las \textbf{costas} comprenden honorarios y todo lo que lleva el proceso (tasa judicial, bono, etcétera).

Además de lo intimatorio, la carta debe abrir una vía de negociación sin usar la fórmula <<métodos alternativos>>: una solicitud de contacto con el estudio, con los medios disponibles ---WhatsApp, correo electrónico---, útil sobre todo cuando el desacuerdo dificulta el contacto directo. No existe una forma estricta de redacción: cada quien usa su tono, sin faltar el respeto.

\begin{important}
\textbf{Contenido mínimo de la carta documento de intimación:}
\begin{enumerate}
\item Fecha y lugar.
\item Remitente y destinatario (una carta por persona), con sus domicilios.
\item Cómo se obligó el destinatario (contrato con fecha y cláusula de destino).
\item Cuál es el incumplimiento (el uso que se da al inmueble).
\item Intimación a que cese, en un plazo, bajo apercibimiento de iniciar la acción de desalojo por cambio de destino, con costas.
\item Constitución de domicilio en el estudio.
\item Datos de contacto para negociar.
\item Cierre: <<Queda usted debidamente notificado>>.
\end{enumerate}
\end{important}

\begin{example}{La posición del abogado del demandado}
Quien es demandado suele esperar a que el otro active el proceso; resulta poco frecuente que su abogado se apresure a comunicarse con la contraparte. En una simulación de práctica, el abogado del locatario respondió con una construcción defensiva: ni confirmó ni negó el incumplimiento, remitió la respuesta a una consulta previa con su cliente y manifestó apertura a buscar vías distintas antes de la instancia judicial (la simulación completa se presenta en el capítulo~\ref{cap:casos}).
\end{example}

\begin{sintesis}
Las diligencias preliminares permiten conocer a quién se demandará y caducan a los treinta días. La intimación es obligatoria en la falta de pago ---mínimo dos períodos y diez días corridos--- y recomendable en las demás causales. La carta documento se despacha por persona, con domicilios correctamente consignados, y debe contener los elementos mínimos y una puerta a la negociación.
\end{sintesis}

\begin{cornell}{Cuadro Cornell del capítulo~\ref{cap:intimacion}}
\cqa{¿Para qué sirve una diligencia preliminar en el desalojo?}{Para focalizar el objeto y conocer a quién demandar (art.\ 323 inc.\ 1.º CPCCN), evitando una excepción de falta de legitimación. Caduca si el proceso principal no se inicia en treinta días.}
\cqa{¿En qué se diferencia de la prueba anticipada?}{La diligencia es restrictiva y residual; la prueba anticipada es un sistema asegurativo probatorio autónomo para no perder una prueba que podría desaparecer.}
\cqa{¿Cuándo es obligatoria la intimación previa?}{Solo en la falta de pago con destino habitacional (art.\ 1222 CCCN): plazo no inferior a diez días corridos desde la recepción, con lugar de pago. Para otras causales es una buena práctica.}
\cqa{¿Cuántos períodos adeudados exige la falta de pago?}{Como mínimo dos, por norma de orden público; el contrato puede ampliar la tolerancia pero no reducirla. La defensa debe alegarla la parte interesada.}
\cqa{¿Por qué conviene imputar el pago a expensas?}{Para que, ante el impago del mes siguiente, queden adeudados dos períodos de canon y el locador pueda demandar.}
\cqa{¿Cuántas cartas se envían si hay varios remitentes o destinatarios?}{Una carta documento por persona.}
\cqa{¿Qué domicilio se consigna para el destinatario y por qué?}{El del inmueble locado, porque se trata de un desalojo.}
\cresumen{Antes de demandar conviene identificar al ocupante, intimar fehacientemente y abrir una vía de negociación. La falta de pago exige intimación y dos períodos de mora; en los demás casos la carta documento es recomendable pero no obligatoria.}
\end{cornell}

\part{El proceso de desalojo}

\chapter{Vías procesales y causales de desalojo}\label{cap:vias}

El proceso de desalojo no se agota en una única forma de tramitar ni en una única causa. Este capítulo ubica el desalojo dentro de la clasificación de los procesos, explica qué vía procesal se le asigna en el orden nacional y en la provincia de Buenos Aires, presenta las etapas del proceso civil y estudia las causales de desalojo, con sus requisitos previos y las figuras de desalojo inmediato, reconocimiento judicial y condena de futuro.

\section{Clasificación procesal del desalojo}

El desalojo se ubica dentro de la clasificación general de los procesos del orden nacional: procesos de conocimiento, procesos universales y procesos especiales. El desalojo es un \textbf{proceso especial}. El Código Procesal le da la posibilidad de tramitar por la vía del proceso de conocimiento \textbf{ordinario}, con debate y prueba amplios, o por la vía del \textbf{sumarísimo}, más acotada. En la práctica, la mayoría de los juzgados asignan el sumarísimo, salvo que existan controversias que justifiquen el ordinario, con un plazo más extenso para ofrecer prueba.

\begin{table}[H]
\centering
\caption{Procesos de conocimiento aplicables al desalojo}
\begin{tabularx}{\textwidth}{>{\raggedright\arraybackslash}X >{\raggedright\arraybackslash}X >{\raggedright\arraybackslash}X}
\toprule
\textbf{Jurisdicción} & \textbf{Procesos de conocimiento disponibles} & \textbf{Vía habitual en el desalojo} \\
\midrule
Orden nacional (Capital Federal) & Ordinario y sumarísimo & Sumarísimo (mayoría de los casos) \\
Provincia de Buenos Aires & Ordinario, sumario y sumarísimo & Depende del criterio judicial en la primera resolución \\
\bottomrule
\end{tabularx}
\end{table}

\begin{important}
El texto histórico del CPCCN establece que el desalojo se rige por las reglas del proceso sumario, pero esa remisión está derogada: la Ley 25.488 (año 2001) eliminó el proceso sumario como categoría autónoma del orden procesal nacional. El operador debe actualizar permanentemente su lectura del Código con cada reforma.
\end{important}

En la provincia de Buenos Aires, la decisión sobre la vía se adopta en la primera resolución del expediente, y conviene identificarla de inmediato ---por ejemplo, con un cartel en la carpeta---, porque los plazos del sumarísimo no tienen nada que ver con los del ordinario. Cuando la situación es simple y poco conflictiva, el juez suele encauzar por la vía sumarísima, de plazos más breves y menos planteos posibles; cuando la demanda presenta una explicación más compleja, tiende al ordinario.

Rige para ambas vías una regla general: con los escritos postulatorios ---demanda y contestación--- deben afirmarse todos los hechos y ofrecerse toda la prueba, incluido, si se quiere, el consultor técnico.

\subsection{El desalojo en otras provincias}

Cada provincia cuenta con su propio código procesal, porque la materia procesal es de competencia no delegada al Congreso Nacional. Existen procesos monitorios en algunas provincias y procesos abreviados en otras. En Córdoba el desalojo tramita por vía abreviada, con una impronta particular sobre el polo vulnerable, vinculada a la perspectiva de género cuando interviene una mujer en situación de vulnerabilidad. Estos códigos son similares entre sí, aunque algunos están más actualizados que otros; la Ciudad de Buenos Aires no necesariamente tiene el más moderno, y hay provincias ---como San Juan--- cuyo código está incluso más actualizado.

\section{Las seis etapas del proceso civil}

\begin{table}[H]
\centering
\caption{Las seis etapas del proceso civil}
\begin{tabularx}{\textwidth}{>{\raggedright\arraybackslash}p{0.22\textwidth} X}
\toprule
\textbf{Etapa} & \textbf{Contenido} \\
\midrule
Introductoria & Demanda, contestación, reconvención y su contestación, y resolución de excepciones previas. \\
Preliminar & Audiencia del art.\ 360: se fijan sujetos y objeto, los hechos controvertidos y los medios de prueba admisibles, se depura el proceso y el juez actúa como conciliador. \\
Probatoria & Producción de los medios admitidos: documental, informativa, pericial, testimonial, confesional y reconocimiento judicial. \\
Conclusional & Alegatos de las partes sobre el mérito de la prueba producida. \\
Recursiva & Dictada la sentencia, eventual interposición de recursos por la parte disconforme. \\
Ejecutoria & Firme la sentencia, cumplimiento y ejecución forzada si no hay cumplimiento voluntario. \\
\bottomrule
\end{tabularx}
\end{table}

\section{Las causales de desalojo}

Las principales causales son la falta de pago, el vencimiento del contrato, el cambio de destino, el comodato vencido, la intrusión y el uso indebido del inmueble. En general son causales \textbf{objetivas}: el incumplimiento del pago o el vencimiento del plazo son hechos verificables, sin mayor margen de discusión salvo que el demandado plantee una reconvención. Las tres causales de uso más frecuente comparten la misma vía procesal, pero requieren pruebas y planteos distintos.

\begin{table}[H]
\centering
\caption{Causales de desalojo, requisito previo y norma de referencia}
\begin{tabularx}{\textwidth}{>{\raggedright\arraybackslash}p{0.17\textwidth} >{\raggedright\arraybackslash}X >{\raggedright\arraybackslash}X >{\raggedright\arraybackslash}X}
\toprule
\textbf{Causal} & \textbf{Fundamento} & \textbf{Requisito previo} & \textbf{Normas citadas} \\
\midrule
Falta de pago & El locatario no abonó el canon en el plazo pactado. & Intimación fehaciente y diez días corridos; dos períodos adeudados. & CCCN, arts.\ 1219 y ss.; CPCCN, art.\ 684 bis. \\
Vencimiento de contrato & Se cumplió el plazo y no se restituyó. & No es obligatoria; se recomienda carta documento. & CCCN, arts.\ 1219 a 1227; CPCCN, art.\ 684. \\
Cambio de destino & El inmueble se usa para un fin distinto del pactado. & No obligatoria; se recomienda intimar al cese. & CPCCN, art.\ 680 ter (reconocimiento judicial). \\
Intrusión & El ocupante ingresó sin título ni contrato. & No requiere intimación contractual; se acredita la ausencia de vínculo. & CPCCN, art.\ 680 bis. \\
\bottomrule
\end{tabularx}
\end{table}

\begin{important}
Los artículos \textbf{680 bis} y \textbf{684 bis} del CPCCN prevén vías prácticamente gemelas de trámite abreviado, pero dirigidas a sujetos distintos: el 680 bis, contra \textbf{intrusos}, sin ningún vínculo contractual; el 684 bis, contra \textbf{locatarios, sublocatarios y ocupantes} que mantuvieron o mantienen alguna relación jurídica con el inmueble. La redacción actual del 684 bis eliminó la exigencia, vigente antes, de destino habitacional: hoy se aplica sin distinción de destino. Contra intrusos la acreditación es más simple porque no hay contrato que analizar. La figura del art.\ 684 bis se desarrolla en el capítulo~\ref{cap:684bis}.
\end{important}

La falta de pago requiere, además, los dos períodos adeudados y la intimación estudiados en el capítulo~\ref{cap:intimacion}. El proyecto de ley que pretendía reducir el plazo de intimación a tres días mantuvo, tras la media sanción, los diez días vigentes. La prueba clave es la exhibición de los últimos recibos y la acreditación del plazo de espera tras la intimación. Para no perder la causal conviene imputar los pagos parciales a servicios, impuestos y expensas, manteniendo siempre vigente la deuda por canon locativo.

\subsection{El reconocimiento judicial del artículo 680 ter}

\begin{formula}{Artículo 680 ter del CPCCN ---reconocimiento judicial}
Cuando el desalojo se fundare en las causales de cambio de destino, deterioro del inmueble, obras nocivas o uso abusivo o deshonesto, el juez deberá realizar antes del traslado de la demanda un reconocimiento judicial dentro de los cinco días de dictada la primera providencia, con asistencia del Defensor Oficial. La misma previsión rige cuando se diera la causal de los artículos 680 bis y 684 bis.
\end{formula}

% TODO: verificar referencia presente en las fuentes originales. La fuente incluye una nota de actualización según la cual el 7 de agosto de 2026 el Senado dio media sanción a un proyecto que propone modificar este artículo; debe verificarse el texto vigente al momento de aplicarlo.
El juez puede delegar la diligencia en otro empleado del juzgado. Su finalidad es que el juzgado verifique directamente los hechos antes de avanzar, lo que permite tramitar estos desalojos con mayor celeridad. Es una modalidad particular que el abogado debe conocer y solicitar expresamente.

\begin{table}[H]
\centering
\caption{Acta notarial y reconocimiento judicial}
\begin{tabularx}{\textwidth}{>{\raggedright\arraybackslash}p{0.20\textwidth} >{\raggedright\arraybackslash}X >{\raggedright\arraybackslash}X}
\toprule
\textbf{Aspecto} & \textbf{Acta notarial} & \textbf{Reconocimiento judicial (art.\ 680 ter)} \\
\midrule
Quién lo realiza & Un escribano contratado por el locador. & El juez, que puede delegar en otro empleado. \\
Carácter & Opcional; tiene fe pública. & Obligatorio en todo cambio de destino. \\
Efecto & Aporta prueba y acelera la acción; no reemplaza el reconocimiento judicial. & Se hace antes del traslado, dentro de los cinco días. \\
\bottomrule
\end{tabularx}
\end{table}

\begin{example}{Cambio de destino: fiesta electrónica en un departamento}
Un joven alquiló un departamento en un edificio con amenities y lo convirtió en lugar de reunión y de fiestas electrónicas, con malestar de los vecinos pese a las denuncias policiales. El reglamento del edificio establece como destino la vivienda, y ese uso configura un cambio de destino que habilita el desalojo. El horario o la frecuencia de la actividad no alteran la calificación jurídica, aunque inciden en la prueba. Otro ejemplo habitual es el consultorio de un profesional de la salud instalado en una unidad cuyo reglamento de copropiedad prohíbe ese uso.
\end{example}

\subsection{Condena de futuro}

\begin{formula}{Artículo 688 del CPCCN ---condena de futuro}
Permite iniciar el desalojo antes de que se produzca el vencimiento del contrato, obteniendo una sentencia que condene al locatario a restituir el inmueble en la fecha futura en que opera el vencimiento contractual. Es una figura que muchos colegas desconocen.
\end{formula}

\begin{example}{La utilidad estratégica de la condena de futuro}
Si el contrato vence en una fecha del año próximo, el locador puede iniciar desde ya el desalojo con fundamento en el art.\ 688, adelantándose a los tiempos del proceso: llegado el vencimiento, ya cuenta con una sentencia que ordena la restitución. Procede solo respecto del vencimiento del contrato, y debe correrse traslado al demandado como en cualquier desalojo. Corresponde explicar al cliente esta posibilidad y sus tiempos reales, para que decida si conviene adelantar el proceso aun sabiendo que la sentencia definitiva puede demorar.
\end{example}

\subsection{Desalojo y daños: acciones distintas}

El desalojo se limita a la restitución de la persona ocupante, sin atender al estado del inmueble. Si el inmueble resulta dañado, el reclamo por daños es una acción distinta, para la que el estado del inmueble deberá acreditarse luego. Tampoco puede el locador desalojar por su cuenta aunque el contrato contenga una cláusula resolutoria: <<no tenés el monopolio de la fuerza estatal>>; hasta quien ocupa tiene derechos (véase el capítulo~\ref{cap:vias-de-hecho}).

\begin{sintesis}
El desalojo es un proceso especial que tramita, en la Nación, por la vía sumarísima en la mayoría de los casos; en la provincia de Buenos Aires la vía se define en la primera resolución. Las causales son objetivas, cada una con requisitos previos propios; el cambio de destino exige reconocimiento judicial previo (art.\ 680 ter) y el vencimiento permite la condena de futuro (art.\ 688).
\end{sintesis}

\begin{cornell}{Cuadro Cornell del capítulo~\ref{cap:vias}}
\cqa{¿Qué vía procesal tiene el desalojo en el orden nacional?}{Es un proceso especial que puede tramitar como ordinario o sumarísimo; la mayoría lo hace por sumarísimo. La remisión histórica al proceso sumario está derogada por la Ley 25.488.}
\cqa{¿Qué cambia en la provincia de Buenos Aires?}{Subsisten ordinario, sumario y sumarísimo, y el juez fija la vía en la primera resolución; conviene identificarla de inmediato por la diferencia de plazos.}
\cqa{¿Cuáles son las etapas del proceso civil?}{Introductoria, preliminar, probatoria, conclusional, recursiva y ejecutoria.}
\cqa{¿Qué distingue a los arts.\ 680 bis y 684 bis?}{Son vías gemelas de trámite abreviado: el 680 bis, contra intrusos; el 684 bis, contra quien tuvo o tiene relación jurídica con el inmueble.}
\cqa{¿Cuándo procede el reconocimiento judicial del art.\ 680 ter?}{En cambio de destino, deterioro, obras nocivas o uso abusivo o deshonesto, y en los casos de los arts.\ 680 bis y 684 bis; antes del traslado, en cinco días, con el Defensor Oficial.}
\cqa{¿Qué es la condena de futuro?}{Iniciar el desalojo antes del vencimiento para obtener una sentencia que condene a restituir en esa fecha; solo procede por vencimiento (art.\ 688).}
\cresumen{Cada causal tiene fundamento, requisito previo y norma propios. Elegir bien la vía y la causal, reunir la prueba desde el inicio y anticipar el reconocimiento judicial o la condena de futuro condiciona el resultado.}
\end{cornell}

\chapter{La demanda de desalojo}\label{cap:demanda}

La demanda es el acto que da inicio formal al proceso, y su redacción concentra buena parte de las decisiones estratégicas: qué hechos se afirman, qué prueba se ofrece y cómo se encuadra la pretensión. Este capítulo estudia el concepto de demanda, su forma conforme al artículo 330 del CPCCN, el modo de presentarla y de notificarla, y la estructura de cada uno de sus apartados.

\section{Concepto y actos previos}

\begin{definition}{Demanda y pretensión}
La demanda es la base del proceso: una petición mediante la cual el actor pide algo concreto al juez, ejerciendo lo que técnicamente se denomina \textbf{pretensión}. Es el acto que inicia formalmente el proceso, aunque existen actos previos, como las diligencias preliminares (véase el capítulo~\ref{cap:intimacion}). El \emph{domicilio} no es, en principio, un dato que justifique una diligencia preliminar: se ubica con medios propios.
\end{definition}

\begin{important}
Conviene escribir lo mínimo indispensable. Cuanto más extenso es un escrito, mayor es el riesgo de inconsistencias entre pasajes, que pueden debilitar la posición procesal y dar argumentos a la contraparte.
\end{important}

\section{La forma de la demanda: artículo 330 del CPCCN}

\begin{formula}{Artículo 330 del CPCCN ---forma de la demanda}
La demanda será deducida por escrito y contendrá: 1) el nombre y domicilio del demandante; 2) el nombre y domicilio del demandado; 3) la cosa demandada, designándola con toda exactitud; 4) los hechos en que se funde, explicados claramente; 5) el derecho expuesto sucintamente, evitando repeticiones innecesarias; 6) la petición en términos claros y positivos.
\end{formula}

Aunque la prueba no integra la enumeración, es en la práctica el elemento que más pesa en el resultado: puede elaborarse una argumentación jurídica sólida, pero si la prueba ofrecida es deficiente la posición procesal queda comprometida. Con el escrito postulatorio debe acompañarse toda la documentación que esté en poder de la parte; la que no lo esté debe individualizarse para su incorporación.

\section{Presentación: firma, sorteo y encabezado}

El sorteo del juzgado no debe solicitarse antes de que la demanda esté firmada por todos los clientes, porque hasta entonces podrían cambiar de abogado, llegar a un arreglo o arrepentirse, y el sorteo asigna el número de expediente. Con patrocinio letrado la firma es hológrafa y deben firmar todos los actores. Si el escrito no se sube al sistema luego del sorteo, este queda sin efecto y debe repetirse.

\subsection{Carátula y objeto}

\begin{definition}{Carátula}
Es el título del escrito: se ubica al margen, sin sangría, a diferencia del resto de los párrafos, que llevan sangría de primera línea. Ejemplos: <<Promueve demanda de desalojo por falta de pago>>, <<por vencimiento de contrato>>. La carátula oficial la asigna el juzgado: nombre del actor, <<c/>> o <<contra>>, nombre del demandado y objeto (<<Fulano c/ Mengano s/ desalojo por falta de pago>>); se consigna primero el apellido, luego el nombre, y <<y otros>> cuando hay más partes.
\end{definition}

La expresión <<otras causales>> proviene del listado de objetos del Colegio de la Abogacía usado para caratular el formulario de sorteo. En el rubro desalojo ese listado distingue falta de pago, vencimiento de contrato, comodato, intrusos y otras causales; el cambio de destino se ubica en esta última, y el listado se actualiza periódicamente.

\subsection{Encabezado y domicilios}

El encabezado identifica a quien se presenta (por derecho propio, patrocinado o mediante apoderado), el domicilio real y el domicilio constituido, el domicilio electrónico y el objeto de la presentación. Debe constituirse siempre el domicilio procesal o \emph{ad litem}, dentro del radio del juzgado donde se sortea la causa, y el domicilio electrónico (en el orden nacional, el CUIT del abogado; en el provincial, el nombre seguido del CUIT). No corresponde consignar el domicilio legal, que responde a otros supuestos (societario, del curador y del curado, del representante legal de menores). Incluir WhatsApp o correo no es obligatorio, pero facilita que la contraparte continúe negociando. En el encabezado se consigna, entre paréntesis, la zona de notificación del domicilio constituido.

La fórmula <<respetuosamente digo>> es de cortesía y no resulta jurídicamente indispensable, ya que el trato respetuoso debe darse por sentado; por lo demás, el artículo 5 de la Ley 23.187 equipara a los abogados con los magistrados en cuanto a la consideración y respeto que se les debe.

\begin{table}[H]
\centering
\caption{Formas de comparecer en el proceso}
\begin{tabularx}{\textwidth}{>{\raggedright\arraybackslash}p{0.22\textwidth} X}
\toprule
\textbf{Forma} & \textbf{Explicación} \\
\midrule
Derecho propio & La parte se presenta personalmente, asistida por un letrado patrocinante obligatorio (solo abogados o procuradores matriculados). \\
Patrocinio & El abogado asiste técnicamente a la parte, que firma junto con el letrado. \\
Apoderado & Mediante poder, el abogado representa y actúa en nombre de la parte; puede reunir a la vez ambos roles. \\
\bottomrule
\end{tabularx}
\end{table}

Las personas jurídicas no pueden presentarse por derecho propio: necesitan un representante legal ---el presidente de una asociación o un socio con facultades suficientes---, que debe acreditar su carácter con copias certificadas del instrumento constitutivo, del acta de asamblea y del acta de directorio, y actuar asistido por letrado.

\section{Estructura de los apartados}

\subsection{Objeto: la pretensión}

El objeto debe indicar a qué se viene, sin que el juez deba inferirlo: <<cambio de destino>>, contra el locatario y sus subinquilinos y ocupantes, con el domicilio del inmueble. Técnicamente la pretensión se compone de tres elementos: el \textbf{subjetivo} (quién reclama y a quién), el \textbf{objetivo} (qué se pide: en el desalojo, la restitución del inmueble) y la \textbf{causa} (los hechos y el fundamento de derecho, mencionados genéricamente y desarrollados luego). Un ejemplo: que se requiere la restitución del inmueble en el plazo de ley, por la falta de pago de los alquileres. Puede cerrarse en <<con costas>> o agregar una breve referencia normativa.

\subsection{Competencia y legitimación activa}

La competencia debe incluirse siempre y exige leer antes el contrato completo, por la posible prórroga (véase el capítulo~\ref{cap:competencia}); a falta de ella se aplica el artículo 5 del CPCCN: es competente el juez del lugar de situación del inmueble. El apartado de legitimación activa es optativo: se incluye como estrategia para anticiparse a la excepción de falta de legitimación (véase la sección~\ref{sec:leg-activa}).

\subsection{Los hechos}

\begin{example}{La demanda como película}
La demanda no es un cuento que debe anticipar de entrada cómo termina: se arma como una película, con una fotografía de la situación inicial, luego qué pasó y recién al final la conclusión. Los hechos deben encuadrar en la conclusión que se pretende, teniendo siempre como norte la norma invocada a favor del actor.
\end{example}

\begin{definition}{Subsunción jurídica}
Al sentenciar, el juez verifica si los hechos alegados y acreditados encuadran en el supuesto de hecho abstracto de una norma; si encajan, aplica esa norma general al caso, dictando una sentencia individual. Constatar los hechos, verificar que una norma los contempla y aplicarla cierra la norma sobre el caso concreto.
\end{definition}

Se describe en apartados separados la relación locativa (plazo, cláusula y destino pactado) y el incumplimiento (por ejemplo, el cambio de destino). En el desalojo por vencimiento, el relato incluye el contrato con sus fechas y la falta de restitución, y se complementa con la intimación realizada por carta documento.

\subsection{Subinquilinos y ocupantes}

Conviene incluir en toda demanda una fórmula que manifieste que se desconoce la existencia de sublocatarios u otras personas en el inmueble, distintas de los codemandados. Si la demanda reconoce a otros ocupantes, deben individualizarse; no reconocerlos evita que el juzgado, por el principio de defensa en juicio del artículo 18 de la Constitución Nacional, exija notificar individualmente a cada uno. Así, la demanda se dirige <<contra subinquilinos y ocupantes>>, y es el oficial de justicia quien constata el día del lanzamiento quién habita el inmueble (véase el capítulo~\ref{cap:terceros}).

\subsection{El derecho}

Se citan las normas del CCCN sobre locación y las del Código Procesal; también pueden citarse doctrina y jurisprudencia. Si se invoca una norma equivocada, la demanda no se tiene por no presentada, por el principio \emph{iura novit curia}; de todos modos conviene evitar errores.

\subsection{La planilla de hechos y medios de prueba}

Como herramienta práctica se recomienda armar para cada escrito una planilla de dos columnas: los hechos alegados, numerados, y frente a cada uno, el medio de prueba ofrecido. Los medios son documental, pericial, informativa, confesional, testimonial y reconocimiento judicial (véase el capítulo~\ref{cap:prueba}). Toda la prueba va con la demanda: con posterioridad solo pueden agregarse documentos de fecha posterior, dado que el sistema vigente concentra la prueba en los escritos postulatorios tras la eliminación del juicio sumario. Conviene ofrecer toda la prueba conducente frente a los hechos que previsiblemente controvertirá la contraparte; por ejemplo, prueba informativa mediante oficio a las cámaras del Gobierno de la Ciudad. Antes, en el proceso ordinario, primero se corrían los traslados y se planteaban las excepciones, y recién después se producía la prueba.

\section{Domicilio desconocido y nulidad de la notificación}

\begin{important}
Cuando el actor no conoce el domicilio del demandado debe agotar las diligencias razonablemente exigibles antes de notificar por edictos a persona incierta: si se libró un oficio o cédula con resultado negativo, corresponde oficiar a organismos públicos ---por ejemplo, la Policía Federal y la Cámara Nacional Electoral--- para dar con el domicilio real.
\end{important}

\begin{example}{La chapa borrada}
En los desalojos es habitual que el notificador sea informado por un vecino de que el demandado no vive allí, o que alguien haya retirado o modificado la numeración de la puerta. Por eso las cédulas y mandamientos suelen exigir que se individualice el inmueble con puntos de referencia adicionales.
\end{example}

Si el oficial notificador no logra ingresar al edificio, debe procurar el ingreso con auxilio de la fuerza pública; si aun así no lo dejan, deja aviso y vuelve al día siguiente; si en esa segunda oportunidad tampoco es atendido, deja la cédula por debajo de la puerta y el demandado queda notificado. La imposibilidad material de ingresar no puede paralizar el proceso. Se aconseja respetar el domicilio contractual o electrónico pactado por razones de seguridad jurídica.

\begin{important}
Si años después de iniciado el proceso se comprueba que la notificación inicial fue defectuosa, se anula lo actuado y debe notificarse nuevamente, volviendo a empezar. Aun así, la demanda mantiene su efecto interruptivo de la prescripción ---lo que se pierde es tiempo procesal, no el derecho de fondo--- y las costas del proceso nulo las afronta quien resulte responsable de la notificación defectuosa.
\end{important}

\begin{sintesis}
La demanda debe cumplir el art.\ 330 CPCCN, presentarse firmada por todos los actores y contener un objeto claro, hechos ordenados como relato, derecho sucinto y toda la prueba. Debe cuidarse la notificación, porque su defecto puede anular años de proceso.
\end{sintesis}

\begin{cornell}{Cuadro Cornell del capítulo~\ref{cap:demanda}}
\cqa{¿Qué exige el art.\ 330 CPCCN?}{Escrito con nombre y domicilio de las partes, cosa demandada con exactitud, hechos claros, derecho sucinto y petición clara y positiva.}
\cqa{¿Por qué no se sortea antes de la firma de todos los actores?}{Porque el sorteo asigna el número de expediente y hasta la firma los clientes pueden desistir o cambiar de abogado; si no se sube el escrito, el sorteo queda sin efecto.}
\cqa{¿Qué domicilios se constituyen y cuál no?}{Procesal o \emph{ad litem} y electrónico, siempre; el domicilio legal no se coloca.}
\cqa{¿Cuáles son los tres elementos de la pretensión?}{Subjetivo (partes), objetivo (qué se pide: restitución) y causa (hechos y derecho).}
\cqa{¿Por qué conviene no reconocer subinquilinos y ocupantes individualmente?}{Para evitar que el juzgado, por la defensa en juicio (art.\ 18 CN), exija notificar a cada uno; el oficial de justicia constata quién habita el día del lanzamiento.}
\cqa{¿Cuándo debe ofrecerse la prueba?}{Con la demanda; después solo documentos de fecha posterior.}
\cqa{¿Qué efecto tiene la nulidad por notificación defectuosa?}{Se anula lo actuado y se notifica de nuevo, pero la demanda conserva el efecto interruptivo de la prescripción y las costas las asume el responsable.}
\cresumen{La demanda ordena hechos, derecho y prueba en un escrito preciso y firmado por todos los actores. Su éxito depende de la prueba ofrecida desde el inicio y de una notificación correcta.}
\end{cornell}

\chapter{Traslado de la demanda y rebeldía}\label{cap:traslado}

Notificado el traslado de la demanda, el demandado enfrenta una sucesión de decisiones: comparecer, contestar, oponer excepciones. Este capítulo estudia los dos efectos de esa notificación ---citación y emplazamiento---, la noción de carga procesal y la preclusión, y la rebeldía, con sus requisitos y sus consecuencias sobre los hechos y los documentos.

\section{Citación y emplazamiento}

El traslado de la demanda notificado al demandado tiene dos caras que conviene distinguir.

\begin{definition}{Citación y emplazamiento}
La \textbf{citación para estar a derecho} impone al demandado la carga de comparecer y constituir domicilio (<<ponerse a derecho>>), bajo apercibimiento de rebeldía. El \textbf{emplazamiento para contestar} impone la carga de contestar la demanda, oponer excepciones y ofrecer prueba, dentro de un plazo determinado. Son dos cargas autónomas: quien comparece y constituye domicilio cumple con la citación aunque todavía no conteste, y el plazo del emplazamiento corre por separado.
\end{definition}

\begin{definition}{Carga procesal y preclusión}
La \textbf{carga procesal} es un imperativo de interés propio, no una obligación: se cumple si se quiere; si no se hace, se pierde la facultad, pero nadie puede exigirla coactivamente. Se diferencia de la \textbf{obligación} en sentido estricto. La \textbf{preclusión} es la pérdida, extinción o consumación de una facultad procesal por no haberse ejercido en el momento debido: el proceso avanza por etapas sucesivas y, superada una, no se puede volver atrás, salvo nulidad. Los plazos son perentorios: vencido el plazo, se pierde la facultad sin necesidad de pedido de la contraparte.
\end{definition}

Presentarse tarde tiene consecuencias: los actos procesales que ya precluyeron no pueden retrotraerse.

\section{La rebeldía del demandado}

\begin{definition}{Rebeldía}
Para que exista rebeldía en sentido técnico es necesario que el traslado haya sido notificado en un \textbf{domicilio conocido} ---con constancia cierta de que la persona vive allí--- y que, pese a ello, el demandado no comparezca a estar en derecho. No es rebelde quien no está fehacientemente notificado.
\end{definition}

\begin{important}
\textbf{Notificación en domicilio conocido más incomparecencia} configura rebeldía en sentido propio. \textbf{Notificación por edictos a persona incierta} (domicilio desconocido, agotadas las diligencias) \textbf{más incomparecencia} no es rebeldía en sentido estricto: corresponde la intervención del defensor oficial, porque sería demasiado grave que el proceso avanzara sin contralor sobre una persona de la que ni siquiera se sabe si conoció la demanda.
\end{important}

La rebeldía implica no presentarse, no simplemente no contestar: si el demandado se presenta ---aunque sea solo constituyendo domicilio--- ya no puede ser declarado rebelde, por más que no conteste. Si compareció por apoderado y luego este renuncia o le es revocado el mandato, no corresponde declarar la rebeldía sin más: el demandado debe ser intimado a designar un nuevo apoderado.

La rebeldía no opera de oficio: requiere pedido de parte, y conviene al actor porque facilita las medidas cautelares, que se pueden dictar sin acreditar peligro en la demora ni prestar contracautela. Altera además el régimen de notificaciones: desde la declaración, todas las resoluciones se notifican por ministerio de la ley (por nota), salvo la sentencia definitiva, que debe notificarse personalmente o por cédula. La resolución que declara la rebeldía queda firme, salvo apelación, al día hábil siguiente, y es ella misma la que se notifica. Ante la duda sobre si el demandado tuvo conocimiento efectivo, existe una primera presunción a favor del actor.

\subsection{Efectos de la incontestación sobre hechos y documentos: artículo 356}

\begin{formula}{Artículo 356 ---el matiz entre <<podrá>> y <<deberá>>}
Respecto de los \textbf{hechos}: si el demandado no los contesta categóricamente o guarda silencio, el juez \textbf{podrá} tenerlos por reconocidos. No es un efecto automático sino una facultad del juez, que debe valorar la prueba del expediente. Respecto de los \textbf{documentos} acompañados con la demanda: si el demandado no se expide expresamente sobre ellos, se los \textbf{deberá} tener por auténticos; es prueba tasada, con efecto predeterminado por la ley.
\end{formula}

\begin{example}{La demanda <<de puro derecho>> rechazada}
Un actor, ante la rebeldía del demandado, pidió que se dictara sentencia <<de puro derecho>>, sin abrir a prueba, confiado en que la rebeldía le garantizaba el triunfo. La sentencia rechazó la demanda: al pedirla de puro derecho el actor renunció a que se abra la causa a prueba, y la sola rebeldía, sin ningún principio de prueba que la respalde, no alcanzó para tener los hechos por acreditados.
\end{example}

\begin{important}
Aun frente a un demandado rebelde conviene abrir la causa a prueba y producir, por ejemplo, la absolución de posiciones, para instalar en el juez una duda razonable a favor del actor. No hace falta una certeza categórica: basta mostrar algo que, sin alcanzar para una sentencia favorable directa, sirva para generar esa duda. El rebelde pierde además la posibilidad de ofrecer prueba fuera de la oportunidad procesal debida, por la preclusión.
\end{important}

\begin{sintesis}
La notificación del traslado produce citación y emplazamiento, dos cargas distintas. La rebeldía exige domicilio conocido e incomparecencia y no equivale a no contestar. Sus efectos sobre los hechos dependen de la valoración del juez (<<podrá>>), por lo que conviene producir prueba igualmente.
\end{sintesis}

\begin{cornell}{Cuadro Cornell del capítulo~\ref{cap:traslado}}
\cqa{¿Qué diferencia hay entre citación y emplazamiento?}{La citación impone comparecer y constituir domicilio; el emplazamiento, contestar dentro de un plazo. Son cargas autónomas.}
\cqa{¿Qué es una carga procesal y en qué se diferencia de la obligación?}{Es un imperativo del propio interés: si no se cumple se pierde la facultad, pero no puede exigirse coactivamente.}
\cqa{¿Cuándo hay rebeldía en sentido técnico?}{Con notificación en domicilio conocido e incomparecencia; con edictos a persona incierta interviene el defensor oficial.}
\cqa{¿Puede ser rebelde quien se presentó pero no contestó?}{No: la rebeldía implica no presentarse; presentarse y constituir domicilio la impide.}
\cqa{¿Qué efectos tiene la declaración de rebeldía?}{Facilita cautelares (sin peligro en la demora ni contracautela) y las resoluciones se notifican por nota, salvo la sentencia definitiva.}
\cqa{¿Qué diferencia hay entre <<podrá>> y <<deberá>> en el art.\ 356?}{Los hechos no contestados el juez podrá tenerlos por ciertos (facultad); los documentos no desconocidos deberá tenerlos por auténticos (prueba tasada).}
\cresumen{Ponerse a derecho y contestar son cargas distintas. La rebeldía solo opera con notificación cierta y pedido de parte, y aun entonces el actor debe producir prueba para sostener su pretensión.}
\end{cornell}

\chapter{Excepciones previas y contestación de la demanda}\label{cap:excepciones}

Contestar la demanda es la decisión de mayor aplicación práctica para el abogado del demandado: qué defensas oponer, cómo negar los hechos sin quedar expuesto a una presunción en contra y cómo organizar el escrito. Este capítulo explica los presupuestos procesales, la clasificación de las defensas, el cómputo de los plazos, la estructura de la contestación y la carga de negar categóricamente del artículo 356.

\section{Presupuestos procesales y excepciones previas}

Las excepciones previas se vinculan con los \textbf{presupuestos procesales}: las condiciones que deben darse para que el proceso avance válidamente. Se requiere un juez competente ---si no lo hay, no hay proceso---; no puede haber dos procesos por la misma causa (litispendencia) y se exige capacidad procesal de las partes. Las defensas se agrupan en dos familias: las procesales (incompetencia, falta de legitimación, litispendencia) y las que remiten al derecho de fondo (cosa juzgada, pago, prescripción, excepción de incumplimiento contractual). A ellas se suma la caducidad de instancia, que opera cuando la parte que debía impulsar el proceso no lo hace. Todo ello ---contestación y excepciones--- constituye el derecho de defensa del demandado.

\begin{definition}{La reconvención}
Más que una defensa, es un contraataque: el demandado introduce, dentro del mismo proceso, una pretensión propia contra el actor, que debe ser conexa con la de la demanda principal. No existe la <<reconvención de la reconvención>>.
\end{definition}

Se llaman excepciones de \textbf{previo y especial pronunciamiento} porque se plantean y resuelven antes de la sentencia, mediante un trámite especial: no tendría sentido tramitar un proceso entero para advertir al final que el juez era incompetente.

\section{Clasificación doctrinaria: impedimentos, excepciones y defensas}

La obra de referencia es \emph{La demanda civil}, de Carlos Carli, que clasifica, retomando también la posición de Palacio, las defensas en tres categorías.

\begin{table}[H]
\centering
\caption{Impedimentos procesales, excepciones y defensas (Carli, con Palacio)}
\begin{tabularx}{\textwidth}{>{\raggedright\arraybackslash}p{0.22\textwidth} >{\raggedright\arraybackslash}X >{\raggedright\arraybackslash}X}
\toprule
\textbf{Categoría} & \textbf{Qué ataca} & \textbf{Ejemplos} \\
\midrule
Impedimentos procesales & Faltan los presupuestos procesales; no se dice nada del derecho ni de la acción. & Incompetencia, litispendencia, falta de personería, falta de arraigo, defecto legal. \\
Excepciones propiamente dichas & La acción, no el proceso ni el derecho de fondo. & Prescripción, cosa juzgada, transacción y conciliación previas. \\
Defensas & Directamente el derecho del actor. & Pago de la obligación reclamada. \\
\bottomrule
\end{tabularx}
\end{table}

\begin{table}[H]
\centering
\caption{Presupuesto procesal y excepción correspondiente}
\begin{tabularx}{\textwidth}{>{\raggedright\arraybackslash}X >{\raggedright\arraybackslash}X}
\toprule
\textbf{Presupuesto procesal} & \textbf{Excepción} \\
\midrule
Juez competente & Incompetencia \\
Proceso único (no litispendencia) & Litispendencia \\
Capacidad procesal / representación & Falta de personería \\
Arraigo (hoy casi en desuso, salvo domicilio o bienes fuera del país) & Falta de arraigo \\
Demanda clara y precisa & Defecto legal en el modo de proponer la demanda \\
\bottomrule
\end{tabularx}
\end{table}

El pago, aunque cuestión de fondo, puede oponerse como excepción previa cuando es de fácil comprobación. Ciertas excepciones nominadas ---falta de legitimación, prescripción--- solo tramitan como previas cuando la cuestión es \textbf{manifiesta}; si requiere prueba, se resuelve con el fondo en la sentencia.

\begin{important}
Toda excepción previa se opone junto con la contestación, en el mismo escrito. Si se omite esa carga en la primera oportunidad, opera la preclusión. En el proceso \textbf{sumarísimo} ---por ejemplo, en el desalojo--- no se admiten excepciones de previo y especial pronunciamiento: todas las defensas, previas o de fondo, se resuelven juntas en la sentencia.
\end{important}

\section{Cómputo de plazos y efectos del planteo de excepciones}

\begin{exercise}{Cómputo de plazos}
\textbf{Supuesto:} se notifica la demanda el lunes 31 y el día hábil siguiente es feriado; el plazo para contestar en el ordinario es de quince días hábiles. \textbf{Regla:} en todos los medios de notificación (nota, cédula, acta notarial, telegrama) el día de la notificación no se cuenta; el plazo corre desde el día hábil siguiente, sin contar sábados, domingos ni feriados. \textbf{Resolución:} si el día inmediato es feriado, no se cuenta, y el próximo día hábil es el <<día uno>> del plazo. Es el punto donde más se equivocan en los exámenes: se confunde el día de la notificación con el primer día del plazo.
\end{exercise}

Al oponer excepciones junto con la contestación, el plazo para contestar queda \textbf{suspendido} hasta que se resuelvan; al reanudarse corre el remanente.

\begin{important}
Si se \textbf{rechaza} la excepción de defecto legal, el plazo se reanuda por el remanente, los días que faltaban al oponerla. Si se \textbf{hace lugar} a ella (porque la demanda era realmente oscura y el actor debió corregirla), el plazo se interrumpe y el demandado cuenta con el plazo completo desde la notificación de la demanda corregida.
\end{important}

\begin{note}
La excepción de defecto legal puede tener finalidad estratégica: aunque se sepa que será rechazada, plantearla permite ganar tiempo, porque mientras se sustancia el plazo queda suspendido. Se cargará con las costas del incidente, pero se habrá ganado tiempo para preparar la defensa; por eso se la utiliza en ocasiones más como mecanismo dilatorio que con su finalidad propia.
\end{note}

\section{Estructura del escrito de contestación}

El escrito guarda paralelo con la demanda: en lugar de <<interpone demanda>> se dice <<contesta demanda>>, con el mismo destinatario.

\begin{important}
\textbf{Estructura del escrito de contestación:}
\begin{enumerate}
\item Sumilla (<<Contesta demanda. Opone excepciones. Ofrece prueba>>).
\item Carátula del expediente y número, si se conoce.
\item Personería: por derecho propio o en representación, acompañando el poder.
\item Domicilios: procesal (físico y electrónico) y real, con datos del letrado.
\item Objeto: contestar en tiempo y forma, solicitando el rechazo con costas.
\item Capítulo de excepciones previas, si corresponde.
\item Contestación propiamente dicha: reconocimiento o negativa de los hechos, versión propia y encuadre jurídico.
\item Ofrecimiento de prueba, con capítulos separados para la de las excepciones y la del fondo.
\item Petitorio: es el capítulo que más se lee.
\end{enumerate}
\end{important}

\begin{example}{Prórroga de competencia planteada en la contestación}
Si el demandado entiende que el juez no es competente, en lugar de contestar sin más plantea una cuestión de competencia pidiendo la remisión al juez que considera competente. Si un inmueble está en Córdoba y el contrato pactó la competencia de los tribunales de la Ciudad Autónoma de Buenos Aires, la prórroga territorial es válida por tratarse de una cuestión patrimonial disponible. Distinto es cuando la competencia no es prorrogable, como la federal frente a un asunto local: el juez puede declararse incompetente de oficio.
\end{example}

\section{El artículo 356: la carga de negar categóricamente}

\begin{formula}{Artículo 356 del CPCCN ---carga de reconocer o negar los hechos}
En la contestación el demandado opondrá todas las excepciones o defensas de que intente valerse y deberá: 1) reconocer o negar categóricamente cada uno de los hechos expuestos en la demanda, la autenticidad de los documentos acompañados que se le atribuyeren y la recepción de las cartas y telegramas a él dirigidos cuyas copias se acompañen; su silencio, sus respuestas evasivas o la negativa meramente general podrán estimarse como reconocimiento de la verdad de los hechos pertinentes y lícitos a que se refieren, y los documentos se tendrán por reconocidos o recibidos; el defensor oficial y el sucesor a título universal de quien participó en los hechos o suscribió los documentos pueden reservar su respuesta definitiva para después de producida la prueba; 2) especificar con claridad los hechos que alegare como fundamento de su defensa; 3) observar, en lo aplicable, los requisitos del artículo 330.
\end{formula}

No basta una fórmula genérica del tipo <<niego todos y cada uno de los hechos por no constarme>>: debe seguirse con una negativa hecho por hecho (<<especialmente niego\ldots>>). Por precaución conviene negar incluso los hechos aparentemente secundarios, porque un hecho sin negar puede completar, por presunción, la cadena de razonamiento del juez en contra del demandado. El reconocimiento o desconocimiento de los documentos debe ser expreso y puntual, documento por documento.

\begin{table}[H]
\centering
\caption{Actitudes frente a los hechos de la demanda (art.\ 356)}
\begin{tabularx}{\textwidth}{>{\raggedright\arraybackslash}p{0.28\textwidth} X}
\toprule
\textbf{Actitud} & \textbf{Consecuencia procesal} \\
\midrule
Reconocer el hecho & Queda fuera del debate probatorio (si el derecho es disponible). \\
Negar el hecho & Pasa a ser controvertido y queda dentro del objeto de la prueba: quien lo alegó deberá probarlo. \\
Silencio, negativa genérica o respuesta evasiva & <<Admisión técnica>>: el juez podrá tenerlo por cierto. Como <<podrá>> no es <<deberá>>, siempre conviene ofrecer y producir prueba. \\
\bottomrule
\end{tabularx}
\end{table}

\begin{example}{La norma y sus hechos antecedentes}
La norma dice que si se da A debe ser B, y B es lo que el actor reclamó: los hechos que integran el antecedente A deben estar probados para que el juez aplique la consecuencia B. Cada parte prueba el presupuesto de la norma que invoca a su favor, como fundamento de su pretensión o de su defensa.
\end{example}

\begin{sintesis}
Las defensas se ubican según ataquen el proceso (impedimentos), la acción (excepciones) o el derecho (defensas). En el desalojo sumarísimo todas se resuelven en la sentencia. La contestación debe negar categóricamente hecho por hecho; el silencio expone a una presunción en contra.
\end{sintesis}

\begin{cornell}{Cuadro Cornell del capítulo~\ref{cap:excepciones}}
\cqa{¿Qué ataca cada categoría de defensa según Carli?}{Impedimentos procesales, los presupuestos del proceso; excepciones propiamente dichas, la acción (prescripción); defensas, el derecho del actor (pago).}
\cqa{¿Qué presupuesto corresponde a cada excepción?}{Juez competente, incompetencia; proceso único, litispendencia; capacidad, falta de personería; arraigo, falta de arraigo; demanda clara, defecto legal.}
\cqa{¿Se admiten excepciones previas en el desalojo sumarísimo?}{No: todas las defensas se resuelven juntas en la sentencia.}
\cqa{¿Cómo se cuenta el plazo para contestar?}{En días hábiles, sin contar el día de la notificación; se suspende al oponer excepciones y, si se acoge el defecto legal, se interrumpe y corre completo de nuevo.}
\cqa{¿Qué exige el art.\ 356 al demandado?}{Reconocer o negar categóricamente cada hecho y la autenticidad de los documentos; el silencio o la negativa genérica pueden tenerse por reconocimiento.}
\cqa{¿Para qué sirve el defecto legal como estrategia?}{Para ganar tiempo, porque suspende el plazo de contestación, aunque se rechace y se carguen las costas del incidente.}
\cresumen{El demandado debe oponer en un solo escrito excepciones, contestación y prueba, negando categóricamente cada hecho. En el desalojo sumarísimo las defensas se resuelven en la sentencia, y los plazos exigen un cómputo cuidadoso.}
\end{cornell}

\chapter{La prueba en el juicio de desalojo}\label{cap:prueba}

Sin prueba no hay demanda que resista un juicio, y la prueba es el elemento que más pesa en el resultado. Este capítulo estudia la oportunidad de ofrecerla, la diferencia entre prueba directa e indiciaria y cada medio probatorio relevante en el desalojo: documental (incluida la impugnación de instrumentos públicos y privados), confesional, pericial, testimonial e informativa.

\section{Oportunidad y carga de la prueba}

El sistema procesal vigente exige acompañar u ofrecer toda la prueba junto con la demanda, y la contraria debe hacerlo con la contestación, cualquiera sea el tipo de proceso. Quien alega un hecho debe probarlo; cada parte prueba el presupuesto de la norma que invoca a su favor. Si la contraparte niega todo o el juez declara la causa de puro derecho, puede dictarse sentencia sin otra prueba cuando lo alegado ya surge acreditado; pero sin ninguna prueba el juez carece de modo de verificar, por ejemplo, que el contrato está vencido.

\begin{important}
Al ofrecer prueba debe anticiparse qué hechos negará la contraparte y qué medio los acreditará. Un documento privado no basta por sí solo si no es reconocido: se respalda con otro medio. La oportunidad de ofrecer prueba precluye y no puede subsanarse después.
\end{important}

\section{Prueba directa, hechos negados e indicios}

Conviene distinguir la prueba directa de un hecho de la prueba por indicios y presunciones. Los \textbf{indicios} no son prueba directa: son otros hechos que, interpretados armónicamente, de forma contundente y categórica, hacen presumir la existencia de un hecho distinto que no pudo probarse directamente. No se <<ofrecen>> como medio autónomo: se construyen en el alegato, a partir de la prueba ya producida.

\begin{example}{Cómo se prueba una simulación sin contradocumento}
Si un fiador transfirió simuladamente un local a su sobrino mediante un boleto de compraventa para eludir su fianza, el contradocumento por el cual el sobrino reconoce que la venta no es real nadie lo exhibiría. Se prueba reuniendo indicios precisos y concordantes ---nunca uno solo---: (1) la falta de capacidad económica del comprador, que podría verificarse requiriendo información a entidades bancarias o a la AFIP sobre el ingreso del dinero declarado; (2) el vínculo de parentesco entre tío y sobrino, indicio de connivencia no determinante por sí solo; y (3) la ausencia de entrega efectiva de la posesión, ya que el supuesto vendedor siguió explotando el local. No se prueba directamente la simulación, pero sí tres hechos que, en conjunto, permiten presumirla. La prueba indiciaria también se utiliza en otros terrenos, como los juicios de filiación, donde se suma el examen genético.
\end{example}

\section{Prueba documental}

\begin{definition}{Documento}
Es toda representación material apta para exteriorizar el pensamiento humano, sin importar el soporte físico. Un ejemplo es el mojón: una piedra que delimita dos terrenos en una mensura \emph{representa} algo humano con valor jurídico, y por eso constituye prueba documental.
\end{definition}

La prueba documental actual incluye comprobantes de pago, contratos de locación y otros documentos literales; mensajes de WhatsApp y correos electrónicos; y videos y fotografías, siempre que puedan acreditarse en su origen y contenido. Lo que importa es el soporte. La regla general es extraer la prueba de la fuente más directa posible: para un video, pedir informativa a la empresa que administra las cámaras de seguridad o resguardar el disco original, para no perder la cadena de custodia. Para probar el estado de un local conviene contar con fotos del antes y del después, tomadas con las partes presentes, incluso firmadas.

\begin{example}{El riesgo de no resguardar la prueba digital}
En un caso, un video relevante estaba grabado en un celular que luego se sobreescribió por no haberse resguardado a tiempo, y se perdió definitivamente. En otro, dentro de un proceso laboral, el abogado contrario envió un pendrive con un video; se solicitó formalmente a la Cámara que lo resguardara bajo su custodia hasta la audiencia de vista de causa, para incorporarlo válidamente sin que la otra parte pudiera desconocer su origen.
\end{example}

\subsection{Instrumentos públicos y privados: cómo se impugnan}

\begin{table}[H]
\centering
\caption{Impugnación de documentos según su naturaleza}
\begin{tabularx}{\textwidth}{>{\raggedright\arraybackslash}p{0.24\textwidth} >{\raggedright\arraybackslash}X >{\raggedright\arraybackslash}X}
\toprule
\textbf{Documento} & \textbf{Cómo se impugna} & \textbf{Observaciones} \\
\midrule
Instrumento público & Únicamente por \textbf{redargución de falsedad}, al contestar el traslado y dentro de los diez días de esa impugnación. & Goza de fe pública y se presume auténtico hasta que se lo declare falso por ese incidente. \\
Instrumento privado & Basta desconocerlo al contestar el traslado, sin incidente autónomo; si se desconoce la firma, se ofrece pericial caligráfica en ese acto. & Si se reconoce la firma, queda reconocido el documento íntegro, incluido el monto o la cláusula. \\
Instrumento privado con firma reconocida & Puede plantearse falsedad ideológica (el contenido no responde a la verdad) o falsedad material (el documento fue adulterado). & --- \\
Carta documento cuya autenticidad se discute & Prueba informativa al correo sobre autenticidad y fecha de recepción. & --- \\
\bottomrule
\end{tabularx}
\end{table}

Dentro de la fe pública, la \textbf{falsedad ideológica} supone que el escribano faltó a la verdad de lo que dice haber presenciado (por ejemplo, que se entregó el dinero en el acto): para atacarla hace falta el incidente de redargución. Las simples manifestaciones de las partes recogidas por el escribano ---no los hechos que presenció--- pueden impugnarse por prueba en contrario, sin ese incidente. Frente a cada traslado hay que analizar qué tipo de documento es.

\begin{important}
Si la firma del contrato de locación está certificada ante escribano público ---práctica especialmente recomendada---, la pericial caligráfica resulta innecesaria, porque un funcionario con fe pública avala la firma. La fuente indicó, al momento de la clase, un costo aproximado de sesenta mil pesos por esa certificación, e indicó además que con el contrato certificado queda preparada la vía ejecutiva para el cobro de alquileres.
\end{important}

\section{Prueba confesional}

\begin{definition}{Prueba confesional y pliego de posiciones}
Se provoca por vía oral mediante afirmaciones, llamadas \textbf{posiciones}, que la contraparte debe reconocer como ciertas o negar bajo juramento. Se redactan en un \textbf{pliego}, que se presenta en sobre cerrado, firmado y notificado a la contraria, \textbf{hasta media hora antes de la audiencia}; pasado ese momento <<el pliego se murió>>. Quien lo redacta es el \textbf{ponente}; quien lo contesta, el \textbf{absolvente}.
\end{definition}

Las proposiciones son afirmativas (<<Para que jure como es cierto que\ldots>>), versan sobre los hechos centrales (que no abonó el alquiler, que ocupa indebidamente el inmueble, que está obligado a desalojar) y deben ser favorables a quien las formula y desfavorables a quien responde. Si el citado, debidamente notificado, no comparece, se lo tiene por confeso de todas las afirmaciones (\emph{ficta confessio}). La ausencia del absolvente no impide la celebración de la audiencia preliminar. La prueba resulta útil incluso frente a un demandado rebelde. Cada posición implica en sí misma un reconocimiento para quien la formula. No debe confundirse con el interrogatorio a testigos, que se dirige a terceros ajenos al litigio.

\begin{note}
La fuente informa que, aunque el Código prevé la presentación escrita del pliego, las posiciones pueden formularse a viva voz en la audiencia, y que la absolución de posiciones cayó en desuso: los abogados suelen acordar no tomarlas y los clientes llegan preparados para responder.
\end{note}

\section{Prueba pericial y consultores técnicos}

La pericial se reserva para cuestiones técnicas, como el desconocimiento de una firma. El perito oficial se designa por sorteo entre los inscriptos en la lista correspondiente (hoy informático; antes, con bolillas en un bolillero). Cada parte puede designar además su \textbf{consultor técnico}, que propone los puntos de pericia, controla la tarea del perito oficial y puede presentar un informe propio si discrepa. Un profesional del derecho no está capacitado para dictaminar sobre puntos médicos, de mala praxis o caligráficos. Importa fijar con precisión los puntos de pericia. Quienes incumplen sus deberes ---escribanos o peritos--- pueden ser sometidos a proceso disciplinario o penal, sin que ello altere el valor de la fe pública ni de la pericial como institución.

\section{Prueba testimonial e informativa}

Puede ofrecerse el testimonio de personas ajenas al proceso; la contraparte también puede pedir que se cite a un tercero, y si este no comparece, esa omisión puede jugar como indicio desfavorable a quien lo propuso. La \textbf{prueba informativa} se solicita por oficios a organismos públicos o entidades privadas para que informen datos objetivos en su poder. El petitorio final sintetiza todo lo pedido y debe consignarse completo, porque muchas veces es lo único que se lee con detenimiento. El cargo judicial certifica fecha y hora de presentación; con la presentación electrónica lo asigna automáticamente el sistema.

\begin{sintesis}
Toda la prueba se ofrece con los escritos postulatorios. La documental es el pilar del desalojo; su impugnación depende de si el instrumento es público o privado. La confesional se prepara con pliego hasta media hora antes de la audiencia, y la indiciaria se construye en el alegato.
\end{sintesis}

\begin{cornell}{Cuadro Cornell del capítulo~\ref{cap:prueba}}
\cqa{¿Cuándo se ofrece la prueba?}{Con la demanda y con la contestación; precluye y no puede subsanarse luego.}
\cqa{¿Qué es un indicio y cómo se usa?}{Un hecho que, junto con otros precisos y concordantes, permite presumir un hecho no probado directamente; se construye en el alegato, no se ofrece como medio.}
\cqa{¿Qué es un documento?}{Toda representación material apta para exteriorizar el pensamiento humano, cualquiera sea el soporte: contratos, mensajes, videos, fotos.}
\cqa{¿Cómo se impugna un instrumento público y uno privado?}{El público, solo por redargución de falsedad en diez días; el privado, desconociéndolo al contestar el traslado y, si se desconoce la firma, con pericial caligráfica.}
\cqa{¿Qué efecto tiene la certificación notarial de la firma?}{Hace innecesaria la pericial caligráfica y prepara la vía ejecutiva para cobrar alquileres.}
\cqa{¿Hasta cuándo se presenta el pliego de posiciones y qué pasa si no comparece el absolvente?}{Hasta media hora antes de la audiencia; si no comparece, estando notificado, queda confeso (\emph{ficta confessio}).}
\cqa{¿Qué hace el consultor técnico?}{Propone puntos de pericia, controla al perito oficial y puede presentar informe propio.}
\cresumen{Probar bien exige ofrecer toda la prueba al inicio, resguardar el soporte original, elegir el modo correcto de impugnar según el tipo de documento y preparar la confesional con tiempo.}
\end{cornell}

\chapter{La audiencia preliminar del artículo 360}\label{cap:audiencia}

Entre la etapa introductoria y la probatoria se ubica la audiencia preliminar. Es la primera y única oportunidad de inmediación del proceso y, bien aprovechada, puede resolver el conflicto o depurarlo. Este capítulo estudia su función polifuncional, el texto del artículo 360, su práctica real, la inapelabilidad de las resoluciones sobre prueba (art.~379), los dos momentos de la causa de puro derecho y la apelación en el sumarísimo.

\section{Utilidad y sentido estratégico}

La audiencia preliminar define qué hechos se prueban, qué prueba entra y si el caso puede cerrarse con acuerdo. Quien la domina evita juicios largos y costosos y no pierde prueba por errores de oportunidad. La fuente la presenta con la metáfora del \textbf{caballo de Troya}: algo sumamente estratégico dentro del proceso, un momento culminante en la fijación del conflicto, de los hechos y de los medios de prueba. Los conflictos se resuelven mejor cuando las personas hablan cara a cara frente a alguien que ordena la conversación, con hechos claros y opciones de acuerdo.

\begin{definition}{Inmediación en la audiencia preliminar}
Es la posibilidad, primera y única en el proceso, de que la jueza o el juez se encuentre con los abogados y las partes, presencial o virtualmente. Es la primera vez que hay un intercambio de ida y vuelta entre todos los intervinientes. Puede transformar la estructura del conflicto: es el momento en que el proceso comienza a revelar qué puede solucionarse. Si no, lo resolverá finalmente el juez.
\end{definition}

La estrategia responde al problema de los \textbf{plazos no razonables}: el proceso civil se excede siempre en el tiempo y, cuando la justicia llega, es tardía. Argentina fue condenada por la Corte Interamericana de Derechos Humanos por incumplir el plazo razonable, por lo que conviene usar las herramientas que creó el legislador para solucionar el conflicto o resolverlo de manera acotada y depurada. La jueza es una \textbf{directora del proceso}.

\section{Roles: jueza, abogados y partes}

Un error frecuente es silenciar al cliente (<<no abras la boca, hablo yo>>) o prescindir de su presencia mediante un poder. La parte es quien estuvo en el conflicto y está lesionada; el abogado es un tercero que la asiste. La jueza no puede trabajar bien sin la parte presente, sea virtual o presencial, porque puede reflotar otros temas y abrir el camino a la conciliación. La palabra es importante, siempre con respeto y apertura, y es más provechoso que las partes estén presentes en el despacho que solo sus abogadas.

\section{Función polifuncional del artículo 360}

\begin{formula}{Artículo 360 del CPCCN ---audiencia preliminar}
El juez citará a las partes a una audiencia, que presidirá con carácter \textbf{indelegable}; si no se hallare presente, no se realizará y se dejará constancia en el libro de asistencia. En ella: 1) invitará a las partes a una conciliación o a otra forma de solución de conflictos, y podrá derivarlas a mediación, con suspensión del procedimiento por treinta días desde la notificación del mediador; 2) recibirá las manifestaciones de las partes sobre el art.\ 361 y resolverá en el acto; 3) oídas las partes, fijará los hechos conducentes sobre los que versará la prueba; 4) recibirá la prueba confesional si fue ofrecida, y la ausencia de los absolventes no impide la audiencia; 5) proveerá las pruebas admisibles y concentrará en una sola audiencia la testimonial, con presencia del juez ---obligación delegable únicamente en el secretario o prosecretario letrado---; 6) si correspondiere, decidirá en el acto que la cuestión es de \textbf{puro derecho}, con lo que la causa queda concluida para definitiva.
\end{formula}

% TODO: verificar referencia presente en las fuentes originales. La fuente enumera los incisos de memoria durante la exposición e indica que conviene cotejar la redacción con el texto oficial vigente (Ley 17.454, con reformas de las Leyes 25.488 y 26.589).

La audiencia es \textbf{polifuncional}: no solo concilia ni solo oraliza. Se fijan los hechos controvertidos y conducentes y los medios probatorios admitidos; se desarrolla la confesional, se fija la testimonial, se sortean peritos, se escuchan las oposiciones y puede declararse la causa de puro derecho. Luego se dicta una resolución, visible en el expediente electrónico, que indica los hechos conducentes y controvertidos y los medios de prueba admitidos. Aunque las partes lleguen a un acuerdo en presencia de la jueza, esta debe dictar la resolución que lo \textbf{homologa}. Lo que se firma es un acta de comparecencia, o a veces ni siquiera ella: queda plasmado en la videograbación. El acuerdo debe estar descripto, porque es lo que se firmará.

\begin{example}{Un acuerdo alcanzado en la audiencia}
En una audiencia del art.\ 360 las partes pueden acordar el desalojo de forma pacífica, con costas por su orden, evitando el lanzamiento.
\end{example}

\subsection{La audiencia en la práctica}

\begin{table}[H]
\centering
\caption{La audiencia preliminar según la ley y en la práctica judicial}
\begin{tabularx}{\textwidth}{>{\raggedright\arraybackslash}p{0.17\textwidth} >{\raggedright\arraybackslash}X >{\raggedright\arraybackslash}X}
\toprule
\textbf{Aspecto} & \textbf{Lo que establece el Código} & \textbf{Lo que la fuente describe de la práctica} \\
\midrule
Quién la toma & La jueza o el juez, con carácter indelegable. & A veces un empleado o <<audiencista>> del juzgado. \\
Conciliación & Se invita a conciliar y se trabaja el conflicto. & Se pregunta <<¿hay conciliación?>> y, si no, se pasa a la etapa siguiente. \\
Hechos y prueba & Se fijan los hechos conducentes y se depuran los medios con las partes presentes. & Se deja todo para después: <<ya proveerán>>. \\
Duración & Momento muy activo del proceso. & Puede durar cinco minutos. \\
Sanción & El Código dispone el carácter indelegable. & No hay sanción para el magistrado y los abogados no suelen denunciar ante el Consejo de la Magistratura. \\
\bottomrule
\end{tabularx}
\end{table}

Según lo dicho en unas jornadas, ingresan ciento cuarenta mil causas por año. Sin audiencia del 360, la jueza igualmente debe fijar la prueba en esa instancia: no se acepta toda la prueba ofrecida. % TODO: contenido incompleto o ambiguo en las notas originales: queda abierto el cuestionamiento sobre cuántos empleados judiciales trabajan efectivamente.

\section{Prueba, impugnación y apelación: el artículo 379}

Las partes pueden impugnar \textbf{únicamente} cuando la jueza rechaza un \emph{hecho} articulado en la demanda, la contestación o la reconvención para que sea objeto de prueba. Esa impugnación se hace \textbf{in voce}, a viva voz en la audiencia. No pueden, en cambio, impugnar el rechazo de un \emph{medio de prueba} (por ejemplo, una pericial médica): ese medio solo se replantea ante la cámara, en la etapa del artículo 260, si se apela la sentencia definitiva y se concede el recurso con efecto libre.

\begin{formula}{Artículos 379 y 260 inc.\ 2 del CPCCN}
\textbf{Art.\ 379:} serán inapelables las resoluciones del juez sobre producción, denegación y sustanciación de las pruebas; y si se hubiere negado alguna medida probatoria, la parte interesada podrá solicitar a la cámara que la diligencie cuando el expediente le fuera remitido para conocer del recurso concedido contra la sentencia definitiva. \textbf{Art.\ 260 inc.\ 2:} dentro del quinto día de notificada la providencia y en un solo escrito, las partes deberán indicar las medidas probatorias denegadas en primera instancia que tengan interés en replantear en los términos de los artículos 379 y 385 in fine; la petición será fundada y resuelta sin sustanciación.
\end{formula}

% TODO: verificar referencia presente en las fuentes originales. El tramo final del art. 379 se completó conforme a la lectura realizada en clase, y el art. 260 se transcribe parcialmente; conviene cotejarlos con el texto oficial. La fuente también menciona <<artículo 320, algo así>> y lo descarta.

\begin{table}[H]
\centering
\caption{Qué se puede impugnar y qué no (art.\ 379)}
\begin{tabularx}{\textwidth}{>{\raggedright\arraybackslash}X >{\raggedright\arraybackslash}p{0.15\textwidth} >{\raggedright\arraybackslash}X}
\toprule
\textbf{Situación} & \textbf{¿Se impugna?} & \textbf{Qué hacer} \\
\midrule
La jueza rechaza un hecho articulado & Sí & Impugnar a viva voz en la audiencia. \\
La jueza no admite un medio de prueba & No (art.\ 379) & Replantear ante la cámara (art.\ 260) si se apela la sentencia y se concede libremente. \\
Proceso sumarísimo & Apelación muy limitada (art.\ 498) & No procede replantear la prueba por el art.\ 379. \\
\bottomrule
\end{tabularx}
\end{table}

\subsection{Proceso sumarísimo y recursos}

El desalojo tramita por proceso ordinario o sumarísimo, y la mayoría por sumarísimo. En este, la apelación se concede en relación, es mucho más limitada y la prueba denegada se pierde, sin posibilidad de replantearla por el artículo 379.

\begin{formula}{Artículo 498, inciso 6 ---proceso sumarísimo: apelación}
Solo serán apelables la sentencia definitiva y las providencias que decreten o denieguen medidas precautorias. La apelación se concederá en relación, en efecto devolutivo, salvo [\ldots].
\end{formula}

Quien en su primer desalojo pida el recurso <<libremente concedido>> y el replanteo del art.\ 379 en un sumarísimo <<está en el horno>>. La fuente señala que este tema suele preguntarse en el examen y que debe leerse con atención el artículo 498.

\section{Los dos momentos de la causa de puro derecho}

\begin{definition}{Causa de puro derecho}
Es aquella en la que no hay apertura a prueba y solo corresponde determinar qué derecho se aplica a la actora y a la demandada: la causa se declara de puro derecho y el expediente pasa a resolverse para definitiva.
\end{definition}

\begin{table}[H]
\centering
\caption{Dos momentos para declarar la causa de puro derecho}
\begin{tabularx}{\textwidth}{>{\raggedright\arraybackslash}X >{\raggedright\arraybackslash}p{0.17\textwidth} >{\raggedright\arraybackslash}X}
\toprule
\textbf{Momento} & \textbf{Norma} & \textbf{Qué ocurre} \\
\midrule
Trabada la litis y resueltas las excepciones previas & Art.\ 359 & Se decide que es de puro derecho antes de la audiencia; firme la providencia se llama autos para sentencia. \\
Dentro de la audiencia preliminar & Art.\ 360 inc.\ 6 & La jueza lo declara en el acto; la causa queda concluida para definitiva. \\
Efecto común & Ambos & Ya no se discuten medios de prueba sino el derecho aplicable. \\
\bottomrule
\end{tabularx}
\end{table}

El art.\ 359 dispone además que, si se hubiesen alegado hechos conducentes sin conformidad entre las partes, aunque estas no lo pidan, el juez recibirá la causa a prueba conforme al art.\ 360.

\begin{sintesis}
La audiencia del art.\ 360 es la única oportunidad de inmediación: el juez concilia, fija hechos y admite prueba. Solo puede impugnarse el rechazo de hechos, no el de medios de prueba (art.\ 379), y en el sumarísimo la apelación es muy limitada (art.\ 498).
\end{sintesis}

\begin{cornell}{Cuadro Cornell del capítulo~\ref{cap:audiencia}}
\cqa{¿Por qué es estratégica la audiencia preliminar?}{Es la primera y única inmediación del proceso: allí puede solucionarse el conflicto o resolverse de modo más acotado y depurado, para acercarse al plazo razonable.}
\cqa{¿Qué funciones cumple según el art.\ 360?}{Conciliar o derivar a mediación, fijar hechos controvertidos y conducentes, recibir la confesional, proveer y concentrar la testimonial, sortear peritos y declarar de puro derecho.}
\cqa{¿Puede tomarla un empleado?}{El Código la declara indelegable; en la práctica a veces la toma un audiencista, sin sanción para el magistrado.}
\cqa{¿Qué se puede impugnar en la audiencia?}{El rechazo de un hecho articulado, a viva voz; no el rechazo de un medio de prueba (art.\ 379).}
\cqa{¿Cómo se replantea una prueba denegada?}{Ante la cámara, por el art.\ 260 inc.\ 2, si se apela la sentencia definitiva y se concede libremente; en el sumarísimo no procede.}
\cqa{¿Cuáles son los dos momentos para declarar la causa de puro derecho?}{El art.\ 359 (antes de la audiencia) y el art.\ 360 inc.\ 6 (en la audiencia).}
\cqa{¿Hace falta homologar un acuerdo logrado en la audiencia?}{Sí: la jueza debe dictar la resolución que lo homologa.}
\cresumen{La audiencia del art.\ 360 concentra conciliación, fijación de hechos y depuración de prueba. Conocer los límites del art.\ 379 y del art.\ 498 evita perder prueba y recursos, sobre todo en el sumarísimo.}
\end{cornell}

\chapter{Costas, honorarios y tasa de justicia}\label{cap:costas}

El desalojo tiene un costo económico que el abogado debe explicar al cliente antes de iniciarlo. Este capítulo distingue costos y costas, describe sus componentes, el modo de pagar la tasa de justicia, el cálculo según el monto del proceso y la práctica del convenio de honorarios.

\section{Principio general y composición de las costas}

El principio general es que quien reclama judicialmente algo sin que el demandado haya incurrido en mora, ni se haya allanado a la primera intimación, no puede cobrarle las costas, porque este no incumplió. Si el demandado fue intimado y no cumplió, y después es demandado, carga con ellas. Los \textbf{costos} son lo que genera de gasto; las \textbf{costas} comprenden honorarios y todo lo que lleva el proceso.

\begin{important}
\textbf{Componentes de las costas:}
\begin{itemize}
\item \textbf{Honorarios profesionales}: según una escala porcentual \emph{regresiva} sobre el monto del proceso: a mayor monto, menor porcentaje.
\item \textbf{Tasa de justicia}: aproximadamente el tres por ciento del monto del juicio, que se abona para iniciar la causa y queda en beneficio del Poder Judicial.
\item \textbf{Gastos causídicos}: desembolsos concretos que cada parte debió hacer para el proceso ---por ejemplo, contratar a un ingeniero para constatar un desperfecto---, acreditados con la factura.
\end{itemize}
\end{important}

Quien pierde el juicio paga los gastos causídicos que la otra parte afrontó. Los honorarios regulados se ejecutan por los procedimientos de ejecución de sentencia, incluso contra el abogado de la parte vencida en lo que hace a sus propios honorarios.

\section{Tasa de justicia}

\begin{formula}{Cálculo de la tasa de justicia}
En los procesos de \textbf{monto determinado}, se calcula como un porcentaje de la suma reclamada (en el caso trabajado, varios meses de canon locativo). En los de \textbf{monto indeterminado} ---por ejemplo, una acción declarativa de certeza sobre un límite entre inmuebles--- corresponde un monto fijo, actualizado periódicamente por acordada, salvo que se invoque el beneficio de litigar sin gastos.
\end{formula}

El pago admite dos vías: el formulario en papel, que se abona en una sucursal bancaria habilitada, y el pago electrónico por banca judicial, para el cual el cliente necesita una tarjeta de débito propia o de un familiar. Muchos clientes prefieren el trámite presencial por desconfianza a ingresar los datos de la tarjeta; por eso conviene ofrecer siempre ambas opciones desde el principio, anticipando los inconvenientes de cada una. Si el cliente tiene la posibilidad de recuperar el inmueble y alquilarlo o venderlo, los honorarios regulados judicialmente suelen ser bajos en relación con el trabajo y el desgaste, por lo que es habitual un convenio de honorarios con el cliente, de modo comparable a la comisión de una inmobiliaria. El petitorio de la demanda incluye el pedido de desocupación inmediata, la condena en costas y, cuando corresponda, la reserva del caso federal.

\section{La presentación electrónica}

El inicio de una causa ante el Poder Judicial de la Nación comienza con el envío de un formulario de sorteo a la Cámara del fuero, que asigna aleatoriamente el juzgado y el número de expediente. Este puede conocerse por la respuesta automática enviada a la casilla desde la que se remitió el formulario, o consultando el apartado de radicaciones del sistema, donde figura el listado de las últimas causas sorteadas.

\begin{note}
Sobre la excepción de arraigo, las fuentes difieren. Una la describe como excepción en desuso, salvo domicilio o bienes fuera del país; otra afirma que la reforma del CCCN la derogó, y que consistía en una caución a quien demandaba sin domicilio ni bienes en el país para garantizar las costas, pensada para quien litiga desde el exterior sin vínculo con el país y no como regla contra los extranjeros.
\end{note}
% TODO: verificar consistencia entre fuentes originales. La clase del 07/09 trata el arraigo como presupuesto <<hoy casi en desuso>>; la del 24/09 afirma que fue derogado por el CCCN.

\begin{cornell}{Cuadro Cornell del capítulo~\ref{cap:costas}}
\cqa{¿Qué diferencia hay entre costos y costas?}{Los costos son lo que genera de gasto; las costas comprenden honorarios y todo lo que lleva el proceso (tasa, bono, etc.).}
\cqa{¿Cómo se componen las costas?}{Honorarios (escala regresiva), tasa de justicia (alrededor del tres por ciento) y gastos causídicos acreditados con factura.}
\cqa{¿Cuándo no se cobran las costas al demandado?}{Si no incurrió en mora ni se allanó a la primera intimación.}
\cqa{¿Cómo se calcula la tasa de justicia?}{En monto determinado, como porcentaje de lo reclamado; en monto indeterminado, un monto fijo actualizado por acordada, salvo beneficio de litigar sin gastos.}
\cqa{¿Qué vías hay para pagarla?}{Formulario en papel en sucursal bancaria o pago electrónico por banca judicial; conviene ofrecer ambas.}
\cresumen{Antes de litigar hay que explicar al cliente qué costas puede generar y a cargo de quién, ofrecer ambas vías de pago de la tasa y, si corresponde, pactar honorarios por convenio.}
\end{cornell}

\chapter{El lanzamiento}\label{cap:lanzamiento}

La sentencia firme que ordena la restitución se hace efectiva mediante el lanzamiento: la desocupación forzada del inmueble. Es la etapa ejecutoria y la más expuesta a imprevistos. Este capítulo reúne las dificultades prácticas, los oficios que conviene solicitar y la intervención de los auxiliares descriptos en el capítulo~\ref{cap:organos}.

\section{La diligencia en la práctica}

Al lanzar, el oficial de justicia debe constatar el domicilio, habitualmente identificado con el número de la puerta. Si el cartel con el número fue retirado poco antes de la diligencia ---estrategia habitual para dilatar alegando que <<no hay domicilio>>---, pueden usarse medios prácticos para volver a señalizar el número y permitir la diligencia.

\begin{important}
El \textbf{flete} es obligatorio en la diligencia: lo impone el propio proceso para evitar que las pertenencias de quien es desalojado terminen en la calle. Es obligación del locador, que recupera el inmueble, hacerse cargo de al menos un mes de depósito de esas pertenencias; transcurrido ese plazo sin que sean retiradas, pueden donarse.
\end{important}

\begin{example}{Oficios de ambulancia y veterinario}
Al solicitar el oficio para el lanzamiento es habitual pedir también uno para ambulancia y otro para atención veterinaria, porque si hay animales no es posible dejarlos en la calle, y es frecuente que alguien se descompense durante la diligencia, por el esfuerzo físico o, a veces, simulando una afección grave.
\end{example}

\begin{example}{Lanzamiento con intervención policial armada}
En casos difíciles, con repercusión en medios, puede ser necesaria la intervención de personal policial armado. El abogado ingresa junto con el oficial de justicia y la fuerza pública, mientras el cliente permanece afuera, a diferencia de la policía, sin armas ni chaleco, en una situación de tensión considerable.
\end{example}

Frente a estas dificultades siempre conviene intentar una salida pacífica, aunque a veces la contraparte no tenga ningún interés en negociar (véase el capítulo~\ref{cap:masc}). Si hay personas vulnerables en el inmueble, el lanzamiento puede exigir intervención del Ministerio Público de la Defensa y, si el inmueble está en otra jurisdicción, el oficio de la Ley 22.172 (capítulo~\ref{cap:organos}). La estrategia recomendada es solicitar la totalidad de los auxiliares disponibles.

\begin{cornell}{Cuadro Cornell del capítulo~\ref{cap:lanzamiento}}
\cqa{¿Por qué es obligatorio el flete?}{Para evitar que las pertenencias del desalojado terminen en la calle; el locador debe afrontar al menos un mes de depósito, tras el cual pueden donarse.}
\cqa{¿Qué oficios conviene pedir?}{Ambulancia y atención veterinaria, además de cerrajero y fuerza pública según el caso.}
\cqa{¿Qué hacer si retiran el número de la puerta?}{Recurrir a medios prácticos para volver a señalizarlo y permitir que la diligencia se lleve a cabo.}
\cqa{¿Qué hay que hacer si hay menores en el inmueble?}{Dar intervención al Ministerio Público de la Defensa y, si el inmueble está en otra jurisdicción, tramitar el oficio de la Ley 22.172.}
\cresumen{El lanzamiento exige anticipar imprevistos: solicitar todos los auxiliares, prever el depósito de pertenencias y no descartar una salida pacífica incluso en el último momento.}
\end{cornell}

\part{Tutela anticipada, vulnerabilidad y acceso a la justicia}

\chapter{Medidas cautelares, contracautela y caución}\label{cap:cautelares}

El tiempo que insume el proceso puede afectar la efectividad de la sentencia. Este capítulo estudia la respuesta del sistema a ese riesgo: las medidas cautelares, sus presupuestos y su caducidad, la diferencia entre cautela y condena, la medida innovativa del artículo 230 con el precedente <<Camacho Acosta>>, y las modalidades de caución y contracautela. El capítulo~\ref{cap:684bis} aplica estas ideas al desalojo inmediato.

\section{La medida cautelar y sus presupuestos}

Los procesos judiciales insumen tiempo, y ese tiempo puede jugar en contra del titular del derecho. Como vía expeditiva puede mencionarse la acción de amparo; en general, el mecanismo para proteger un derecho amenazado mientras dura el proceso es solicitar una medida cautelar.

\begin{definition}{Medida cautelar}
Las medidas cautelares son \textbf{provisionales} y \textbf{accesorias} de un proceso principal: no constituyen un proceso en sí mismas, y su función es proteger el resultado eventual de aquel mientras se sustancia (por ejemplo, un embargo o un secuestro de mercadería). Exigen como presupuestos \textbf{verosimilitud del derecho} y \textbf{peligro en la demora}. Como se resuelven sin escuchar a la contraparte (\emph{inaudita parte}), exigen como contrapartida una \textbf{contracautela}, en resguardo de los daños que la medida pudiera causar si finalmente no asistiera razón al peticionante.
\end{definition}

Las medidas cautelares constituyen un subsistema autónomo del sistema procesal y requieren siempre reunir sus propios presupuestos con independencia del proceso principal (véase la sección sobre el fiador en el capítulo~\ref{cap:terceros}).

\subsection{Caducidad de las cautelares y su interacción con la mediación}

\begin{normacitada}
\textbf{Caducidad de las medidas cautelares.} Por su carácter accesorio, cuando una medida cautelar se traba antes de la demanda, quien la obtuvo debe iniciar el proceso principal dentro de los \textbf{diez días} contados desde que quedó trabada. Vencido el plazo, la medida caduca de pleno derecho y quien la solicitó puede quedar obligado a responder por los daños causados por la traba injustificada.
\end{normacitada}

Si para demandar es necesario transitar antes la mediación obligatoria, y esta se inicia dentro de los diez días, el plazo se detiene mientras dure la mediación y vuelve a correr recién a los veinte días del acta de cierre, igual que la prescripción (véase el capítulo~\ref{cap:masc}).

\begin{important}
La estrategia óptima es iniciar la mediación en los \textbf{primeros días} del plazo de diez, no en los últimos: si se inicia el día nueve, queda apenas un día hábil tras el acta de cierre para redactar e interponer la demanda; si se inicia el primer día, quedan prácticamente los diez días completos.
\end{important}

\begin{example}{El boleto de compraventa y la multa diaria}
Un boleto de compraventa incluía una cláusula penal: una multa por cada día de demora en escriturar. Si el titular de la cautelar la deja caducar por no iniciar a tiempo la demanda, deberá responder por esa demora, porque dejar caducar la medida hace presumir que no existía razón suficiente para haberla trabado. Si la demanda se inició en tiempo y forma, no hay responsabilidad por la demora.
\end{example}

\begin{note}
Debe notificarse al mediador el resultado del proceso, porque sus honorarios se calculan sobre el monto reclamado y los paga la parte requerida si hay acuerdo, o quien resulte condenado en costas si el conflicto termina en juicio.
\end{note}

\section{Cautela y condena}

La resolución de \textbf{condena} es definitiva, aunque pueda tener luego una etapa recursiva. La \textbf{cautela} es provisoria: puede quedar firme, modificarse o revocarse. Esta distinción es relevante porque, en ciertos casos, el juez resuelve en el plano cautelar lo mismo que resolvería al finalizar el proceso: en el desalojo, con una cautela puede ordenar la desocupación inmediata, y la condena final ordena precisamente la desocupación. Es la misma consecuencia práctica en dos planos distintos, uno provisorio y otro definitivo.

\begin{important}
Que el juez adelante una decisión cautelar no implica prejuzgamiento, porque puede modificar su criterio al dictar la sentencia definitiva.
\end{important}

\section{La medida innovativa del artículo 230 y «Camacho Acosta»}

\begin{formula}{Artículo 230 del CPCCN ---medida cautelar innovativa}
Procede cuando se acreditan los presupuestos generales de toda medida cautelar ---verosimilitud del derecho y peligro en la demora--- y requiere contracautela o caución. La doctrina de Calamandrei la identifica como el sistema matriz de todas las cautelares: ante cualquier situación de peligro que no encuadre en una figura específica (embargo, inhibición general de bienes, prohibición de contratar), puede recurrirse a esta vía general.
\end{formula}

La medida admite dos variantes: la de \textbf{no innovar}, que pide al juez mantener la situación tal como está para que no se modifique en perjuicio de quien la solicita, y la \textbf{innovativa} propiamente dicha, que pide modificar la situación actual porque está causando un daño. La distinción es aplicable a todo proceso. Antes de que existiera el artículo 684 bis, este sistema general ya existía, y algunos profesionales recurrían a él para pedir un lanzamiento anticipado, aunque no era una práctica cotidiana.

\begin{example}{El leading case «Camacho Acosta» (Fallos 320:1633)}
Un trabajador pierde un miembro superior en un accidente de trabajo. Mientras avanza el proceso, el muñón se va endureciendo y necesita con urgencia una prótesis, cara y de difícil consecución, sin poder esperar la sentencia de condena. La Corte ordenó, por una medida cautelar innovativa, que el empleador entregara la prótesis antes de que terminara el proceso. El caso marcó que, ante situaciones de esta naturaleza, los juzgados deben proteger primero y conocer después: proteger a la parte que atraviesa la urgencia y analizar luego si existía derecho.
\end{example}

\section{Caución y contracautela en la práctica}

La forma de fijar las cauciones no está reglada por la ley, lo que genera criterios diversos según cada juzgado.

\begin{table}[H]
\centering
\caption{Modalidades de caución en la práctica profesional}
\begin{tabularx}{\textwidth}{>{\raggedright\arraybackslash}p{0.25\textwidth} X}
\toprule
\textbf{Modalidad} & \textbf{Descripción} \\
\midrule
Autoembargo & El propio locador ofrece su inmueble en garantía; al terminar el proceso se levanta el embargo. Permite constituir la caución sin disponer de dinero en efectivo. Si la sentencia es desfavorable, los daños se pagan en dinero y, si no se abonan, el inmueble se ejecuta. \\
Caución real sobre otro bien & Si el locador no quiere comprometer el inmueble en disputa, ofrece otro bien de su propiedad (por ejemplo, un terreno). \\
Depósito dinerario judicial & Se deposita una suma en el banco de los tribunales. El monto varía según el criterio de cada juzgado: algunos piden seis meses de alquiler, otros tres, y otros hasta casi dos años, lo que suele generar resistencia del cliente. \\
\bottomrule
\end{tabularx}
\end{table}

\begin{important}
La caución no tiene un monto ni una forma reglados por la ley: depende del criterio de cada juzgado. Las alternativas más usadas son el autoembargo del propio inmueble en disputa, la caución real sobre otro bien y el depósito dinerario judicial.
\end{important}

\begin{sintesis}
Las cautelares protegen provisoriamente el resultado del proceso con verosimilitud del derecho, peligro en la demora y contracautela, y caducan a los diez días si no se demanda. El artículo 230 es el sistema matriz; la caución varía según cada juzgado.
\end{sintesis}

\begin{cornell}{Cuadro Cornell del capítulo~\ref{cap:cautelares}}
\cqa{¿Qué presupuestos exige una cautelar y por qué exige contracautela?}{Verosimilitud del derecho y peligro en la demora; como se dicta sin oír a la contraparte, la contracautela resguarda los daños si el peticionante no tenía razón.}
\cqa{¿En qué se diferencia cautela de condena?}{La cautela es provisoria y revocable; la condena es definitiva. En el desalojo ambas pueden ordenar la desocupación, sin que adelantar la cautela implique prejuzgamiento.}
\cqa{¿Cuándo caduca una cautelar trabada antes de la demanda?}{A los diez días de trabada si no se inicia el proceso principal; con mediación, el plazo se detiene y vuelve a correr a los veinte días del acta de cierre.}
\cqa{¿Qué es el artículo 230 y qué enseña «Camacho Acosta»?}{El sistema matriz de las cautelares (con variantes de no innovar e innovativa); el caso enseña a proteger primero y conocer después.}
\cqa{¿Qué modalidades de caución existen?}{Autoembargo, caución real sobre otro bien y depósito dinerario judicial; la ley no fija monto ni forma.}
\cresumen{La cautela adelanta una protección provisoria a cambio de una garantía. Hay que conocer su plazo de caducidad, elegir bien la caución y entender que el artículo 230 sirve de marco general para figuras especiales como la del artículo 684 bis.}
\end{cornell}

\chapter{La desocupación inmediata: artículo 684 bis}\label{cap:684bis}

El artículo 684 bis permite adelantar el lanzamiento sin esperar la sentencia. Este capítulo explica su sistema, sus tres causales taxativas, el sentido técnico de <<intruso>>, el modo de acreditar el abandono, un caso práctico de falta de pago aparente y el régimen de impugnación con su jurisprudencia.

\section{El sistema del artículo 684 bis}

\begin{formula}{Artículo 684 bis del CPCCN ---desocupación inmediata}
Incorporado al Código Procesal de 1967 para los procesos de desalojo, habilita al juez a disponer la desocupación inmediata del inmueble en forma anticipada a la sentencia definitiva, una vez vencido el plazo para contestar la demanda y cumplidas las cargas de citación y emplazamiento, previa fijación de una caución.
\end{formula}

Su función consiste en adelantarse a la condena y llegar antes con un lanzamiento. Por el principio de defensa en juicio siempre existe traslado a la otra parte. Vencido el plazo para contestar y cumplidas las cargas de citación y emplazamiento, el juez decide si exige caución para dar lugar a la medida. En términos generales se exige: si quien obtuvo el uso y goce anticipadamente resulta, en el fallo definitivo, haber pedido el lanzamiento de una persona que pagaba, estaba al día y no cambió el destino, deberá responder con la caución por todos los daños y perjuicios.

\begin{important}
A diferencia de una cautelar del artículo 230, el 684 bis exige que la otra parte ya haya sido notificada de la demanda y que haya vencido su plazo para contestar: no se dicta <<a espaldas>> del demandado, como sí puede ocurrir con un embargo cautelar tradicional. La fuente también lo describe como una figura que, sin ser técnicamente una cautelar del artículo 230, comparte con ella la exigencia de contracautela. El 230 es el sistema general; el 684 bis, un subsistema especial incorporado para el desalojo.
\end{important}

\section{Causales de procedencia}

\begin{formula}{Causales del artículo 684 bis}
1) Falta de pago del canon locativo. 2) Vencimiento del plazo contractual. 3) Intrusión: ocupación del inmueble sin consentimiento, bajo ninguna forma, de quien tiene derecho a otorgar su uso y goce.
\end{formula}

No toda persona que ocupa un inmueble sin contrato vigente es un intruso: si existe un contrato, por ejemplo, del año 2000, se trata de una locación continuada a la que nunca se renovó el contrato, y no de una intrusión. El intruso es quien ocupa sin consentimiento ni derecho alguno.

Cuando se invoca el \textbf{abandono} del inmueble, debe comprobarse: corresponde solicitar al juez un mandamiento de constatación, con oficial de justicia de la zona ---y, eventualmente, cerrajero---, para corroborar que en el interior no hay nada, o que hay bienes en mal estado. No corresponde autorizar al cliente a ingresar por su cuenta: el contrato puede seguir vigente y la persona puede regresar.

\begin{example}{El cambio de cuenta bancaria y la falta de pago}
Un locador de edad avanzada, por desconfianza hacia su banco, abrió una cuenta en otra entidad y avisó informalmente a su inquilina que depositara allí los alquileres, sin formalizarlo. La inquilina pagó siempre en la nueva cuenta y enviaba capturas de los comprobantes por WhatsApp, que el locador, según su relato, no revisaba. El locador inició el desalojo por falta de pago sin saber que la inquilina pagaba. Esta, ya asesorada, aportó el extracto bancario y las capturas; se contactó al abogado contrario para esclarecer la situación, y el proceso, iniciado sin causa real, se desactivó sin ejecutar la caución. La enseñanza es metodológica: hay que dejar que el cliente cuente su versión sin dirigir la entrevista.
\end{example}

\begin{formula}{Recomendación de redacción contractual}
Todo cambio del domicilio de pago o de la cuenta bancaria designada en el contrato debe formalizarse por escrito. Se recomienda incorporar una cláusula de comunicación periódica sobre el estado del inmueble y del pago.
\end{formula}

\section{Impugnación y jurisprudencia}

La resolución que dispone la desocupación inmediata es impugnable. En el \textbf{efecto devolutivo}, la orden se cumple igual aunque el recurso esté pendiente ante la Cámara; en el \textbf{efecto suspensivo}, el juez suspende la ejecución hasta que el superior resuelva. En los precedentes citados el recurso se concede con efecto devolutivo, y el suspensivo es la excepción.

\begin{example}{Jurisprudencia citada sobre el artículo 684 bis}
\textbf{Sala H de la Cámara Nacional Civil:} concedió el recurso con efecto devolutivo, de modo que la resolución impugnada se cumple igual. \textbf{Sala D de la Cámara Nacional Civil:} confirmó la procedencia del 684 bis en un caso en el que ya se había concretado el lanzamiento con auxilio de la fuerza pública, continuando adelante pese a la oposición. \textbf{Departamento judicial de La Plata:} un precedente provincial que también reconoce la figura.
\end{example}

\begin{cornell}{Cuadro Cornell del capítulo~\ref{cap:684bis}}
\cqa{¿Qué permite el artículo 684 bis y cuándo?}{Disponer la desocupación inmediata antes de la sentencia, vencido el plazo para contestar y cumplidas citación y emplazamiento, previa caución.}
\cqa{¿En qué se diferencia de una cautelar tradicional?}{Exige que el demandado ya esté notificado y haya vencido su plazo: no se dicta a sus espaldas.}
\cqa{¿Cuáles son sus causales?}{Falta de pago, vencimiento del plazo contractual e intrusión; son taxativas.}
\cqa{¿Quién es intruso?}{Quien ocupa sin consentimiento alguno del titular; un contrato vencido o una locación continuada no lo convierten en intruso.}
\cqa{¿Cómo se acredita el abandono?}{Con mandamiento de constatación ordenado por el juez, no con el ingreso del cliente por su cuenta.}
\cqa{¿Qué efecto suele tener el recurso contra la desocupación inmediata?}{Devolutivo: el lanzamiento se ejecuta mientras la Cámara resuelve; el suspensivo es la excepción.}
\cresumen{El 684 bis adelanta el lanzamiento con caución y solo procede por tres causales taxativas. Escuchar al cliente y verificar los hechos evita desalojos infundados, como el del cambio de cuenta bancaria.}
\end{cornell}

\chapter{Vulnerabilidad, género y acceso a la justicia en el desalojo}\label{cap:acceso}

La perspectiva de vulnerabilidad presentada en el capítulo~\ref{cap:principios} se manifiesta en el proceso de maneras concretas: en el lanzamiento, en la jurisprudencia de la Corte, en el lenguaje de las notificaciones y en la tensión entre las formas y la tutela judicial efectiva. Este capítulo reúne esas manifestaciones y cierra con la idea de que el desalojo exige priorizar la solución del conflicto.

\section{La vulnerabilidad en ambos lados de la puerta}

La vulnerabilidad puede presentarse de ambos lados de la relación locativa: el locador también puede ser la parte vulnerable, no solo el locatario. Hay personas sin familiares, o con familiares sin recursos, que viven del alquiler de un departamento o local pequeño junto con una jubilación mínima. No todos los locadores son vulnerables ni todos los locatarios lo son: algunos pagan cánones muy altos, incluso en dólares, por lo que deben analizarse las circunstancias de cada caso. En el desalojo suele mirarse el lado interno de la puerta y no el externo.

\begin{example}{El lanzamiento demorado por la presencia de menores}
En un proceso con sentencia firme, ya en la etapa de lanzamiento, se presentó una familia con menores a cargo alegando ocupar el inmueble, sin poder acreditarlo con facilidad. El juzgado exigió la partida de nacimiento que probara el vínculo y, acreditado este, ofició a organismos estatales ---nacionales y municipales--- para determinar qué ayuda o reubicación podía ofrecerse, condicionando el lanzamiento a esa respuesta. Los organismos oficiados se desentendieron en la práctica, cada área remitiendo la competencia a otra. En el mismo expediente se detectaron indicios de que el ocupante original había cedido el inmueble a esa familia deliberadamente para dilatar el lanzamiento, lo que fue calificado como posible abuso del proceso, pasible de sanción.
\end{example}

\begin{important}
Los juzgados pueden condicionar la ejecución de un lanzamiento a la respuesta de organismos estatales cuando hay menores u otras personas vulnerables. El uso de la ocupación con fines dilatorios, para retrasar artificialmente un lanzamiento con sentencia firme, puede constituir un abuso del proceso sancionable.
\end{important}

En la Ciudad de Buenos Aires existe un mecanismo por el cual, a través de los Centros de Gestión y Participación (CGP) de las comunas, la persona interesada lleva su contrato, acredita que tiene trabajo y el Gobierno de la Ciudad puede otorgar una subvención para evitar que la familia quede en la calle o buscar una reubicación. Requiere estar radicado en la Ciudad; en otros distritos, como La Matanza, esa herramienta no está tan desarrollada.

\subsection{Perspectiva de género y de familia}

Existen precedentes de perspectiva de género y de familia en materia de desalojo, entre ellos un fallo de la provincia de Corrientes que rechazó un desalojo entre familiares ---dirigido contra una cuñada--- por entender que, detrás de un planteo patrimonial, había un conflicto propio del fuero de familia. Cuando el conflicto entre parientes tiene una raíz familiar y no patrimonial, corresponde encauzarlo por el fuero de familia.

\section{La jurisprudencia de la Corte: «para habitar hay que habitar con derechos»}

\begin{formula}{Cómo citar un fallo de la Corte Suprema}
Fallos: 335:452. Se ingresa a la página de la Corte, a la sección de sumarios y se localiza por tomo (335) y página (452) de la colección oficial de Fallos. No es necesario memorizar el nombre completo de las partes para citar correctamente un fallo.
\end{formula}

Este fallo se analiza desde una perspectiva de vulnerabilidad. Su enseñanza central es que para habitar hay que habitar con derechos, vinculando la locación de vivienda con la perspectiva de género. Es una mirada relativamente reciente que ha ido creciendo progresivamente.

\begin{important}
En el caso, las críticas se dirigieron a los jueces de instancias inferiores por no haber respetado la igualdad y no discriminación. La Corte sostuvo que la igualdad es un principio constitucional y una garantía, pero que en el caso concreto no se encontraba violentada, y que hubo un error en la forma de plantear el recurso extraordinario, lo que dejó la cuestión de fondo fuera de análisis. Un error de forma ---por ejemplo, respecto de los recaudos del artículo 280 del CPCCN--- puede impedir que la Corte llegue siquiera a analizar lo sustancial (véase el capítulo~\ref{cap:principios}).
\end{important}

\section{Lenguaje claro y derecho a la información procesal}

\begin{example}{La cédula que no se entendió a tiempo}
Una señora recibió la cédula de notificación de una demanda de desalojo, pero no comprendió la urgencia del documento y no lo envió por foto a su hijo, que estaba de viaje con su nieto en un viaje de egresados de séptimo grado. No quiso interrumpir el viaje. El domicilio estaba correctamente notificado ---el hijo, soltero, figuraba con su documento allí---, por lo que la notificación era formalmente válida; sin embargo, la falta de comprensión por parte de una persona sin formación jurídica generó la pérdida de un plazo procesal relevante.
\end{example}

Quien no tiene formación jurídica con frecuencia no advierte que en una notificación ---percibida como un documento hostil, en lenguaje técnico--- hay una fecha fijada por el oficial notificador. En algunos casos la redacción compleja de los autos resulta a propósito, práctica incorrecta, porque toda persona tiene derecho ---garantizado por la Constitución Nacional, la Convención Americana sobre Derechos Humanos y el bloque de constitucionalidad--- a recibir información clara sobre los caminos procesales que puede tomar.

\begin{important}
El derecho a la información procesal, en un lenguaje claro y comprensible, es un derecho constitucional y convencional. Una notificación formalmente válida puede fracasar en su función si el destinatario no comprende su contenido ni su urgencia.
\end{important}

\section{Flexibilización de las formas y tutela judicial efectiva}

Existen fallos provinciales, en general más innovadores, que permitieron contestar la demanda fuera del plazo. En un caso, una jueza de familia dejó contestar fuera de término, dejando en claro esa circunstancia y aclarando que no adelantaba su resolución: solo buscaba escuchar por qué no se contestó dentro de los cinco días del sumarísimo y qué prueba se ofrecía. Esta flexibilización, orientada a la verdad jurídica material, entra en tensión con la preclusión y puede abrir la puerta a la inseguridad jurídica. Las circunstancias de tiempo y lugar importan ---no es lo mismo litigar en una gran ciudad que en un pueblo alejado--- y el fuero de familia admite, en general, flexibilidades que no existen en el fuero patrimonial.

Como contrapunto, las fuentes señalan que las propuestas de mayor oralidad e inmediación no implican por sí solas una mejora si no se acompañan de recursos humanos y materiales: de nada sirve declarar el proceso oral si la primera audiencia demora años. La Corte Interamericana de Derechos Humanos ha condenado a numerosos países, entre ellos la Argentina, por incumplir el plazo razonable. En el sistema argentino rige una preclusión estricta y las costas y cauciones suelen ser bajas; en el \emph{pretrial} estadounidense (con \emph{discovery}) y en el \emph{disclosure} británico, en cambio, los hechos y la prueba se exhiben entre las partes antes del juicio y quien demanda sin fundamento corre el riesgo de costas elevadas. Preferible a multiplicar juzgados es reforzar los existentes con personal efectivamente disponible.

\section{La solución del conflicto por sobre la resolución impuesta}

La balanza de la justicia sintetiza un paralelo constante: buscar soluciones antes que resoluciones impuestas, y agilizar los procesos en aras de la tutela judicial efectiva, por ejemplo para el jubilado que vive de un ingreso mínimo y necesita recuperar rápido su pequeño inmueble para volver a alquilarlo. La abogacía es también un servicio de mediación entre las partes: mediar, conciliar y autocomponer son caminos preferibles, salvo la violencia. Las fuentes agregan que el derecho y el arte ---con ejemplos como <<Las dos Fridas>>, que muestra a la misma persona representada de dos formas, como demanda y contestación narran los mismos hechos desde dos perspectivas--- sirven como herramienta pedagógica para pensar que una misma realidad admite más de una representación legítima.

\begin{sintesis}
La vulnerabilidad se evalúa en ambos lados del conflicto y puede condicionar el lanzamiento; la forma procesal convive con el derecho a comprender la notificación; la flexibilización de plazos debe equilibrarse con la seguridad jurídica; y la prioridad es resolver el conflicto.
\end{sintesis}

\begin{cornell}{Cuadro Cornell del capítulo~\ref{cap:acceso}}
\cqa{¿Quién puede ser vulnerable en el desalojo?}{Cualquiera de las partes: el locatario y también el locador que vive del alquiler con ingresos mínimos; se analiza cada caso.}
\cqa{¿Qué puede ocurrir en el lanzamiento si hay menores?}{El juzgado puede exigir acreditar el vínculo y oficiar a organismos estatales, condicionando el lanzamiento a su respuesta; la ocupación cedida para dilatar puede ser abuso del proceso.}
\cqa{¿Qué enseña Fallos 335:452?}{Que para habitar hay que habitar con derechos; la Corte no entró al fondo por un error en la forma de plantear el recurso.}
\cqa{¿Por qué importa el lenguaje claro?}{Porque una notificación válida puede fracasar si el destinatario no comprende su urgencia; la información clara es un derecho constitucional y convencional.}
\cqa{¿Qué tensión plantea flexibilizar plazos?}{Enfrenta la verdad jurídica material con la seguridad jurídica de la preclusión; el fuero de familia admite más flexibilidad que el patrimonial.}
\cresumen{El desalojo exige mirar a ambos lados de la puerta, cuidar la comprensión de las notificaciones y balancear formas y tutela efectiva, priorizando la solución del conflicto sobre la resolución impuesta.}
\end{cornell}

\part{Práctica profesional integrada}

\chapter{Práctica profesional: expediente electrónico y planteos procesales}\label{cap:practica}

El conocimiento de las normas no alcanza para llevar un desalojo adelante: el trámite depende también de sistemas electrónicos, de formatos de archivo y de decisiones sobre prueba y planteos constitucionales. Este capítulo reúne esas herramientas de práctica, que complementan el estudio de la demanda (capítulo~\ref{cap:demanda}), de la prueba (capítulo~\ref{cap:prueba}) y de las costas (capítulo~\ref{cap:costas}).

\section{Más allá del artículo 330: el expediente electrónico}

El artículo 330 del CPCCN fija los requisitos formales mínimos de la demanda, pero no constituye una plantilla suficiente para casos complejos: un caso con varias causales y varios sujetos exige apartados que ese artículo no prevé expresamente, y conviene desarrollarlos según la complejidad de la demanda.

Al presentar la demanda con la prueba documental, el sistema de gestión electrónica exige que cada archivo adjunto no supere un tamaño máximo. Cuando la documentación excede ese peso, se divide en partes (<<documental uno>>, <<documental dos>>) y cada parte se adjunta por separado. El límite se fue actualizando con las acordadas de la Corte Suprema de Justicia de la Nación dictadas desde 2020: comenzó en 5 MB y hoy se encuentra en 20 MB.

\begin{table}[H]
\centering
\caption{Fuentes normativas del expediente electrónico}
\begin{tabularx}{\textwidth}{>{\raggedright\arraybackslash}p{0.27\textwidth} >{\raggedright\arraybackslash}p{0.27\textwidth} X}
\toprule
\textbf{Fuente normativa} & \textbf{Qué regula} & \textbf{Vigencia} \\
\midrule
Acordadas de la Corte Suprema (desde 2020) & Gestión electrónica del expediente. & Se actualizan periódicamente; el límite de archivo vigente es de 20 MB. \\
Código Procesal Civil y Comercial de la Nación & Formas y plazos de los escritos (sistema originalmente analógico). & Vigente; se aplica junto con las acordadas electrónicas. \\
Reglamento de la Justicia Nacional & Formas complementarias de los escritos. & De origen anterior, pero aún vigente y de aplicación práctica. \\
\bottomrule
\end{tabularx}
\end{table}

La transición del sistema analógico al electrónico trajo un problema práctico: la demora de los juzgados en aceptar y despachar los escritos electrónicos, que quedan en espera de que el personal los reciba desde su propio sistema (la <<dependencia>>). Hay juzgados que despachan en plazos prudenciales y otros que acumulan demoras de hasta un mes; ante un escrito urgente, una práctica habitual es concurrir a la mesa de entradas y solicitar el despacho directamente.

\begin{important}
Cumplir el artículo 330 y las acordadas de gestión electrónica es condición necesaria, pero no suficiente: conocer el límite de tamaño de archivo, saber dividir la documental y tener presente que el despacho depende de que el personal del juzgado acepte el escrito son datos de práctica que no figuran en ningún código, pero determinan si un escrito llega a horario.
\end{important}

\section{Fuentes de prueba y medios de prueba}

La demanda es un acto procesal de postulación, distinto de los actos de notificación ---a cargo del oficial notificador--- y de los actos de obtención, mediante los cuales se solicita un informe. Antes de redactarla existe la etapa de diligencias preliminares del artículo 323 (véase el capítulo~\ref{cap:intimacion}). Como premisa de la práctica, la parte demandada suele negar y desconocer todos los hechos y la totalidad de la documental, incluso la certificada por escribano, para dilatar el proceso, por ejemplo pidiendo una pericial caligráfica improcedente cuando la firma ya está certificada por un funcionario público.

\begin{definition}{Fuente de prueba y medio de prueba}
La \textbf{fuente de prueba} es el lugar donde queda registrado el hecho que se quiere probar: es anterior y externa al proceso (por ejemplo, la percepción directa de un vecino). El \textbf{medio de prueba} es el mecanismo que el Código habilita para incorporar esa fuente al expediente ---documental, testimonial, informativa, informática---; se produce dentro del proceso.
\end{definition}

\begin{example}{El vecino como fuente, el testigo como medio}
Los videos y las fotografías son fuentes ya registradas. Pero lo que percibieron los vecinos ---olores, cantidad de gente, cánticos--- no se incorpora librando un oficio al vecino: corresponde ofrecer prueba testimonial, porque la fuente inmediata es su propia percepción y el medio que el Código prevé para incorporarla es el testimonio.
\end{example}

\begin{important}
Tener razón sobre el fondo ---que el contrato prohíbe la actividad comercial y que hay quejas documentadas--- no garantiza ganar el juicio si la prueba se ofrece por el medio equivocado. Para litigar un proceso especial como el desalojo hay que dominar antes los institutos generales del proceso civil.
\end{important}

\section{Inconstitucionalidad y reserva del caso federal}

El control de constitucionalidad en el sistema argentino es \textbf{difuso}: no es necesario esperar a que la causa llegue a la Corte Suprema para plantear la inconstitucionalidad de una norma, que puede formularse ya ante el juez de primera instancia.

\begin{important}
El planteo de inconstitucionalidad debe formularse en la \textbf{primera presentación} ---típicamente, al contestar la demanda---; si no, resulta extemporáneo. Una vez formulado, debe mantenerse expresamente en los escritos posteriores para no perder la vía del recurso extraordinario ante la Corte Suprema: eso es la \textbf{reserva del caso federal}.
\end{important}

Las fuentes sostienen un criterio propio frente a la cláusula de estilo <<se hace reserva del caso federal>> incluida automáticamente en cualquier escrito: los derechos no se reservan sino que se ejercen para obtener su protección, y reservar un eventual daño sin una lesión constitucional concreta es una fórmula vacía, copiada de otros escritos. Solo corresponde formularla y sostenerla cuando efectivamente hay un planteo de inconstitucionalidad en danza, hecho en tiempo y forma desde la primera presentación.

\section{El circuito de presentación: nación y provincia}

Para la tasa de justicia y su pago véase el capítulo~\ref{cap:costas}. El circuito ante el Poder Judicial de la Nación (PJN) tiene tres momentos que no pueden saltearse: sorteo, carga y seguimiento.

\textbf{Sorteo.} El inicio de una causa comienza con el envío de un formulario de sorteo a la Cámara del fuero, que asigna el juzgado y el número de expediente. Este se conoce por la respuesta automática al correo desde el que se envió el formulario o por el apartado de radicaciones del sistema. El sorteo responde a que la distribución de la competencia es una cuestión de orden público, ajena a la voluntad de las partes; el Código prevé herramientas para apartar a un magistrado: la recusación sin causa (art.\ 14), limitada en el desalojo, y la recusación con causa (art.\ 17), siempre disponible ante motivos objetivos, como el parentesco con una de las partes.

\textbf{Carga.} La demanda y su documentación se cargan dentro del sistema, en la sección de escritos nuevos, identificando el expediente por su \emph{número} y no por el nombre de las partes. Toda la documentación debe subirse en PDF; no se admite Word ni imágenes JPG; se respeta el límite de 20 MB, dividiendo en partes si hace falta.

\textbf{Seguimiento.} El sistema permite localizar la causa por el apellido del actor y distingue los objetos posibles (falta de pago, vencimiento de contrato, otras causales, donde se encuadra el cambio de destino). Al regir la preclusión, es indispensable el seguimiento periódico de cada causa. Además, únicamente los días martes y viernes el sistema habilita <<dejar nota>> respecto de resoluciones no notificadas por otra vía, lo que evita considerarse notificado por ministerio de la ley.

\begin{table}[H]
\centering
\caption{Presentación electrónica: Nación y Provincia de Buenos Aires}
\begin{tabularx}{\textwidth}{>{\raggedright\arraybackslash}p{0.20\textwidth} >{\raggedright\arraybackslash}p{0.35\textwidth} X}
\toprule
\textbf{Aspecto} & \textbf{Poder Judicial de la Nación} & \textbf{Provincia de Buenos Aires (SCBA)} \\
\midrule
Acceso & Sistema del PJN, con formulario de sorteo a la Cámara. & Token para quienes tienen domicilio digital habilitado, o Mesa de Entradas Virtual de la Suprema Corte. \\
Inicio de la causa & Formulario de sorteo enviado previamente. & Carga directa de los campos, incluido el departamento judicial (según competencia territorial y eventual foro prorrogado). \\
Texto de la demanda & Se adjunta como PDF. & Se copia y pega en un campo; solo la versión firmada de forma hológrafa se adjunta como archivo. \\
Documentación & Archivos de hasta 20 MB, divididos si es necesario. & Un PDF por cada documento (contrato, informe de dominio, fotografías, por separado). \\
Firma y envío & Firma hológrafa escaneada; el trámite pasa por la aceptación del juzgado. & Firma digital registrada; envío directo. \\
\bottomrule
\end{tabularx}
\end{table}

El sistema provincial elimina varios pasos manuales del nacional: no se envía correo previo ni se consulta radicaciones, ni hay que recordar subir la demanda una vez sorteada la causa ---omisión que, en el sistema nacional, provoca la caducidad del sorteo y obliga a sortear de nuevo---. Ambos exigen el mismo rigor: conocer el circuito de cada fuero, respetar los formatos y seguir el expediente hasta que la resolución llegue a destino.

\begin{cornell}{Cuadro Cornell del capítulo~\ref{cap:practica}}
\cqa{¿Qué hacer si la prueba documental supera el límite de tamaño?}{Dividirla en partes (<<documental uno>>, <<dos>>) y adjuntar cada una por separado; el límite actual es de 20 MB.}
\cqa{¿Qué diferencia hay entre fuente y medio de prueba?}{La fuente es el lugar donde está registrado el hecho, anterior y externo al proceso; el medio es el mecanismo del Código para incorporarla.}
\cqa{¿Cuándo debe plantearse la inconstitucionalidad?}{En la primera presentación, típicamente al contestar la demanda, y mantenerse luego para reservar el caso federal.}
\cqa{¿Qué tres momentos tiene el circuito ante el PJN?}{Sorteo, carga (en PDF, por número de expediente) y seguimiento periódico, con los días de nota martes y viernes.}
\cqa{¿Cómo se inicia una causa en la Provincia de Buenos Aires?}{Por carga directa de los campos del formulario electrónico, sin sorteo previo, con firma digital registrada.}
\cresumen{La práctica del desalojo exige dominar los sistemas electrónicos de cada fuero, elegir bien fuente y medio de prueba y plantear la inconstitucionalidad en tiempo, evitando fórmulas vacías como la reserva automática del caso federal.}
\end{cornell}

\chapter{Casos integradores y talleres}\label{cap:casos}

Este capítulo reúne los casos y talleres que integran lo estudiado: un caso de falta de pago y mejoras no permitidas, y el caso de desalojo por cambio de destino, trabajado desde la consulta inicial y la intimación hasta la demanda y la audiencia preliminar. Los nombres y la ambientación de los casos fueron elegidos, según las fuentes, con intención didáctica.

\section{Falta de pago y mejoras no permitidas}

\begin{exercise}{Falta de pago y mejoras no permitidas}
\textbf{Datos:} contrato de locación de inmueble a diez años; precio con actualización mensual conforme a una tasa de referencia; el contrato prohíbe expresamente al locatario realizar mejoras. \textbf{Incumplimientos:} (1) el locatario realizó mejoras no permitidas; (2) adeuda el alquiler conforme a la forma pactada. \textbf{Consulta:} el locador pregunta cómo proceder frente a ambos incumplimientos.
\end{exercise}

\subsection{Primer análisis: la ley aplicable}

Lo primero es determinar qué ley estaba vigente cuando se firmó el contrato (art.\ 7 CCCN; véase el capítulo~\ref{cap:normativo}). Dos regímenes importan aquí: la Ley 27.551 introdujo la actualización \textbf{semestral} por un índice legal; el DNU 70/2023 volvió a permitir pactar libremente el método de actualización, incluso la mensual.

\begin{important}
La actualización mensual del caso solo sería válida si el contrato se firmó \textbf{después del DNU 70/2023}. Si se hubiera firmado durante la vigencia de la Ley 27.551, la actualización mensual libremente pactada habría sido \textbf{nula} por tratarse de una norma de orden público, y correspondería aplicar la actualización semestral por el índice legal.
\end{important}

\subsection{Falta de pago e intimación}

Ante la falta de pago, el locador debe intimar fehacientemente al locatario a abonar la suma debida en un plazo no menor a diez días, especificando el monto exacto. Solo vencido ese plazo sin regularizar el pago, y con al menos dos períodos adeudados, queda expedita la acción de desalojo por falta de pago (véase el capítulo~\ref{cap:intimacion}).

\subsection{Las mejoras no permitidas}

El CCCN prevé como causal de desalojo, además de la falta de pago, el incumplimiento de la obligación de conservar la cosa en el estado en que fue entregada. La regla general es que, si el contrato no prohíbe las mejoras, el locatario puede realizarlas; aquí el contrato las prohibía, de modo que la mejora es un \textbf{incumplimiento contractual autónomo}. Se prueba con el inventario del estado inicial, fotografías del estado inicial y posterior, y acta de constatación notarial.

\begin{important}
Para que el acta de constatación tenga valor probatorio pleno debe citarse a la parte contraria a presenciarla. Si no concurre, conviene llevar testigos: de lo contrario el acta es una prueba unilateral, realizada sin contradicción, y pierde fuerza frente al juez.
\end{important}

\begin{table}[H]
\centering
\caption{Conservación ordinaria y mejora en sentido propio}
\begin{tabularx}{\textwidth}{>{\raggedright\arraybackslash}p{0.25\textwidth} X >{\raggedright\arraybackslash}p{0.28\textwidth}}
\toprule
\textbf{Categoría} & \textbf{Ejemplo en un contrato de diez años} & \textbf{¿Alcanzada por la prohibición de mejoras?} \\
\midrule
Conservación ordinaria & Pintar la casa durante la locación; el contrato puede exigir devolverla en condiciones similares a las de la entrega, pero eso no equivale a prohibir pintar. & No. \\
Mejora en sentido propio & Tirar una pared para ampliar un ambiente, instalar una cocina nueva, obras de alteración estructural. & Sí. \\
\bottomrule
\end{tabularx}
\end{table}

\section{Desalojo por cambio de destino: el caso Riquel}

\subsection{La consulta y los hechos}

La unidad se desarrolla mediante una simulación: un grupo de estudiantes conforma un estudio jurídico y representa una consulta ante una abogada, siguiendo un guion. Esa elección pedagógica busca observar cómo llega la información desde el cliente, con sus intereses entrecruzados, antes de traducirla en categorías jurídicas. Concurren tres clientes: Marta Riquel (madre, viuda) y sus hijos Juan Román y Claudia Riquel, copropietarios en carácter de herederos de un inmueble que perteneció a los abuelos paternos de Juan Román y Claudia, quienes lo adquirieron con el fruto de su trabajo tras emigrar desde Italia. La sucesión ya estaba iniciada.

\begin{important}
\textbf{Dos conflictos, una sola consulta.} Conviven una cuestión sucesoria (Claudia y su madre quieren conservar la propiedad por su valor sentimental; Juan Román necesita hacerse de su parte indivisa para afrontar deudas personales) y una cuestión contractual y procesal (recuperar la tenencia frente al locatario). Conviene distinguir ambos planos desde el primer contacto y priorizar el que requiere acción inmediata, sin descuidar el otro.
\end{important}

El inmueble, ubicado en el barrio de Núñez, Ciudad Autónoma de Buenos Aires, estaba alquilado a Martín Milán y a quien se identifica como su media hermana. El contrato estaba vigente y el alquiler al día, lo que descarta fundar el reclamo en la falta de pago. La cláusula tercera establece: <<Destino: vivienda familiar exclusivamente>>; la abogada la verifica leyendo el contrato original que le acercan.

La pareja del locatario, identificada como Laura Bataglia, instaló en la vivienda un emprendimiento de venta de productos de un club de fútbol. Según el relato: hay compradores que forman filas porque el comercio no vende en línea; se venden camisetas, buzos, equipos deportivos, tazas y banderines; se usan la cocina, el horno y el gas para preparar productos alimenticios con los colores del club; entre cincuenta y cien personas concurren los fines de semana; hay vehículos en doble fila; y se escuchan música y cánticos a cualquier hora. La familia cuenta con fotografías, videos y mensajes de vecinos. El locatario reconoce que la familia quiere recuperar la tenencia, pero se niega a restituir el inmueble invocando que el contrato está vigente, y atribuye la actividad a su pareja.

El contrato estaba garantizado con una garantía real constituida por un matrimonio titular del inmueble ofrecido, con firmas certificadas por escribano y con informe de dominio solicitado previamente.

\subsection{Fundamento jurídico}

El derecho a exigir el desalojo por cambio de destino es una causal autónoma, distinta de la falta de pago y del vencimiento, que subsiste aunque el canon esté al día y el plazo no haya vencido (véase el capítulo~\ref{cap:vias}). El uso intensivo y comercial de la cocina y del suministro, la venta a público masivo y las quejas vecinales son la prueba concreta de que el inmueble dejó de destinarse a vivienda familiar. Como elemento adicional puede requerirse por oficio a la dependencia de habilitaciones del gobierno local un informe sobre si la zona habilita actividad comercial. Antes de demandar conviene evaluar si el cliente busca recuperar el inmueble (desalojo) o solo hacer cesar el uso indebido conservando la locación (intimación de cese). Se define como paso siguiente reunir la documentación y, en la medida de lo posible, intentar una solución sin litigio.

\subsection{El taller en tres actos}

\begin{example}{Los tres actos del taller de práctica}
Sobre un caso ficticio de desalojo por cambio de destino ---una vivienda alquilada como <<destino familiar únicamente>> que el inquilino usa para elaborar y vender comida---, el taller recorre lo que hace un estudio jurídico, trabajando dos planos en paralelo: el mundo judicial, regido por la preclusión, y el mundo de la negociación entre las partes. \textbf{Acto 1:} redactar la carta documento de intimación, del lado del locador (capítulo~\ref{cap:intimacion}). \textbf{Acto 2:} conversación entre estudios, por teléfono o videollamada, para negociar antes del juicio; no se espera acuerdo en esta etapa, pues la negociación formal se proyecta hacia la audiencia del artículo 360. \textbf{Acto 3:} redacción de la demanda de desalojo por cambio de destino (capítulo~\ref{cap:demanda}), con el artículo 680 ter como pieza clave, hasta el formulario de sorteo y el número de expediente.
\end{example}

El taller entrena habilidades transversales ---redactar con precisión, elegir domicilios, prever qué negará la contraparte y ofrecer toda la prueba desde el inicio--- y permite evitar errores caros: sortear sin la firma de todos los clientes, ignorar el reconocimiento judicial del artículo 680 ter o descuidar la prueba. También enseña que negociar con datos, plazos y respeto suele resolver más conflictos que un juicio.

\subsection{Simulacro de la audiencia preliminar}

\begin{exercise}{La audiencia preliminar del caso}
\textbf{Autos:} <<Riquel Juan Román y otros contra Milán Martín y otros sobre desalojo por cambio de destino>>, expediente 12314, Juzgado Nacional de Primera Instancia en lo Civil n.º 5. \textbf{Actuación:} audiencia preliminar (art.\ 360) iniciada a las 12:10 del jueves 1 de octubre de 2026. \textbf{Intervinientes:} la parte actora (los locadores herederos), la parte demandada (el locatario y su familia), el fiador, los letrados, la jueza y las secretarías. En el simulacro la jueza dirige la audiencia y las secretarias la ayudan si se traba.
\end{exercise}

\textbf{Apertura.} La secretaría indica hora, fecha, autos, número de expediente y juzgado. Se verifica la identidad y la comparecencia, comenzando por la parte actora: cada compareciente indica nombre y apellido y exhibe su documento; los letrados indican su carácter. Se informa que la audiencia será grabada y su registro se incorporará al expediente electrónico. En el caso, el fiador no estaba presente.

\textbf{Exposiciones.} La actora relató que heredaron el inmueble y lo alquilaron hace un año para vivienda, que se enteraron por videos de redes y por una vecina de que se lo usa para otras cosas ---gente, ruidos, olores, música y colectivos---, que el alquiler es de \$750.000 y que la casa es muy querida porque crecieron en ella, y que están destruyendo la cocina. La demandada reconoció un emprendimiento pequeño de repostería, iniciado tras perder el empleo; sostuvo que la venta es normalmente en línea, que hubo una excepción por un partido importante, que la situación es temporal y que tramitan el permiso para un local aparte. Los letrados de la demandada propusieron un acuerdo conciliatorio con costas por su orden.

\textbf{Hechos conducentes.} Conforme al inciso 3 del artículo 360 y a los escritos postulatorios, la jueza fijó como hechos conducentes y controvertidos el uso del inmueble y el cambio de destino, dado que el contrato estaba destinado únicamente a vivienda.

\textbf{Prueba.} Admitió la documental (contrato, cartas, fotografías, videos y capturas de chat) y la testimonial, y rechazó el reconocimiento judicial y la pericial ofrecidas. Ante el rechazo, algún abogado debía oponerse, y la jueza leer el artículo 379 y explicar que la resolución es inapelable (véase el capítulo~\ref{cap:audiencia}). Se fijó la testimonial para el 10 de diciembre.

\textbf{Conciliación.} La jueza invitó a conciliar, señalando que un acuerdo suele ser más rápido, menos costoso y más adecuado, y dio cinco minutos. La actora pidió el cese total del cambio de destino, el arreglo de la cocina y que el contratante viva en el inmueble; rechazó toda prórroga. La demandada ofreció el reintegro del inmueble en las condiciones en que se entregó, un plazo de tres o seis meses para tramitar el permiso, ajustar el alquiler, hacer envíos para evitar la afluencia de gente, instalar extractores para los olores y un convenio aparte para que un locador obtenga ingresos.

\textbf{Acuerdo.} Las partes informaron que se cambiará el lugar donde se fabrican las tortas, que se permitirá por tres semanas que la gente siga yendo a buscar la pastelería, pero no la producción, y que se colocará doble vidrio para aislar el sonido. La demandada manifestó su desacuerdo con la prohibición de producir en su domicilio. Respecto de las costas, cada parte pagaría las suyas. La jueza dispuso que el acuerdo conste en un acta, pase a homologación y se suba al sistema.

\subsection{Devolución sobre el simulacro}

\begin{itemize}
\item En la audiencia del artículo 360 el juez \textbf{puede proponer fórmulas conciliatorias}; como el intercambio no daba la impresión de que las partes fueran a acordar, la jueza debió intervenir proponiendo. El acuerdo conciliatorio del juzgado se distingue de la mediación privada: el juez puede dirigir como conciliador, por ejemplo preguntando a la actora si le parece prudente que el demandado se quede un mes más y luego deje la actividad, y puede advertir qué tendrá en cuenta al sentenciar.
\item \textbf{La litis debe integrarse}: el fiador es parte, y sin litis integrada es difícil acordar. En la realidad, si el fiador notificado no concurre, la audiencia se celebra igual y se le notifica el acuerdo para que preste conformidad u oposición, resguardando la defensa en juicio y la igualdad de las partes (capítulo~\ref{cap:terceros}).
\item La testimonial se fija en otro momento; la confesional no fue ofrecida, pero de haberlo sido correspondía invitar a su producción.
\item En la Capital los abogados se identifican con tomo y folio, y no con <<matrícula>>; no es relevante para el fondo. Una de las demandadas no concurrió, por lo que la litis no estaba integrada.
\item Según la devolución recibida, no había hechos controvertidos, porque ninguna parte negó lo dicho por la otra, y habría podido pasarse a trámite directo o declararse de puro derecho; sin embargo, las fotografías no estaban certificadas por escribano, por lo que no podía prescindirse de la audiencia.
% TODO: contenido incompleto o ambiguo en las notas originales: discusión de las docentes sobre trámite directo y puro derecho.
\end{itemize}

\begin{table}[H]
\centering
\caption{Acuerdo judicial y mediación privada (según lo expuesto en las fuentes)}
\begin{tabularx}{\textwidth}{>{\raggedright\arraybackslash}p{0.20\textwidth} >{\raggedright\arraybackslash}X >{\raggedright\arraybackslash}X}
\toprule
\textbf{Aspecto} & \textbf{Acuerdo conciliatorio en el juzgado} & \textbf{Mediación privada} \\
\midrule
Quién conduce & El juez puede proponer fórmulas y dirigir como conciliador. & El mediador. \\
Homologación & El juez debe homologar el acuerdo. & Según lo dicho, el acuerdo del mediador no necesita homologación. \\
Interés del juez & Terminar el proceso por una forma normal y tener un expediente menos (se mencionaron ciento cuarenta mil ingresos anuales). & --- \\
\bottomrule
\end{tabularx}
\end{table}

\begin{cornell}{Cuadro Cornell del capítulo~\ref{cap:casos}}
\cqa{¿Qué ley rige la actualización mensual del caso de mejoras?}{Solo es válida si el contrato se firmó después del DNU 70/2023; durante la Ley 27.551 sería nula por orden público y correspondería la actualización semestral.}
\cqa{¿Cómo se prueba una mejora no autorizada?}{Con inventario inicial, fotografías antes y después y acta de constatación notarial con la contraparte citada o testigos.}
\cqa{¿Qué diferencia hay entre conservación ordinaria y mejora?}{Pintar es conservación ordinaria; tirar una pared o instalar una cocina nueva es mejora en sentido propio.}
\cqa{¿Qué causal se invoca en el caso Riquel y por qué?}{Cambio de destino, porque el contrato está vigente y el alquiler al día; no corresponde falta de pago ni vencimiento.}
\cqa{¿Qué dos conflictos conviven en la consulta de la familia?}{El sucesorio entre coherederos y el contractual con el locatario; se prioriza el que exige acción inmediata.}
\cqa{¿Qué hechos conducentes fijó la jueza en el simulacro?}{El uso del inmueble y el cambio de destino, conforme al art.\ 360 inc.\ 3 y los escritos postulatorios.}
\cqa{¿Qué ocurre si el fiador no concurre a la audiencia?}{Se celebra igual y se le notifica el acuerdo para que preste conformidad u oposición.}
\cresumen{Los casos integran norma aplicable, causal, prueba, intimación, demanda y audiencia: la ley vigente al firmar el contrato define las opciones y la audiencia preliminar puede cerrar el conflicto con un acuerdo homologado.}
\end{cornell}

\backmatter

\chapter{Referencias citadas en las fuentes}

Las fuentes originales no incluyen una bibliografía formal. Se conservan únicamente las referencias que aparecen mencionadas en ellas, sin completarlas con información externa.

% TODO: verificar referencias presentes en las fuentes originales; todas aparecen incompletas (sin editorial, año ni lugar) y se conservan tal como fueron mencionadas.

\section*{Doctrina}
\begin{itemize}
\item Carli, Carlos. \emph{La demanda civil}. Obra de referencia citada para la clasificación de impedimentos procesales, excepciones y defensas (con la posición de Palacio).
\item Entelman, Remo. Concepto de conflicto como interferencia intersubjetiva de intereses.
\item Calamandrei. Doctrina sobre el artículo 230 como sistema matriz de las cautelares.
\item Lorenzetti, Ricardo. Postura sobre la existencia de consumidor y usuario en la locación.
\item Salgado. Artículo publicado en la revista \emph{Derecho Procesal} sobre diligencias preliminares y \emph{discovery}/\emph{disclosure}.
\end{itemize}

\section*{Jurisprudencia}
\begin{itemize}
\item CSJN, <<Camacho Acosta>>, Fallos 320:1633.
\item CSJN, Fallos 335:452.
\item Cámara Nacional de Apelaciones en lo Comercial en pleno, plenario del 27 de mayo de 2003, <<Recusación sin causa en juicios ejecutivos s/convocatoria a plenario>>.
\item Cámara Nacional Civil, Salas D y H, y Departamento judicial de La Plata (art.\ 684 bis).
\end{itemize}

\section*{Normativa citada}
\begin{itemize}
\item Código Civil y Comercial de la Nación (Ley 26.994): arts.\ 7, 24, 958, 1017, 1093, 1187 y ss., 1200 a 1204, 1219 a 1227, 1533 y ss.
\item Código Procesal Civil y Comercial de la Nación: arts.\ 5, 14, 17, 34 a 36, 47, 230, 260, 280, 323, 330, 356, 359, 360, 379, 498, 680, 680 bis, 680 ter, 684 bis, 688.
\item Ley 27.551, Ley 27.737, DNU 70/2023, Ley 25.488, Ley 26.589, Ley 22.172, Ley 23.091, Ley 23.187, Ley 24.522.
\item Acordada 5/2009 de la CSJN (Reglas de Brasilia).
\end{itemize}

\chapter{Apéndice: indicaciones para la evaluación final}

Las fuentes consignan, para la evaluación final oral, las siguientes indicaciones: repasar el proceso general para saber en qué punto se está parado y seguir cada una de las clases y presentaciones compartidas. La forma del examen es oral, ante dos o tres docentes, y se elige un número de una grilla que corresponde a una pregunta.

\begin{example}{Pregunta de examen: la causa de puro derecho}
<<¿En qué momento se puede declarar la causa de puro derecho en el proceso de desalojo?>> \textbf{Respuesta esperada:} hay dos momentos: el artículo 359, una vez trabada la litis, y el artículo 360, inciso 6, dentro de la audiencia. Declarar de puro derecho significa que el juez verá cuál es el derecho aplicable: ya no se discuten medios de prueba, porque no hay prueba por producir (fue producida como prueba anticipada, las partes desistieron o solo hay documental). Véase el capítulo~\ref{cap:audiencia}.
\end{example}

\end{document}
